% Lectura y comentario : discriminations et violences contre les femmes (la familia hispánica) — Espagnol Tle A — Cours signé OnBuch+
% Programme officiel MINESEC. Compiler avec : tectonic E01-discriminations-violences-femmes.tex
\newcommand{\DOCMATIERE}{Espagnol}
\newcommand{\DOCNIVEAU}{Tle A}
\newcommand{\DOCMODULE}{Módulo 1 : Vie familiale et intégration sociale}
\newcommand{\DOCLECON}{E01}
\newcommand{\DOCTITRE}{Lectura y comentario : discriminations et violences contre les femmes (la familia hispánica)}
% OnBuch+ — préambule « cours signés » ENRICHI (compilé avec tectonic / XeLaTeX)
% Système visuel « Pop » d'OnBuch+ : fond crème (jamais blanc), bordures encre
% épaisses, ombres « dures » décalées, titres Archivo Black en capitales.
% Version MATHÉMATIQUES : graphes (pgfplots/tikz), boîtes pédagogiques
% (définition, propriété, méthode, attention, le savais-tu, exercice résolu,
% exercice, TP, bilan), page de garde et signature OnBuch+ ancrée en bas.
%
% Macros attendues dans main.tex AVANT \input{preamble} :
%   \DOCMATIERE  \DOCNIVEAU  \DOCMODULE  \DOCLECON (numéro)  \DOCTITRE
\documentclass[11pt]{article}

\usepackage{fontspec}
\usepackage[spanish,french]{babel} % français = langue principale ; l'espagnol via \esp{..} / espagnol
\usepackage[a4paper,margin=2cm,top=3.4cm,bottom=2.8cm,headheight=1.4cm,footskip=1.3cm]{geometry}
\usepackage{amsmath,amssymb,amsfonts}
\usepackage{mathtools} % \xmapsto, \coloneqq... souvent utilisés par les modèles
\usepackage{cancel} % simplifications barrées (bilans de forces)
% \abs{z} (module, valeur absolue) et \norm{u} (norme), utilisés par les modèles
\providecommand{\abs}{}\let\abs\relax\DeclarePairedDelimiter{\abs}{\lvert}{\rvert}
\providecommand{\norm}{}\let\norm\relax\DeclarePairedDelimiter{\norm}{\lVert}{\rVert}
\usepackage[table]{xcolor} % [table] : fournit aussi \rowcolors (alternance de lignes)
\usepackage{pagecolor}
\usepackage{tikz}
\usetikzlibrary{arrows.meta,positioning,calc,shapes.geometric,shapes.misc,shapes.arrows,decorations.pathmorphing,decorations.pathreplacing,decorations.markings,patterns,shadows,mindmap,trees,fit,backgrounds,matrix,intersections,angles,quotes}
\usepackage{pgfplots}
\usepackage[version=4]{mhchem} % équations nucléaires \ce{^{235}_{92}U}
\usepackage[european,siunitx]{circuitikz} % schémas électriques
\pgfplotsset{compat=1.18}
\usepgfplotslibrary{fillbetween,groupplots} % aires entre courbes (calcul intégral)
\usepackage{enumitem}
\usepackage{fancyhdr}
% [most] charge toutes les bibliothèques tcolorbox (listings, minted, raster,
% documentation...) et gonfle inutilement la mémoire TeX sur un document long
% (~300 boîtes) : on ne charge que ce qui est réellement utilisé.
\usepackage[skins,breakable]{tcolorbox}
\usepackage{graphicx}
\usepackage{colortbl}
\usepackage{booktabs}
\usepackage{tabularx}
\usepackage{diagbox} % case d'en-tête diagonale des tableaux à double entrée
% Colonnes extensibles centrée / gauche / droite (C, L, R), souvent utilisées par les modèles
\newcolumntype{C}{>{\centering\arraybackslash}X}
\newcolumntype{L}{>{\raggedright\arraybackslash}X}
\newcolumntype{R}{>{\raggedleft\arraybackslash}X}
\usepackage{array}
\usepackage{multicol}
\usepackage{siunitx}
% parse-numbers=false : accepte aussi \SI{2,5 \times 10^{-3}}{...}
\sisetup{locale=FR,output-decimal-marker={,},group-minimum-digits=4,parse-numbers=false}
\DeclareSIUnit\ohm{\ensuremath{\symup{\Omega}}} % U+2126 absent de la police
\DeclareSIUnit\division{div} % réglages d'oscilloscope (ms/div, V/div)
\DeclareSIUnit\tr{tr} % tours (vitesse de rotation, ex. \SI{1500}{\tr\per\minute})
\DeclareSIUnit\molar{M} % concentration molaire (ex. \SI{3}{\molar}, \SI{1}{\micro\molar})
\DeclareSIUnit\an{an} % année (le modèle confond parfois avec \year, un compteur TeX) — ex. \SI{5}{\kilo\meter\per\an}
\DeclareSIUnit\FCFA{FCFA} % franc CFA (ex. \SI{500}{\FCFA\per\kilo\gram})
\DeclareSIUnit\degreeBrix{\textdegree Brix} % teneur en sucre d'un sirop/jus (réfractométrie)
\DeclareSIUnit\month{mois} % le modèle utilise parfois \month, un compteur TeX, comme unité de durée
\usepackage{qrcode}
% \degree : commande fréquemment utilisée par les modèles pour un angle ou une
% température en dehors de \SI{}{} (non fournie par les paquets chargés).
\providecommand{\degree}{^{\circ}}
% \qed (fin de démonstration), utilisé par les modèles sans amsthm
\providecommand{\qed}{\hfill$\square$}
% \barre : double barre des valeurs interdites dans un tableau de variations (tkz-tab)
\providecommand{\barre}{\ensuremath{\|}}
% environnement proof (amsthm non chargé) utilisé par les modèles
\ifdefined\proof\else\newenvironment{proof}[1][Démonstration]{\par\noindent\textit{#1.}\ }{\qed\par}\fi
\usepackage{lastpage}
\usepackage{hyperref}
\hypersetup{hidelinks}

% ============================================================
% Palette « Pop » OnBuch+ — mêmes hex que la classe P (plus_ui.dart)
% ============================================================
\definecolor{poporange}{HTML}{F59321}
\definecolor{popdark}{HTML}{E07A0C}
\definecolor{popcream}{HTML}{FFF6E8}
\definecolor{popink}{HTML}{1C1714}
\definecolor{poppurple}{HTML}{7A5AE0}
\definecolor{popgreen}{HTML}{2E9E6A}
\definecolor{poppink}{HTML}{E0457B}
\definecolor{popblue}{HTML}{2F7DE1}
\definecolor{popgold}{HTML}{C98A12}
\definecolor{popmuted}{HTML}{8A7F76}
\definecolor{popline}{HTML}{EADFD2}
\definecolor{popgray}{HTML}{8A7F76}
\colorlet{popgrey}{popgray}
% Alias tolérants (le modèle nomme parfois les couleurs différemment)
\colorlet{popwhite}{white}
\colorlet{popblack}{popink}
\colorlet{popred}{poppink}
\colorlet{popyellow}{popgold}
% Teintes claires (fonds de tableaux/encadrés) — le modèle nomme parfois
% les couleurs avec un suffixe L (ex. popblueL) sans les avoir définies.
\colorlet{poporangeL}{poporange!20}
\colorlet{popdarkL}{popdark!20}
\colorlet{poppurpleL}{poppurple!20}
\colorlet{popgreenL}{popgreen!20}
\colorlet{poppinkL}{poppink!20}
\colorlet{popblueL}{popblue!20}
\colorlet{popgoldL}{popgold!20}
\colorlet{popredL}{popred!20}
\colorlet{popgrayL}{popgray!20}
% Teintes claires pour fonds de tableaux / zones
\colorlet{popblueL}{popblue!10!white}
\colorlet{poporangeL}{poporange!14!white}
\colorlet{popgreenL}{popgreen!12!white}
\colorlet{poppurpleL}{poppurple!12!white}
\colorlet{poppinkL}{poppink!12!white}
\colorlet{popgoldL}{popgold!14!white}

\pagecolor{popcream}

% ============================================================
% Typographie
% ============================================================
\defaultfontfeatures{Ligatures=TeX}
\setmainfont{PlusJakartaSans-Medium.ttf}[Path=../../../fonts/,BoldFont=PlusJakartaSans-Bold.ttf]
% Maths en sans-serif (Fira Math) avec chiffres et lettres droites de Plus
% Jakarta, pour que les formules se fondent dans le texte.
\usepackage[math-style=ISO,bold-style=ISO]{unicode-math}
\setmathfont{FiraMath-Regular.otf}[Path=../../../fonts/]
\setmathfont{PlusJakartaSans-Medium.ttf}[Path=../../../fonts/,range={up/num,up/{Latin,latin},\mathup}]
\usepackage{icomma} % virgule décimale sans espace en mode math
% Caractères mathématiques Unicode absents des polices de texte (Plus Jakarta,
% Archivo Black) : rendus par leur équivalent mathématique, en texte comme en
% formule — sinon ils s'affichent en carré vide (ex. « Arithmétique dans ℤ »).
\usepackage{newunicodechar}
\newunicodechar{ℝ}{\ensuremath{\mathbb{R}}}
\newunicodechar{ℤ}{\ensuremath{\mathbb{Z}}}
\newunicodechar{ℂ}{\ensuremath{\mathbb{C}}}
\newunicodechar{ℕ}{\ensuremath{\mathbb{N}}}
\newunicodechar{ℚ}{\ensuremath{\mathbb{Q}}}
\newunicodechar{𝔻}{\ensuremath{\mathbb{D}}}
\newunicodechar{α}{\ensuremath{\alpha}}
\newunicodechar{β}{\ensuremath{\beta}}
\newunicodechar{γ}{\ensuremath{\gamma}}
\newunicodechar{δ}{\ensuremath{\delta}}
\newunicodechar{ε}{\ensuremath{\varepsilon}}
\newunicodechar{θ}{\ensuremath{\theta}}
\newunicodechar{λ}{\ensuremath{\lambda}}
\newunicodechar{μ}{\ensuremath{\mu}}
\newunicodechar{σ}{\ensuremath{\sigma}}
\newunicodechar{τ}{\ensuremath{\tau}}
\newunicodechar{φ}{\ensuremath{\varphi}}
\newunicodechar{ψ}{\ensuremath{\psi}}
\newunicodechar{ω}{\ensuremath{\omega}}
\newunicodechar{Σ}{\ensuremath{\Sigma}}
% majuscules produites par \MakeUppercase dans les titres (α-aminés -> Α)
\newunicodechar{Α}{\ensuremath{\alpha}}
\newunicodechar{Β}{\ensuremath{\beta}}
\newunicodechar{Γ}{\ensuremath{\Gamma}}
\newunicodechar{Δ}{\ensuremath{\Delta}}
\newunicodechar{Ε}{\ensuremath{\varepsilon}}
\newunicodechar{Θ}{\ensuremath{\Theta}}
\newunicodechar{Λ}{\ensuremath{\Lambda}}
\newunicodechar{Μ}{\ensuremath{\mu}}
\newunicodechar{Τ}{\ensuremath{\tau}}
\newunicodechar{Φ}{\ensuremath{\Phi}}
\newunicodechar{Ψ}{\ensuremath{\Psi}}
\newunicodechar{Ω}{\ensuremath{\Omega}}
\newunicodechar{↦}{\ensuremath{\mapsto}}
\newunicodechar{∈}{\ensuremath{\in}}
\newunicodechar{∉}{\ensuremath{\notin}}
\newunicodechar{∧}{\ensuremath{\wedge}}
\newunicodechar{⇔}{\ensuremath{\Leftrightarrow}}
\newunicodechar{⟺}{\ensuremath{\Longleftrightarrow}}
\newunicodechar{⇒}{\ensuremath{\Rightarrow}}
\newunicodechar{⟹}{\ensuremath{\Longrightarrow}}
\newunicodechar{∘}{\ensuremath{\circ}}
\newunicodechar{∪}{\ensuremath{\cup}}
\newunicodechar{∩}{\ensuremath{\cap}}
\newunicodechar{≡}{\ensuremath{\equiv}}
\newunicodechar{⊂}{\ensuremath{\subset}}
\newunicodechar{⊥}{\ensuremath{\perp}}
\newunicodechar{∃}{\ensuremath{\exists}}
\newunicodechar{∀}{\ensuremath{\forall}}
\newunicodechar{∛}{\ensuremath{\sqrt[3]{\,}}}
\newunicodechar{∜}{\ensuremath{\sqrt[4]{\,}}}
\newunicodechar{⃗}{\ensuremath{{}^{\to}}}
\newunicodechar{ⁿ}{\ifmmode^{n}\else\textsuperscript{\ensuremath{n}}\fi}
\newunicodechar{ˣ}{\ifmmode^{x}\else\textsuperscript{\ensuremath{x}}\fi}
\newunicodechar{ᵇ}{\ifmmode^{b}\else\textsuperscript{\ensuremath{b}}\fi}
\newunicodechar{ʳ}{\ifmmode^{r}\else\textsuperscript{\ensuremath{r}}\fi}
\newunicodechar{ᵘ}{\ifmmode^{u}\else\textsuperscript{\ensuremath{u}}\fi}
\newunicodechar{ʸ}{\ifmmode^{y}\else\textsuperscript{\ensuremath{y}}\fi}
\newunicodechar{ᵃ}{\ifmmode^{a}\else\textsuperscript{\ensuremath{a}}\fi}
\newunicodechar{ᵉ}{\ifmmode^{e}\else\textsuperscript{\ensuremath{e}}\fi}
\newunicodechar{ᵏ}{\ifmmode^{k}\else\textsuperscript{\ensuremath{k}}\fi}
\newunicodechar{ᵖ}{\ifmmode^{p}\else\textsuperscript{\ensuremath{p}}\fi}
\newunicodechar{ᴺ}{\ifmmode^{N}\else\textsuperscript{\ensuremath{N}}\fi}
\newunicodechar{ᶜ}{\ifmmode^{c}\else\textsuperscript{\ensuremath{c}}\fi}
\newunicodechar{ʰ}{\ifmmode^{h}\else\textsuperscript{\ensuremath{h}}\fi}
\newunicodechar{ᵠ}{\ifmmode^{q}\else\textsuperscript{\ensuremath{q}}\fi}
\newunicodechar{ₙ}{\ifmmode_{n}\else\textsubscript{\ensuremath{n}}\fi}
\newunicodechar{ₐ}{\ifmmode_{a}\else\textsubscript{\ensuremath{a}}\fi}
\newunicodechar{ᵢ}{\ifmmode_{i}\else\textsubscript{\ensuremath{i}}\fi}
\newunicodechar{ᵣ}{\ifmmode_{r}\else\textsubscript{\ensuremath{r}}\fi}
\newunicodechar{ₖ}{\ifmmode_{k}\else\textsubscript{\ensuremath{k}}\fi}
\newunicodechar{ᵦ}{\ifmmode_{\beta}\else\textsubscript{\ensuremath{\beta}}\fi}
\newunicodechar{ₓ}{\ifmmode_{x}\else\textsubscript{\ensuremath{x}}\fi}
\newfontfamily\popdisplay{ArchivoBlack-Regular.ttf}[Path=../../../fonts/]
\newfontfamily\poplogo{SpaceGrotesk-Bold.ttf}[Path=../../../fonts/]
\newfontfamily\popsemi{PlusJakartaSans-SemiBold.ttf}[Path=../../../fonts/]
\newfontfamily\popxbold{PlusJakartaSans-ExtraBold.ttf}[Path=../../../fonts/]
\newfontfamily\popxboldls{PlusJakartaSans-ExtraBold.ttf}[Path=../../../fonts/,LetterSpace=3.0]

\setlist[enumerate]{leftmargin=1.7em,itemsep=5pt,topsep=4pt}
\setlist[itemize]{leftmargin=1.7em,itemsep=4pt,topsep=4pt,label={\color{poporange}\textbullet}}
\setlength{\parindent}{0pt}
\setlength{\parskip}{5pt}
\renewcommand{\arraystretch}{1.25}

% pgfplots : style Pop par défaut
\pgfplotsset{
  every axis/.append style={axis line style={popink,line width=1pt},tick style={popink},
    grid style={popline,line width=0.5pt},label style={font=\small},tick label style={font=\footnotesize}},
  popcurve/.style={poporange,line width=1.8pt,no markers,smooth},
}
% Même style disponible dans un \draw TikZ simple (hors pgfplots)
\tikzset{popcurve/.style={poporange,line width=1.8pt}}
\tikzset{popgrid/.style={popline,very thin}}
\colorlet{popcurve}{poporange} % le nom du style est parfois utilisé comme couleur

% ============================================================
% Pill / squiggle
% ============================================================
\newcommand{\poppill}[3]{%
  \tikz[baseline=-0.55ex]\node[draw=popink,line width=1.2pt,rounded corners=2.6pt,%
    fill=#2,inner xsep=6.5pt,inner ysep=3pt]{%
    \popxboldls\fontsize{7}{7}\selectfont\color{#3}\MakeUppercase{#1}};%
}
\newcommand{\popsquiggle}[1]{%
  \tikz[baseline=-0.4ex]{
    \draw[#1,line width=2.6pt,line cap=round]
      plot [smooth] coordinates {(0,0.09) (0.34,0.01) (0.68,0.21) (1.02,0.01) (1.36,0.10)};
  }%
}
% Mot-clé surligné (marqueur orange sous le texte)
\newcommand{\cle}[1]{{\popxbold\color{popdark}#1}}

% ============================================================
% En-tête / pied
% ============================================================
\pagestyle{fancy}
\fancyhf{}
\renewcommand{\headrulewidth}{0pt}
\renewcommand{\footrulewidth}{0pt}
\lhead{\raisebox{-0.15ex}{\poplogo\fontsize{13}{13}\selectfont\color{popink}OnBuch\color{poporange}+}}
\rhead{\poppill{\DOCMATIERE\ \textbullet\ \DOCNIVEAU\ \textbullet\ Leçon \DOCLECON}{white}{popink}}
\lfoot{\popxbold\fontsize{7.6}{7.6}\selectfont\color{popmuted}\MakeUppercase{Cours signé OnBuch+}\ \textbullet\ {\popsemi\DOCTITRE}}
\rfoot{\tikz[baseline=-0.6ex]\node[rounded corners=8.5pt,fill=popink,minimum height=17pt,minimum width=17pt,inner xsep=5pt,inner sep=0]{\popxbold\fontsize{8}{8}\selectfont\color{popcream}\thepage\,/\,\pageref*{LastPage}};}

% ============================================================
% Titres
% ============================================================
\newcommand{\chaptitre}[2]{%
  \begin{tcolorbox}[enhanced,colback=poporange,colframe=popink,boxrule=2.6pt,arc=11pt,
      drop shadow=popink,left=15pt,right=15pt,top=12pt,bottom=11pt,before skip=3pt,after skip=18pt]
    \poppill{#2}{popink}{popcream}\par\vspace{9pt}
    {\raggedright\hyphenpenalty=10000\exhyphenpenalty=10000\popdisplay\fontsize{22}{24}\selectfont\color{popcream}\MakeUppercase{#1}\par}
  \end{tcolorbox}
}
% Section numérotée : pastille orange avec le numéro + titre Archivo
\newcounter{popsec}
\newcommand{\coursec}[1]{%
  \stepcounter{popsec}\setcounter{popsub}{0}%
  \par\vspace{16pt}\noindent
  \begin{minipage}{\linewidth}
  \tikz[baseline=-0.7ex]\node[draw=popink,line width=1.6pt,fill=poporange,rounded corners=4pt,minimum size=22pt,inner sep=0]{\popdisplay\fontsize{12}{12}\selectfont\color{popcream}\thepopsec};%
  \hspace{8pt}{\popdisplay\fontsize{15}{17}\selectfont\color{popink}\MakeUppercase{#1}}\par
  \vspace{4pt}\noindent{\color{poporange}\rule{36pt}{3.6pt}}
  \end{minipage}\par\nopagebreak\vspace{8pt}
}
\newcounter{popsub}
\newcommand{\courssub}[1]{%
  \stepcounter{popsub}%
  \par\vspace{9pt}\noindent{\popxbold\fontsize{11.5}{13}\selectfont\color{popink}{\color{poporange}\thepopsec.\thepopsub}\hspace{6pt}#1}\par\nopagebreak\vspace{4pt}
}

% ============================================================
% Boîtes pédagogiques — même recette : fond blanc, bordure épaisse colorée,
% ombre dure encre, badge plein. Titre optionnel [#1].
% ============================================================
\tcbset{popbase/.style={breakable,enhanced,colback=white,boxrule=2.3pt,arc=9pt,
  shadow={3pt}{-3pt}{0pt}{popink},left=14pt,right=14pt,top=11pt,bottom=11pt,before skip=11pt,after skip=15pt,
  fontupper=\popsemi\fontsize{10.6}{14.5}\selectfont\color{popink},
  fontlower=\popsemi\fontsize{10.6}{14.5}\selectfont\color{popink}}}
% Coche dessinée (le glyphe ✓ n'existe pas dans les polices du document)
\newcommand{\popcheck}[1]{\tikz[baseline=-0.5ex]\draw[#1,line width=1.6pt,line cap=round,line join=round](-0.13,0.0)--(-0.04,-0.09)--(0.13,0.1);}
\providecommand{\coche}{\popcheck{popgreen}}
\providecommand{\resultat}[1]{\par\smallskip\noindent\fcolorbox{popgreen}{popgreenL}{\parbox{\dimexpr\linewidth-2\fboxsep-2\fboxrule\relax}{\textbf{Résultat.} #1}}\par\smallskip}
\newcommand{\popboxhead}[3]{\poppill{#1}{#2}{white}\ifx\relax#3\relax\else\hspace{6pt}{\popxbold\fontsize{10.6}{12}\selectfont\color{popink}#3}\fi\par\vspace{7pt}}

\newtcolorbox{aretenir}[1][]{popbase,colframe=popgreen,before upper={\popboxhead{À retenir}{popgreen}{#1}}}
% Texte en espagnol (document de lecture, dialogue, poème, extrait) : cadre
% d'étude. Un titre optionnel (source, auteur) s'affiche dans l'en-tête.
\newtcolorbox{textebox}[1][]{popbase,colframe=popblue,before upper={\popboxhead{Texto}{popblue}{#1}\begin{otherlanguage}{spanish}},after upper={\end{otherlanguage}}}
% Marques ✓ / ✗ (forme correcte / fautive) souvent utilisées par les modèles
\providecommand{\cross}{\textcolor{poppink}{\ensuremath{\times}}}
\providecommand{\xmark}{\cross}\providecommand{\crossmark}{\cross}\providecommand{\cmark}{\ensuremath{\checkmark}}
% Ligne de réponse / justification à compléter (inventée par les modèles)
\providecommand{\justification}[1]{\par\noindent\textit{Justification~:} \dotfill\par}
% \textipa (package tipa non chargé) : la police ne couvre pas l'API -> simple passage
\providecommand{\textipa}[1]{#1}
% Mot ou phrase en espagnol à l'intérieur d'un paragraphe français
\newcommand{\esp}[1]{\foreignlanguage{spanish}{\emph{#1}}}
\newtcolorbox{exemplebox}[1][]{popbase,colframe=poporange,before upper={\popboxhead{Exemple}{poporange}{#1}}}
\newtcolorbox{definition}[1][]{popbase,colframe=popblue,before upper={\popboxhead{Définition}{popblue}{#1}}}
\newtcolorbox{propriete}[1][]{popbase,colframe=poppurple,before upper={\popboxhead{Propriété}{poppurple}{#1}}}
\newtcolorbox{methode}[1][]{popbase,colframe=popgold,before upper={\popboxhead{Méthode}{popgold}{#1}}}
\newtcolorbox{attention}[1][]{popbase,colframe=poppink,before upper={\popboxhead{Attention}{poppink}{#1}}}
\newtcolorbox{experience}[1][]{popbase,colframe=popblue,colback=popblueL,before upper={\popboxhead{Expérience}{popblue}{#1}}}
% « Le savais-tu ? » : carte violette pleine (contexte, culture, Cameroun)
\newtcolorbox{savaistu}[1][]{popbase,colback=poppurple,colframe=popink,
  fontupper=\popsemi\fontsize{10.4}{14.2}\selectfont\color{white},
  before upper={\poppill{Le savais-tu\,?}{white}{poppurple}\ifx\relax#1\relax\else\hspace{6pt}{\popxbold\color{white}#1}\fi\par\vspace{7pt}}}
% Exercice résolu : énoncé en haut, \tcblower puis la solution sur fond crème
\newcounter{popexo}\newcounter{popexr}
\newtcolorbox{exoresolu}[1][]{popbase,colframe=poporange,colbacklower=popcream,
  segmentation style={popink,line width=1.2pt,dashed},
  before upper={\stepcounter{popexr}\popboxhead{Exercice résolu \thepopexr}{poporange}{#1}},
  before lower={\poppill{Solution}{popink}{popcream}\par\vspace{6pt}}}
% Exercice (non résolu en place) — la correction est dans la partie Corrigés
\newtcolorbox{exercice}[1][]{popbase,colframe=popink,
  before upper={\stepcounter{popexo}\popboxhead{Exercice \thepopexo}{popink}{#1}}}
% Corrigé
\newtcolorbox{corrige}[1][]{popbase,colframe=popgreen,colback=popgreenL,
  before upper={\popboxhead{Corrigé}{popgreen}{#1}}}
% Objectifs / prérequis / bilan
\newtcolorbox{objectifs}{popbase,colframe=popink,colback=poporangeL,
  before upper={\popboxhead{Objectifs de la leçon}{popink}{}}}
\newtcolorbox{prerequis}{popbase,colframe=popink,colback=popblueL,
  before upper={\popboxhead{Prérequis}{popblue}{}}}
\newtcolorbox{bilan}{popbase,colframe=popink,colback=white,boxrule=2.8pt,
  before upper={\popboxhead{Fiche bilan}{poporange}{L'essentiel à connaître pour l'examen}}}
% Figure encadrée avec légende
\newtcolorbox{popfigure}[1][]{enhanced,breakable=false,colback=white,colframe=popline,boxrule=1.4pt,arc=8pt,
  left=10pt,right=10pt,top=10pt,bottom=8pt,before skip=10pt,after skip=12pt,halign=center,
  lower separated=false,halign lower=center,
  fontlower=\popsemi\fontsize{9}{11.5}\selectfont\color{popmuted},
  #1}
\newcommand{\legende}[1]{\par\vspace{6pt}{\popsemi\fontsize{9}{11.5}\selectfont\color{popmuted}#1}}

% ============================================================
% Signature OnBuch+
% ============================================================
\newcommand{\poplogoinline}[1]{{\poplogo\fontsize{#1}{#1}\selectfont\color{popink}OnBuch\color{poporange}+}}
\newcommand{\signatureonbuch}{%
  \begin{tcolorbox}[enhanced,colback=popink,colframe=popink,boxrule=2.2pt,arc=10pt,
      drop shadow=poporange,left=14pt,right=14pt,top=10pt,bottom=10pt,before skip=0pt,after skip=0pt]
    \tikz[baseline=-0.7ex]\node[circle,fill=popgreen,draw=popcream,line width=1.3pt,minimum size=22pt,inner sep=0]{\popcheck{white}};%
    \hspace{9pt}\begin{minipage}[c]{0.62\linewidth}
      {\popxbold\fontsize{12}{13}\selectfont\color{popcream}Signé\ \textbullet\ {\poplogo OnBuch\color{poporange}+}}\par\vspace{2pt}
      {\raggedright\popsemi\fontsize{8}{10.5}\selectfont\color{popline}\MakeUppercase{Équipe pédagogique OnBuch+}\\\MakeUppercase{Conforme au programme officiel MINESEC}\par}
    \end{minipage}\hfill
    {\poplogo\fontsize{22}{22}\selectfont\color{popcream}OnBuch\color{poporange}+}
  \end{tcolorbox}%
}

% ============================================================
% Page de garde
% #1 titre · #2 matière · #3 niveau · #4 module · #5 n° leçon · #6 durée ·
% #7 description / objectifs (texte ou itemize)
% ============================================================
\newcommand{\pagegarde}[7]{%
  \thispagestyle{empty}
  \newgeometry{a4paper,margin=1.7cm,top=1.6cm,bottom=1.4cm}%
  \begin{tikzpicture}[remember picture,overlay]
    \fill[poporange!18] ([xshift=-3.2cm,yshift=-3.6cm]current page.north east) circle (4.6cm);
    \fill[poppurple!14] ([xshift=2.4cm,yshift=7.6cm]current page.south west) circle (3.8cm);
  \end{tikzpicture}%
  {\poplogo\fontsize{30}{30}\selectfont\color{popink}OnBuch\color{poporange}+}\hfill
  \raisebox{0.6ex}{\poppill{Cours signé \textbullet\ Premium}{popink}{popcream}}\par
  \vspace{1pt}\popsquiggle{poppurple}\par
  \vspace{8pt}
  \begin{tcolorbox}[enhanced,colback=poporange,colframe=popink,boxrule=2.8pt,arc=13pt,
      drop shadow=popink,left=18pt,right=18pt,top=12pt,bottom=13pt,after skip=10pt]
    \poppill{#2 \textbullet\ #3}{popink}{popcream}\hspace{5pt}\poppill{#4}{white}{popink}\par\vspace{8pt}
    {\popdisplay\fontsize{14}{14}\selectfont\color{popink}LEÇON #5}\par\vspace{4pt}
    {\raggedright\hyphenpenalty=10000\exhyphenpenalty=10000\popdisplay\fontsize{25}{27}\selectfont\color{popcream}\MakeUppercase{#1}\par}
    \vspace{8pt}
    {\popxbold\fontsize{9.5}{9.5}\selectfont\color{popink}\MakeUppercase{Durée conseillée : #6}}
  \end{tcolorbox}
  \begin{tcolorbox}[enhanced,colback=white,colframe=popink,boxrule=2.2pt,arc=10pt,
      drop shadow=popink,left=16pt,right=16pt,top=10pt,bottom=10pt,after skip=8pt]
    \poppill{Dans cette leçon}{poppurple}{white}\par\vspace{5pt}
    \popsemi\fontsize{9.3}{12.6}\selectfont\color{popink}#7
  \end{tcolorbox}
  \noindent\begin{minipage}[t]{0.487\linewidth}
    \begin{tcolorbox}[enhanced,colback=popgreen,colframe=popink,boxrule=2.2pt,arc=10pt,
        drop shadow=popink,left=11pt,right=11pt,top=10pt,bottom=10pt,height=3.1cm,valign=center]
      \tcbox[colback=white,colframe=popink,boxrule=1.3pt,arc=5pt,boxsep=0pt,nobeforeafter,
        left=3pt,right=3pt,top=3pt,bottom=3pt,valign=center]{\qrcode[height=1.9cm]{https://play.google.com/store/apps/details?id=com.luvvix.onbuch&hl=fr}}%
      \hspace{8pt}\begin{minipage}[c]{0.5\linewidth}
        {\popdisplay\fontsize{11}{12.5}\selectfont\color{white}\MakeUppercase{Télécharge OnBuch}}\par\vspace{4pt}
        {\raggedright\popsemi\fontsize{8.4}{11}\selectfont\color{white}Gratuit \textbullet\ Google Play\\Tuteur IA, annales, concours, résultats\par}
      \end{minipage}
    \end{tcolorbox}
  \end{minipage}\hfill
  \begin{minipage}[t]{0.487\linewidth}
    \begin{tcolorbox}[enhanced,colback=poppurple,colframe=popink,boxrule=2.2pt,arc=10pt,
        drop shadow=popink,left=11pt,right=11pt,top=10pt,bottom=10pt,height=3.1cm,valign=center]
      \tikz[baseline=-0.7ex]\node[circle,fill=white,draw=popink,line width=1.3pt,minimum size=1.6cm,inner sep=0]{\popdisplay\fontsize{24}{24}\selectfont\color{poppurple}+};%
      \hspace{8pt}\begin{minipage}[c]{0.6\linewidth}
        {\popdisplay\fontsize{11}{12.5}\selectfont\color{white}\MakeUppercase{Passe à OnBuch+}}\par\vspace{4pt}
        {\raggedright\popsemi\fontsize{8.4}{11}\selectfont\color{white}Tous les cours signés, Tuteur IA illimité, fiches PDF — onglet OnBuch+ de l'app\par}
      \end{minipage}
    \end{tcolorbox}
  \end{minipage}
  \vspace{8pt}
  \popcontenu
  \vfill
  \signatureonbuch
  \restoregeometry
  \newpage
}

% Rangée « Ce cours contient » (page de garde)
\newcommand{\poptile}[3]{% #1 couleur #2 gros texte #3 légende
  \begin{tcolorbox}[enhanced,colback=#1,colframe=popink,boxrule=2pt,arc=9pt,shadow={2.5pt}{-2.5pt}{0pt}{popink},
      left=6pt,right=6pt,top=9pt,bottom=9pt,halign=center,valign=center,height=1.75cm,width=\linewidth,nobeforeafter]
    {\popdisplay\fontsize{12.5}{13}\selectfont\color{white}\MakeUppercase{#2}}\par\vspace{3pt}
    {\centering\popsemi\fontsize{7.8}{9.5}\selectfont\color{white}#3\par}
  \end{tcolorbox}}
\newcommand{\popcontenu}{%
  \par\noindent{\popxboldls\fontsize{7.5}{7.5}\selectfont\color{popmuted}CE COURS SIGNÉ CONTIENT}\par\vspace{7pt}
  \noindent\begin{minipage}[t]{0.235\linewidth}\poptile{popblue}{Cours}{complet, illustré, pas à pas}\end{minipage}\hfill
  \begin{minipage}[t]{0.235\linewidth}\poptile{poporange}{Méthodes}{et exercices résolus}\end{minipage}\hfill
  \begin{minipage}[t]{0.235\linewidth}\poptile{poppink}{Exercices}{type examen + corrigés}\end{minipage}\hfill
  \begin{minipage}[t]{0.235\linewidth}\poptile{popgreen}{Fiche bilan}{l'essentiel à retenir}\end{minipage}\par}

% Sommaire
\newtcolorbox{sommaire}{enhanced,colback=white,colframe=popink,boxrule=2.2pt,arc=10pt,
  drop shadow=popink,left=18pt,right=18pt,top=13pt,bottom=13pt,before skip=6pt,after skip=20pt,
  before upper={\poppill{Sommaire}{poppurple}{white}\par\vspace{9pt}\popsemi\fontsize{10.6}{14.5}\selectfont\color{popink}}}

% Signature de fin de cours — poussée tout en bas de la dernière page
\newcommand{\findecours}{%
  \par\vfill\nopagebreak
  \begin{center}\popsquiggle{poporange}\end{center}\vspace{4pt}
  \signatureonbuch
}
\AtBeginDocument{\def\llbracket{\mathopen{[\mkern-3.5mu[}}\def\rrbracket{\mathclose{]\mkern-3.5mu]}}} % intervalles d'entiers (glyphe ⟦ absent de la police)
% Glyphes absents de Fira Math : redéfinis à partir de glyphes disponibles
\AtBeginDocument{%
  \def\setminus{\mathbin{\backslash}}\let\smallsetminus\setminus
  \def\vdots{\mathord{\vbox{\baselineskip3.2pt\lineskiplimit0pt\kern4pt\hbox{.}\hbox{.}\hbox{.}}}}%
  \def\ddots{\mathinner{\mkern1mu\raise6.5pt\hbox{.}\mkern2mu\raise3.6pt\hbox{.}\mkern2mu\raise0.7pt\hbox{.}\mkern1mu}}%
  \def\triangle{\mathord{\scalebox{0.9}{$\symup{\Delta}$}}}%
  \def\nabla{\mathord{\raisebox{\depth}{\scalebox{1}[-1]{$\symup{\Delta}$}}}}%
  \def\checkmark{\popcheck{popdark}}%
  \def\star{\mathbin{\ast}}\def\bigstar{\mathord{\ast}}%
  \def\diamond{\mathbin{\scalebox{0.6}{$\Diamond$}}}%
  \def\bigcup{\mathop{\mathchoice{\raisebox{-0.3em}{\scalebox{1.7}{$\cup$}}}{\scalebox{1.25}{$\cup$}}{\cup}{\cup}}\displaylimits}%
  \def\bigcap{\mathop{\mathchoice{\raisebox{-0.3em}{\scalebox{1.7}{$\cap$}}}{\scalebox{1.25}{$\cap$}}{\cap}{\cap}}\displaylimits}%
  \def\lesseqqgtr{\mathrel{\lessgtr}}\def\bigoplus{\mathop{\oplus}\displaylimits}%
}
\providecommand{\ssi}{si et seulement si} % abréviation fréquente
\AtBeginDocument{\providecommand{\wideparen}{\overparen}} % arc de cercle

\begin{document}

\pagegarde{Lectura y comentario : discriminations et violences contre les femmes (la familia hispánica)}{Espagnol}{Tle A}{Módulo 1 : Vie familiale et intégration sociale}{E01}{2 h}{Cette leçon de lecture et commentaire propose un texte original en espagnol sur les discriminations et les violences faites aux femmes dans le monde hispanique. Elle développe le lexique thématique de la famille et de la condition féminine et entraîne les élèves à la méthode du commentaire de texte exigée au Baccalauréat.
\vspace{4pt}
\begin{multicols}{2}\small
\begin{itemize}
\item À la fin de cette leçon, je sais lire et comprendre globalement et finement un texte espagnol sur les violences faites aux femmes.
\item À la fin de cette leçon, je sais identifier le thème, la thèse et les idées directrices d'un texte argumentatif.
\item À la fin de cette leçon, je sais relever, définir et réemployer le lexique de la famille hispanique et de la condition féminine (violencia de género, machismo, maltrato, denuncia, igualdad, acoso).
\item À la fin de cette leçon, je sais répondre par des phrases complètes à des questions de compréhension en espagnol.
\item À la fin de cette leçon, je sais présenter un texte (nature, source, auteur, thème) selon la méthode du commentaire.
\item À la fin de cette leçon, je sais dégager deux ou trois idées directrices et les justifier par des citations du texte.
\end{itemize}\end{multicols}}

\begin{sommaire}
\begin{multicols}{2}
\begin{enumerate}
\item Texte support : lecture et découverte
\item Compréhension du texte
\item Lexique thématique : la famille et la condition de la femme
\item Méthode du commentaire de texte
\item Commentaire modèle et production guidée
\item Activité d'intégration
\item Méthodes et exercices résolus
\item Exercices
\item Corrigés des exercices
\item Fiche bilan
\end{enumerate}
\end{multicols}
\end{sommaire}

\begin{objectifs}
\begin{itemize}
\item À la fin de cette leçon, je sais lire et comprendre globalement et finement un texte espagnol sur les violences faites aux femmes.
\item À la fin de cette leçon, je sais identifier le thème, la thèse et les idées directrices d'un texte argumentatif.
\item À la fin de cette leçon, je sais relever, définir et réemployer le lexique de la famille hispanique et de la condition féminine (violencia de género, machismo, maltrato, denuncia, igualdad, acoso).
\item À la fin de cette leçon, je sais répondre par des phrases complètes à des questions de compréhension en espagnol.
\item À la fin de cette leçon, je sais présenter un texte (nature, source, auteur, thème) selon la méthode du commentaire.
\item À la fin de cette leçon, je sais dégager deux ou trois idées directrices et les justifier par des citations du texte.
\item À la fin de cette leçon, je sais rédiger un commentaire personnel structuré (avis argumenté, exemples du Cameroun et du monde hispanique).
\item À la fin de cette leçon, je sais produire un court texte engagé (affiche, paragraphe de dénonciation) en réutilisant le lexique étudié.
\end{itemize}
\end{objectifs}

\begin{prerequis}
\begin{itemize}
\item Lexique de base de la famille étudié en Première (padres, esposos, hijos, hogar).
\item Techniques de lecture : repérage du thème et des mots-clés d'un texte (Première).
\item Expression de l'opinion : pienso que, me parece que, estoy de acuerdo con + subjonctif/indicatif.
\item Connecteurs logiques simples : sin embargo, además, por eso, en conclusión.
\item Notions de société : droits de l'homme, égalité, vues en classe de Première et en EMC.
\end{itemize}
\end{prerequis}

\newpage

\coursec{Texte support : lecture et découverte}

\courssub{Situation d'accroche et problématique}

En mars 2023, des milliers de femmes ont défilé dans les rues de Madrid, de Mexico et de Yaoundé pour dire \esp{Ni una menos} (« Pas une de moins »). Banderoles violettes, slogans scandés, chansons : partout, le même message. Au Cameroun comme en Espagne, les statistiques montrent que les violences faites aux femmes restent une réalité quotidienne : harcèlement, mariages forcés, violences conjugales. Or, pour comprendre ce phénomène dans le monde hispanophone, il faut savoir lire la presse espagnole et latino-américaine : comment les journalistes nomment-ils ces réalités ? Quel vocabulaire utilisent-ils ? Quels arguments avancent-ils ?

\textbf{Problématique :} comment la presse hispanophone parle-t-elle des discriminations et des violences contre les femmes, et comment pouvons-nous, à partir d'un article, construire un commentaire de texte digne du Baccalauréat ?

Dans cette première étape de la leçon, nous allons lire un article original, en dégager les informations essentielles, identifier son genre, deviner le sens des mots nouveaux grâce au contexte, puis reformuler son contenu général en deux ou trois phrases.

\courssub{Première lecture : la bonne démarche}

Avant de plonger dans le texte, adoptons la méthode professionnelle du lecteur efficace. Au Baccalauréat, le temps est compté : une première lecture intelligente fait gagner de précieuses minutes.

\begin{methode}[Réussir sa première lecture]
\begin{enumerate}
\item \textbf{Lire le titre} : il annonce le thème et souvent la position de l'auteur. Ici, \esp{Las mujeres siguen luchando} (« Les femmes continuent de lutter ») annonce à la fois un constat (la lutte) et une durée (\esp{seguir} + gérondio = continuer à).
\item \textbf{Identifier la source} : un article de presse n'attend pas le même type de lecture qu'un poème ou qu'un extrait de roman.
\item \textbf{Lire le premier et le dernier paragraphe} : le premier pose le problème, le dernier donne souvent la conclusion ou la solution.
\item \textbf{Lire ensuite le texte entier}, sans dictionnaire, en acceptant de ne pas tout comprendre.
\item \textbf{Souligner les mots répétés} : ce sont les mots-clés du thème. Ici, vous remarquerez que \esp{mujeres}, \esp{violencia}, \esp{igualdad} reviennent plusieurs fois.
\item \textbf{Formuler une première idée générale} en une phrase : « De quoi parle ce texte ? »
\end{enumerate}
\end{methode}

\begin{attention}[Le piège de la traduction mot à mot]
Beaucoup d'élèves commencent par traduire chaque mot inconnu avec le dictionnaire. Erreur ! En compréhension de l'écrit, on lit d'abord \emph{globalement}, on devine, on fait des hypothèses. Le dictionnaire ne sert qu'à \emph{vérifier} ensuite. Un mot inconnu sur dix n'empêche jamais de comprendre un texte.
\end{attention}

\courssub{Lecture silencieuse du texte}

Lisez maintenant silencieusement l'article suivant, paru dans un journal hispanophone. Soulignez au crayon les mots qui vous semblent importants ou inconnus.

\begin{textebox}[Las mujeres siguen luchando — artículo de prensa]
En muchos países hispanos, la violencia de género sigue siendo una realidad cotidiana. Según las estadísticas, una de cada tres mujeres ha sufrido maltrato físico o psicológico en el hogar, muchas veces a manos de su propia pareja. Detrás de las cifras hay rostros, historias y silencios.

El machismo, heredado de una educación tradicional, considera a la mujer inferior al hombre y le niega sus derechos más elementales : el derecho a estudiar, a trabajar, a elegir libremente a su esposo o a decidir sobre su propio cuerpo. En algunas regiones, las niñas abandonan la escuela antes de que los niños lo hagan porque sus familias prefieren invertir en la educación de los hijos varones.

Sin embargo, el problema no se queda en la esfera privada. En la calle y en el trabajo, el acoso sigue siendo frecuente : miradas insistentes, comentarios ofensivos, propuestas humillantes. Muchas víctimas no presentan ninguna denuncia por miedo a las represalias o por vergüenza ante la sociedad. El silencio se convierte así en el mejor aliado del agresor.

Sin embargo, gracias a las leyes de igualdad y al trabajo incansable de las asociaciones y de las ONG, cada vez más mujeres se atreven a hablar, a denunciar y a reclamar justicia. Las redes sociales amplifican sus voces : campañas como \emph{Ni una menos} han movilizado a millones de personas en América Latina y en España.

La lucha contra el acoso y la discriminación es tarea de todos : la familia, la escuela y el Estado deben actuar juntos. Educar a los niños en el respeto, proteger a las víctimas y castigar a los culpables son las tres claves para construir una sociedad más justa, donde hombres y mujeres caminen por fin de la mano, con los mismos derechos y la misma dignidad.
\end{textebox}

\textbf{Traduction du texte (pour vérification après lecture) :}

« Dans de nombreux pays hispaniques, la violence de genre continue d'être une réalité quotidienne. Selon les statistiques, une femme sur trois a subi des mauvais traitements physiques ou psychologiques au foyer, souvent de la part de son propre partenaire. Derrière les chiffres, il y a des visages, des histoires et des silences.

Le machisme, hérité d'une éducation traditionnelle, considère la femme comme inférieure à l'homme et lui refuse ses droits les plus élémentaires : le droit d'étudier, de travailler, de choisir librement son époux ou de décider de son propre corps. Dans certaines régions, les filles abandonnent l'école avant que les garçons ne le fassent parce que leurs familles préfèrent investir dans l'éducation des fils.

Cependant, le problème ne reste pas dans la sphère privée. Dans la rue et au travail, le harcèlement reste fréquent : regards insistants, commentaires offensants, propositions humiliantes. Beaucoup de victimes ne portent pas plainte par peur des représailles ou par honte devant la société. Le silence devient ainsi le meilleur allié de l'agresseur.

Cependant, grâce aux lois sur l'égalité et au travail infatigable des associations et des ONG, de plus en plus de femmes osent parler, dénoncer et réclamer justice. Les réseaux sociaux amplifient leurs voix : des campagnes comme \emph{Ni una menos} ont mobilisé des millions de personnes en Amérique latine et en Espagne.

La lutte contre le harcèlement et la discrimination est la tâche de tous : la famille, l'école et l'État doivent agir ensemble. Éduquer les enfants dans le respect, protéger les victimes et punir les coupables sont les trois clés pour construire une société plus juste, où hommes et femmes marchent enfin main dans la main, avec les mêmes droits et la même dignité. »

\courssub{Lecture à voix haute : prononciation et rythme}

Le professeur lit maintenant le texte à voix haute, puis des élèves lisent chacun un paragraphe. Pendant cette lecture, soyez attentifs aux points de prononciation suivants, typiques des difficultés des francophones :

\begin{itemize}
\item \esp{género}, \esp{psicológico} : la lettre \esp{g} devant \esp{e}/\esp{i} se prononce comme un « r » grasseyé espagnol (son \emph{jota} /x/), jamais « gué ».
\item \esp{vergüenza} : le tréma sur le \esp{ü} indique que le \esp{u} se prononce : « ver-gouèn-za » /beɾˈɣwenθa/. Sans tréma, \esp{gue} se prononcerait « gué » /ge/.
\item \esp{denuncia}, \esp{estadísticas} : respectez l'accent tonique écrit : de-NUN-cia /deˈnunθja/, es-ta-DIS-ti-cas /estaˈðistikas/.
\item \esp{hogar} : le \esp{h} initial est toujours muet en espagnol : « o-gar » /oˈɣaɾ/.
\item \esp{leyes} : le \esp{y} se prononce comme un « y » français dans « payer » /ʝ/.
\end{itemize}

\begin{attention}[Erreurs de prononciation fréquentes des francophones]
\begin{itemize}
\item Ne prononcez jamais le \esp{h} : \esp{hogar}, \esp{hombres}, \esp{hijas} se lisent sans h aspiré.
\item Le \esp{v} espagnol se prononce comme un \esp{b} doux (bilabiale) : \esp{violencia} ne commence pas par le « v » labio-dentaire de « violon ».
\item Le \esp{z} et le \esp{c} devant \esp{e/i} se prononcent « th » anglais (\emph{distinción}) en Espagne (\esp{vergüenza} $\rightarrow$ /beɾˈɣwenθa/), mais « s » (\emph{seseo}) en Amérique latine. Les deux normes sont correctes.
\end{itemize}
\end{attention}

\courssub{Repérage des informations de première lecture}

Après la lecture, répondons aux trois questions fondamentales du premier contact avec un texte.

\textbf{1. ¿De qué habla el texto ? (De quoi parle le texte ?)}

Le texte parle de la \cle{violencia de género} et des discriminations que subissent les femmes dans les pays hispanophones, ainsi que des moyens de lutter contre ce phénomène.

\textbf{2. ¿Quién está concernido ? (Qui est concerné ?)}

Les personnages centraux sont \cle{las mujeres} (les femmes), victimes du machisme et du harcèlement. Mais le texte mentionne aussi d'autres acteurs : la \cle{familia}, l'\cle{escuela} (l'école), l'\cle{Estado} (l'État), les \cle{asociaciones} et les \cle{ONG}.

\textbf{3. ¿Dónde ? (Où ?)}

Le cadre est large : \esp{muchos países hispanos} (de nombreux pays hispaniques), avec des références à l'Amérique latine et à l'Espagne. Les lieux précis évoqués sont le \esp{hogar} (le foyer), la \esp{calle} (la rue), le \esp{trabajo} (le travail) et l'\esp{escuela}.

\begin{exemplebox}[Questions de compréhension orale]
Répondez en espagnol, par des phrases complètes :
\begin{enumerate}
\item ¿Qué tipo de violencia sufren muchas mujeres según el texto ?
\item ¿Por qué muchas víctimas no presentan ninguna denuncia ?
\item ¿Qué tres instituciones deben actuar juntas según el autor ?
\item ¿Qué papel juegan las redes sociales en la lucha de las mujeres ?
\end{enumerate}
\end{exemplebox}

\courssub{Identification du genre du texte}

Observons les indices matériels et linguistiques du texte :

\begin{center}
\begin{tabularx}{\linewidth}{|X|X|}
\hline
\rowcolor{popblueL} \textbf{Indice observé} & \textbf{Ce qu'il révèle} \\
\hline
Titre informatif : \esp{Las mujeres siguen luchando} & Annonce un sujet de société \\
\hline
Référence aux \esp{estadísticas} & Recherche d'objectivité, ton journalistique \\
\hline
Connecteurs d'opinion : \esp{sin embargo} & Le texte ne se contente pas d'informer : il argumente \\
\hline
Phrase finale injonctive : \esp{deben actuar juntos} & Prise de position de l'auteur \\
\hline
Exemples concrets : campagne \esp{Ni una menos} & Ancrage dans l'actualité \\
\hline
\end{tabularx}
\end{center}

\begin{definition}[Article de presse argumentatif]
Un \cle{artículo de prensa} est un texte journalistique qui informe sur un fait d'actualité. Lorsqu'il ne se contente pas de rapporter les faits mais défend une thèse (ici : « la lutte contre les violences est l'affaire de tous »), on parle d'\cle{artículo de opinión} ou de texte \cle{argumentativo}. Ses marques : connecteurs logiques (\esp{sin embargo}, \esp{por eso}), verbes d'opinion et de devoir (\esp{deber}, \esp{reclamar}), exemples et chiffres à l'appui.
\end{definition}

Notre texte est donc un \textbf{article de presse à visée argumentative} : il informe (statistiques, réalités) et il persuade (appel à l'action collective dans le dernier paragraphe).

\courssub{Hypothèses de sens à partir du contexte}

Deviner le sens d'un mot inconnu est une compétence essentielle du Baccalauréat. Voici comment procéder avec trois mots du texte.

\textbf{Exemple 1 :} \esp{maltrato}. Contexte : \esp{una de cada tres mujeres ha sufrido maltrato físico o psicológico}. Les adjectifs « physique ou psychologique » et le verbe \esp{sufrir} (souffrir) indiquent quelque chose de négatif subi par les victimes. Hypothèse : « mauvais traitement ». La famille du mot français « maltraiter » confirme.

\textbf{Exemple 2 :} \esp{vergüenza}. Contexte : \esp{no presentan ninguna denuncia por miedo o por vergüenza}. Le mot est associé à \esp{miedo} (peur) par la coordination \esp{o} : c'est donc un sentiment négatif qui empêche d'agir. Hypothèse : « honte ».

\textbf{Exemple 3 :} \esp{atreverse}. Contexte : \esp{cada vez más mujeres se atreven a hablar}. Le contraste avec le silence des victimes, signalé par \esp{sin embargo}, suggère un changement de comportement courageux. Hypothèse : « oser ».

\begin{methode}[Deviner le sens d'un mot inconnu]
\begin{enumerate}
\item Identifier la nature grammaticale du mot (nom, verbe, adjectif ?).
\item Observer les mots qui l'entourent (verbes, adjectifs, connecteurs).
\item Chercher un mot français de la même famille (attention aux faux amis !).
\item Formuler une hypothèse et la réinsérer dans la phrase : le sens général tient-il ?
\item Vérifier seulement ensuite dans le glossaire ou le dictionnaire.
\end{enumerate}
\end{methode}

\begin{attention}[Faux amis à surveiller dans ce thème]
\begin{itemize}
\item \esp{una denuncia} = une plainte, une dénonciation (au sens judiciaire/administratif, pas « trahison »).
\item \esp{la vergüenza} = la honte (et non « la virginité » ; le verbe \esp{avergonzar} = faire honte / mettre mal à l'aise).
\item \esp{el acoso} = le harcèlement (rien à voir avec « accoster »).
\item \esp{la pareja} = le/la partenaire, le couple (et non « la pareille »).
\item \esp{las represalias} = les représailles (orthographe espagnole : un seul \esp{s} après \esp{a} ; français : double \esp{s}).
\end{itemize}
\end{attention}

\courssub{Glossaire du texte}

Voici les mots difficiles du texte, traduits comme ils le seraient en marge de votre copie.

\begin{definition}[Glossaire — Vocabulario del texto]
el maltrato = le mauvais traitement, la maltraitance\\
el hogar = le foyer, la maison\\
la denuncia = la plainte, la dénonciation\\
el miedo = la peur\\
la vergüenza = la honte\\
la ley = la loi\\
atreverse a + infinitif = oser faire quelque chose\\
la tarea = la tâche\\
juntos = ensemble\\
el acoso = le harcèlement\\
las represalias = les représailles\\
reclamar = réclamer, exiger
\end{definition}

\begin{propriete}[Quelques familles de mots utiles]
\begin{itemize}
\item \esp{maltrato} (nom) $\rightarrow$ \esp{maltratar} (verbe : maltraiter) $\rightarrow$ \esp{maltratador/a} (celui/celle qui maltraite).
\item \esp{denuncia} (nom) $\rightarrow$ \esp{denunciar} (verbe : dénoncer, porter plainte).
\item \esp{igualdad} (nom : égalité) $\rightarrow$ \esp{igual} (adjectif : égal) $\rightarrow$ \esp{desigualdad} (inégalité).
\item \esp{miedo} (nom) $\rightarrow$ \esp{tener miedo} (avoir peur ; attention : \esp{tener}, pas \esp{ser} ni \esp{estar}).
\end{itemize}
\end{propriete}

\courssub{Carte mentale du thème}

Pour mémoriser le vocabulaire et organiser vos idées avant le commentaire, voici une carte mentale du thème central du texte.

\begin{popfigure}
\begin{tikzpicture}[mindmap, grow cyclic,
  every node/.style={concept, text=white, font=\small, execute at begin node=\hskip0pt},
  root concept/.append style={fill=popink, font=\small\bfseries},
  level 1/.append style={level distance=3.2cm, sibling angle=90},
  level 2/.append style={level distance=2.2cm, sibling angle=45, font=\scriptsize}]
\node[root concept]{violencia de género}
  child[concept color=poporange]{ node {causas}
    child { node {machismo} }
    child { node {educación tradicional} }
  }
  child[concept color=popdark]{ node {conséquences}
    child { node {silencio} }
    child { node {miedo y vergüenza} }
  }
  child[concept color=popgreen]{ node {solutions}
    child { node {leyes de igualdad} }
    child { node {denuncia} }
  }
  child[concept color=popblue]{ node {acteurs}
    child { node {familia y escuela} }
    child { node {Estado y ONG} }
  };
\end{tikzpicture}
\legende{Carte mentale lexicale : organiser le thème \esp{violencia de género} en quatre branches.}
\end{popfigure}

\courssub{Reformulation orale : el resumen}

Dernière étape de la découverte : être capable de résumer le texte en deux ou trois phrases, à l'oral, sans le regarder. C'est l'exercice du \cle{resumen}, première étape de tout commentaire de texte.

\begin{exoresolu}[Rédiger un resumen du texte]
Énoncé : Résumez l'article \esp{Las mujeres siguen luchando} en deux ou trois phrases en espagnol, sans recopier de phrase entière du texte.
\tcblower
Solution détaillée : un bon resumen reprend le thème, le problème et la conclusion, avec ses propres mots et des connecteurs.

\esp{El artículo trata de la violencia de género en los países hispanos: muchas mujeres sufren maltrato y acoso, pero pocas se atreven a denunciar. Sin embargo, gracias a las leyes y a las asociaciones, la situación empieza a cambiar. El autor concluye que la familia, la escuela y el Estado deben actuar juntos.}

Traduction : « L'article traite de la violence de genre dans les pays hispaniques : beaucoup de femmes subissent maltraitance et harcèlement, mais peu osent porter plainte. Cependant, grâce aux lois et aux associations, la situation commence à changer. L'auteur conclut que la famille, l'école et l'État doivent agir ensemble. »

Remarquez les outils du resumen : \esp{trata de} (il s'agit de), \esp{sin embargo} (cependant), \esp{el autor concluye que} (l'auteur conclut que).
\end{exoresolu}

\begin{exercice}[À vous de jouer]
\begin{enumerate}
\item Retrouvez dans le texte trois mots de la famille de \esp{igual} et traduisez-les.
\item Devinez le sens de \esp{rostro} (paragraphe 1) à partir du contexte, puis vérifiez.
\item Reformulez oralement le contenu du texte en deux phrases, en commençant par \esp{Según el artículo...}
\item Classez ces mots dans la carte mentale (causas, consecuencias, soluciones, actores) : \esp{el machismo}, \esp{la denuncia}, \esp{la escuela}, \esp{el silencio}, \esp{las leyes}, \esp{la vergüenza}.
\end{enumerate}
\end{exercice}

\begin{corrige}[Exercice : À vous de jouer]
\begin{enumerate}
\item \esp{igualdad} (l'égalité), \esp{igual} (égal), et l'on peut ajouter \esp{desigualdad} (l'inégalité).
\item \esp{Detrás de las cifras hay rostros, historias y silencios} : associé à « histoires » et opposé aux « chiffres » abstraits, \esp{rostro} désigne quelque chose de concret et humain : le visage.
\item Modèle : \esp{Según el artículo, la violencia contra las mujeres sigue siendo un problema grave en el mundo hispano. Pero las leyes, las asociaciones y la educación pueden cambiar esta realidad.}
\item Causas : \esp{el machismo}. Conséquences : \esp{el silencio}, \esp{la vergüenza}. Solutions : \esp{la denuncia}, \esp{las leyes}. Acteurs : \esp{la escuela}.
\end{enumerate}
\end{corrige}

\begin{savaistu}[Le 016 et le 8 de marzo]
En Espagne, le numéro d'urgence \textbf{016} est réservé aux victimes de violences de genre : gratuit, confidentiel, il ne laisse aucune trace sur la facture téléphonique. Le \textbf{8 mars} est la \esp{Jornada Internacional de la Mujer} (\esp{Día Internacional de la Mujer}) : ce jour-là, des manifestations (\esp{manifestaciones}) rassemblent des millions de personnes dans tout le monde hispanophone, de Madrid à Buenos Aires, et le mouvement \esp{Ni una menos}, né en Argentine en 2015, est devenu un symbole mondial de la lutte contre les féminicides. Au Cameroun, la Journée internationale de la femme est aussi célébrée chaque 8 mars avec des défilés et des pagnes spécialement confectionnés : un beau point de rencontre entre nos cultures !
\end{savaistu}

\begin{aretenir}[L'essentiel de cette étape]
\begin{itemize}
\item Première lecture : titre, source, premier et dernier paragraphe, puis lecture entière en soulignant les mots répétés.
\item Trois questions de repérage : ¿De qué habla el texto? ¿Quién? ¿Dónde?
\item Notre texte est un \textbf{article de presse argumentatif} : il informe et persuade.
\item Un mot inconnu se devine par le contexte (nature, entourage, famille de mots) avant toute vérification.
\item Le \esp{resumen} reformule en deux ou trois phrases : thème + problème + conclusion, avec des connecteurs (\esp{trata de}, \esp{sin embargo}, \esp{el autor concluye que}).
\end{itemize}
\end{aretenir}

\coursec{Texte support : lecture et découverte}

\courssub{Anticipation et mise en contexte}

\begin{methode}[Avant la lecture : activer ses connaissances]
\begin{enumerate}
    \item \textbf{Observez le titre} : \esp{« Romper el silencio : la lucha contra el machismo »}. Quels mots-clés identifiez-vous ? \cle{Romper el silencio}, \cle{lucha}, \cle{machismo}.
    \item \textbf{Identifiez le genre textuel} : Article d'opinion / Tribune / Article de presse engagée. Source probable : journal espagnol (\esp{El País}, \esp{El Mundo}) ou ONG (\esp{Amnistía Internacional}).
    \item \textbf{Formulez des hypothèses} sur le contenu : statistiques, causes, conséquences, solutions, appel à l'action.
    \item \textbf{Mobilisez le lexique connu} : \esp{violencia de género, denuncia, igualdad, machismo, maltrato, acoso, protección, ley}.
\end{enumerate}
\end{methode}

\begin{experience}[Activité orale : « Lluvia de ideas » (Brainstorming)]
En classe entière, au tableau, construisez une carte mentale autour de \cle{machismo} et \cle{violencia de género}.
\begin{itemize}
    \item Branches : \textbf{Causes} (éducation, stéréotypes, patriarcat), \textbf{Manifestations} (contrôle, insultes, coups, féminicide), \textbf{Conséquences} (peur, silence, mort, cycle intergénérationnel), \textbf{Solutions} (éducation, lois, juges, refuges, dénonciation).
    \item L'enseignant note en espagnol ; les élèves copient dans leur cahier de lexique.
\end{itemize}
\emph{Objectif :} Activer le schème cognitif, réduire l'anxiété lexicale, créer une « banque de mots » réutilisable pour la compréhension et l'expression.
\end{experience}

\courssub{Première lecture globale (lecture cursive)}

\begin{methode}[Consignes de lecture silencieuse (10 min)]
\begin{enumerate}
    \item Lisez le texte \textbf{une première fois sans dictionnaire} pour saisir le sens global.
    \item Soulignez \textbf{les mots transparents} (cognats) : \esp{violencia, familia, sistema, dominación, inferior, maltrato, silencio, víctimas, psicológico, agresión, educación, estereotipos, recursos, protección, voluntad, política, rigor, igualdad, derecho, humano, innegociable}.
    \item Entourez \textbf{les 5 à 7 mots inconnus indispensables} (ex. : \esp{lacra, macroencuesta, expareja, revictimiza, coeducativa, desmontar}).
    \item Répondez oralement ou par écrit (brouillon) aux questions « repères » :
    \begin{itemize}
        \item De quoi parle le texte ? (Thème)
        \item Quel est le ton ? (Indigné, urgent, analytique, espoir)
        \item Quelle est la structure visible ? (3 paragraphes : Constat $\rightarrow$ Causes/Silence $\rightarrow$ Solutions/Appel)
    \end{itemize}
\end{enumerate}
\end{methode}

\begin{textebox}[Texte support : « Romper el silencio : la lucha contra el machismo » (Original)]
En la familia hispánica, la \textbf{violencia de género} sigue siendo una \textbf{lacra} invisible que destruye vidas detrás de puertas cerradas. Según la \textbf{macroencuesta de 2023}, una de cada tres mujeres mayores de 16 años ha sufrido violencia física o sexual en algún momento de su vida; la mayoría, a manos de su \textbf{expareja}. El \textbf{machismo}, entendido no como un acto aislado sino como un \textbf{sistema de dominación}, considera a la mujer \textbf{inferior} al hombre y justifica el \textbf{maltrato} como herramienta de control.

Sin embargo, el \textbf{silencio} sigue siendo el mayor cómplice. El 70\% de las víctimas no presenta \textbf{denuncia} por \textbf{miedo} a las represalias, por \textbf{vergüenza}, por dependencia económica o porque no confían en una justicia que a menudo las \textbf{revictimiza}. El \textbf{acoso} psicológico —insultos, control del móvil, aislamiento familiar— precede casi siempre a la agresión física, pero es más difícil de probar.

La solución no es solo penal. Urge una \textbf{educación coeducativa} desde la escuela que \textbf{desmonte} los estereotipos: los niños no son « fuertes », las niñas no son « cuidadoras ». Faltan \textbf{juzgados especializados}, faltan recursos para la \textbf{protección} inmediata y falta voluntad política para aplicar la \textbf{Ley Orgánica 1/2004} con rigor. Romper el silencio es el primer paso: \textbf{denunciar}, \textbf{acompañar}, \textbf{educar}. La \textbf{igualdad} no es un favor, es un derecho humano \textbf{innegociable}.
\legende{Texte original d'environ 220 mots, niveau Terminale, vocabulaire ciblé Módulo 1.}
\end{textebox}

\courssub{Deuxième lecture : repérage lexical et structurel}

\begin{exercice}[Repérage guidé (travail individuel puis mise en commun)]
Sur le texte (version papier ou annotée sur tablette) :
\begin{enumerate}
    \item Soulignez \textbf{les connecteurs logiques} : \esp{Sin embargo, Según, ;, pero, no es solo, Urge que, La solución...}.
    \item Mettez en \textbf{gras (ou surligneur couleur 1)} les \textbf{noms du champ lexical de la violence/oppression} : \esp{violencia de género, lacra, machismo, sistema de dominación, inferior, maltrato, control, silencio, cómplice, denuncia, miedo, vergüenza, revictimiza, acoso, agresión}.
    \item Mettez en \textbf{italique (ou surligneur couleur 2)} les \textbf{noms/verbes des solutions/droits} : \esp{educación coeducativa, desmontar, juzgados especializados, protección, voluntad política, aplicar, romper, denunciar, acompañar, educar, igualdad, derecho, innegociable}.
    \item Identifiez \textbf{la thèse} (phrase résumant l'opinion de l'auteur) : \esp{« La igualdad no es un favor, es un derecho humano innegociable »} (dernière phrase) ou \esp{« Urge una educación coeducativa... »}.
\end{enumerate}
\end{exercice}

\begin{definition}[Glossaire indispensable du texte (mots opaques / spécialisés)]
\begin{itemize}
    \item \esp{lacra} (nf) = fléau, plaie sociale, tare (orig. : marque au fer rouge sur le bétail).
    \item \esp{macroencuesta} (nf) = enquête nationale à grande échelle (macro- + enquête).
    \item \esp{expareja} (nf) = ex-partenaire, ex-conjoint (préfixe \esp{ex-} + \esp{pareja}).
    \item \esp{sistema de dominación} (nm) = système de domination (concept sociologique).
    \item \esp{revictimizar} (v) = revictimiser : faire subir une seconde victimisation (ex. lors de l'instruction judiciaire).
    \item \esp{coeducativa} (adj) = coéducatif/ve : éducation mixte, non sexiste, égalitaire.
    \item \esp{desmontar} (v) = démonter, déconstruire (ici : déconstruire des idées reçues).
    \item \esp{innegociable} (adj) = non négociable (préfixe \esp{in-} + \esp{negociable}).
    \item \esp{Ley Orgánica 1/2004} = Loi organique espagnole « Medidas de protección integral contra la violencia de género » (référence juridique clé).
\end{itemize}
\end{definition}

\coursec{Compréhension du texte}

\courssub{Compréhension globale : thème, idée principale, intention de l'auteur}

\begin{methode}[Comment répondre aux questions de compréhension globale]
\begin{enumerate}
    \item \textbf{Lisez le titre et les paratextes} (source, date, auteur) : ils orientent le thème.
    \item \textbf{Identifiez le thème} (sujet général) en un mot ou une expression nominale : \esp{la violencia de género}, \esp{el machismo en la familia}, \esp{la igualdad de derechos}.
    \item \textbf{Formulez l'idée principale} (le message central) en une phrase affirmative : \esp{El texto denuncia la persistencia del machismo y exige una educación para la igualdad.}
    \item \textbf{Précisez l'intention de l'auteur} : \esp{informar}, \esp{denunciar}, \esp{concienciar}, \esp{exigir medidas}, \esp{romper el silencio}.
    \item \textbf{Rédigez la réponse en espagnol}, phrase complète, en reprenant les mots-clés de la question.
\end{enumerate}
\end{methode}

\begin{exemplebox}[Exemples de réponses globales (d'après le texte support)]
\textbf{Pregunta :} ¿Cuál es el tema principal del texto? \\
\textbf{Respuesta :} El tema principal del texto es \textbf{la violencia de género y la discriminación contra la mujer en el ámbito familiar y social}. \\[4pt]
\textbf{Pregunta :} ¿Cuál es la idea principal que defiende el autor? \\
\textbf{Respuesta :} La idea principal es que \textbf{el machismo estructural sigue causando sufrimiento a las mujeres y que solo la educación y la denuncia pueden lograr la igualdad real}. \\[4pt]
\textbf{Pregunta :} ¿Cuál es la intención del autor al escribir este artículo? \\
\textbf{Respuesta :} La intención del autor es \textbf{denunciar la situación de las víctimas, concienciar a la sociedad y exigir la aplicación efectiva de las leyes de protección}.
\end{exemplebox}

\courssub{Compréhension fine : repérage, citation, reformulation}

Les questions de détail portent sur des informations explicites (chiffres, causes, conséquences, solutions) ou implicites (inférences logiques). La règle d'or au Bac : \textbf{jamais de « Sí/No » sec, jamais de réponse en français, jamais de paraphrase vague sans ancre textuelle}.

\begin{propriete}[Technique de la réponse justifiée (Tipo Bac)]
Pour toute question de compréhension détaillée, suivez ce protocole en 4 étapes :
\begin{enumerate}
    \item \textbf{Repérez} le paragraphe ou la phrase exacte qui contient la réponse.
    \item \textbf{Citez} le fragment pertinent \textbf{entre guillemets espagnols} (\esp{« ... »}) — pas tout le paragraphe, seulement le segment nécessaire.
    \item \textbf{Reformulez} en espagnol avec vos mots (synonymes, changement de structure) pour prouver que vous avez compris.
    \item \textbf{Rédigez une phrase complète} qui reprend le sujet de la question.
\end{enumerate}
\end{propriete}

\begin{popfigure}
\begin{tikzpicture}[
    node distance=1.2cm and 1.5cm,
    every node/.style={font=\small, align=center},
    box/.style={draw=popblue, thick, rounded corners, fill=popblueL, minimum width=3.2cm, minimum height=1.4cm},
    arrow/.style={-Stealth, thick, popdark}
]
\node[box, align=center] (q){\textbf{Pregunta}\\\emph{(Leer + Analizar)}};
\node[box, right=of q, align=center] (loc){\textbf{Localizar}\\\emph{(Subrayar en el texto)}};
\node[box, right=of loc, align=center] (cit){\textbf{Citar}\\\emph{(« Fragmento exacto »)}};
\node[box, right=of cit, align=center] (ans){\textbf{Responder}\\\emph{(Frase completa + reforma)}};
\draw[arrow] (q) -- (loc);
\draw[arrow] (loc) -- (cit);
\draw[arrow] (cit) -- (ans);
\node[below=0.3cm of q, text width=3cm, font=\tiny, popmuted] {¿Por qué...?\newline¿Qué cifra...?};
\node[below=0.3cm of loc, text width=3cm, font=\tiny, popmuted] {Párrafo 2, línea 3};
\node[below=0.3cm of cit, text width=3cm, font=\tiny, popmuted] {\esp{« por miedo o vergüenza »}};
\node[below=0.3cm of ans, text width=3cm, font=\tiny, popmuted] {\esp{No denuncian por temor y vergüenza.}};
\end{tikzpicture}
\legende{Organigramme : la chaîne de traitement d'une question de compréhension (méthode type Bac).}
\end{popfigure}

\begin{methode}[Rédiger la réponse complète : modèles de phrases]
Utilisez ces structures pour commencer votre phrase en reprenant la question :
\begin{itemize}
    \item \textbf{Question \esp{¿Por qué...?}} $\rightarrow$ \esp{La razón es que... / Se debe a que... / Porque...}
    \item \textbf{Question \esp{¿Qué...? / ¿Cuál...?}} $\rightarrow$ \esp{Es... / Se trata de... / Lo que ocurre es que...}
    \item \textbf{Question \esp{¿Cómo...?}} $\rightarrow$ \esp{Se hace... / Lo hace de la siguiente manera: ...}
    \item \textbf{Question \esp{Según el texto...}} $\rightarrow$ \esp{Según el texto, ... / El texto afirma que... / El autor señala que...}
\end{itemize}
\end{methode}

\begin{exemplebox}[Questions fines et réponses modèles (basées sur le texte support)]
\textbf{1. Chiffres / Données précises} \\
\textbf{Pregunta :} ¿Qué porcentaje de mujeres ha sufrido violencia física o sexual según los datos citados? \\
\textbf{Respuesta :} \textbf{Según el texto, « una de cada tres mujeres » ha sufrido violencia física o sexual.} (Reformulation : \esp{Aproximadamente el 33\% de las mujeres...}) \\[4pt]
\textbf{2. Causes du machisme} \\
\textbf{Pregunta :} ¿Cuáles son las causas profundas del machismo que menciona el autor? \\
\textbf{Respuesta :} \textbf{El autor señala que el machismo tiene su origen en « una educación patriarcal » y en « estereotipos de género transmitidos desde la infancia ».} \\[4pt]
\textbf{3. Raisons du silence des victimes} \\
\textbf{Pregunta :} ¿Por qué muchas víctimas no denuncian a sus agresores? \\
\textbf{Respuesta :} \textbf{Muchas víctimas no denuncian « por miedo a las represalias » y « por vergüenza o culpa internalizada ».} \\[4pt]
\textbf{4. Solutions proposées} \\
\textbf{Pregunta :} ¿Qué medidas propone el texto para erradicar la violencia de género? \\
\textbf{Respuesta :} \textbf{El texto propone « una educación coeducativa desde la escuela », « el fortalecimiento de los juzgados especializados » y « campañas de concienciación social ».}
\end{exemplebox}

\begin{attention}[Erreurs fréquentes des francophones — \textbf{À éviter absolument au Bac}]
\begin{itemize}
    \item \textbf{Répondre en français.} \emph{Consigne :} Toute réponse doit être rédigée \textbf{en espagnol}. Zéro point si la réponse est en français, même si l'idée est juste.
    \item \textbf{Répondre par « Sí » ou « No » sans justification.} \emph{Consigne :} Les questions « ¿Cree usted que...? », « ¿Está de acuerdo...? » exigent une justification textuelle. \esp{Sí, porque el texto dice que...}
    \item \textbf{Recopier tout le paragraphe.} \emph{Consigne :} Citez \textbf{uniquement le fragment utile} (3 à 10 mots max). Recopier 5 lignes montre que vous n'avez pas repéré l'info précise.
    \item \textbf{Ne pas mettre de guillemets} ou utiliser des guillemets français \guillemotleft \guillemotright. \emph{Consigne :} Guillemets espagnols obligatoires : \esp{« ... »}.
    \item \textbf{Confondre « idea principal » et « tema ».} \emph{Rappel :} Le thème = nom (la violence) ; l'idée principale = phrase (la violence persiste à cause de l'éducation).
    \item \textbf{Calque syntaxique :} \esp{*La víctima no denuncia por el miedo.} $\rightarrow$ \textbf{Correct :} \esp{La víctima no denuncia por miedo} (sans article devant nom abstrait après \esp{por}).
    \item \textbf{Confusion \esp{ser}/\esp{estar} :} \esp{*La mujer es inferior.} (Caractère essentiel/identité $\rightarrow$ \esp{ser}). \esp{*La víctima está asustada.} (État passager $\rightarrow$ \esp{estar}). Le texte utilise \esp{considera a la mujer inferior} (verbe + attribut $\rightarrow$ \esp{ser} implicite).
    \item \textbf{Subjonctif après \esp{urge que}, \esp{falta que}, \esp{es necesario que} :} \esp{Urge una educación que desmonte} (subjonctif), \esp{Faltan juzgados que protejan} (subjonctif).
\end{itemize}
\end{attention}

\courssub{Entraînement type Bac : Vrai / Faux justifié et recherche lexicale}

\begin{propriete}[Exercice « Verdadero / Falso » justifié (Modalité Bac)]
\begin{enumerate}
    \item Lisez l'affirmation et décidez si elle est \textbf{Verdadera (V)} ou \textbf{Falsa (F)} d'après \textbf{exclusivement le texte}.
    \item \textbf{Justifiez toujours} en citant le passage exact (\esp{« ... »}) et en expliquant brièvement.
    \item Si c'est \textbf{Falso}, corrigez l'information fausse avec la bonne info du texte.
    \item Réponse en espagnol, phrase complète.
\end{enumerate}
\end{propriete}

\begin{exemplebox}[Modèles V / F justifiés]
\textbf{Afirmación :} « El machismo considera a la mujer superior al hombre. » \\
\textbf{Respuesta :} \textbf{Falso.} El texto afirma lo contrario: \esp{« El machismo considera a la mujer inferior al hombre »}. \\[4pt]
\textbf{Afirmación :} « Las leyes actuales son suficientes para proteger a las víctimas. » \\
\textbf{Respuesta :} \textbf{Falso.} El autor denuncia que \esp{« la aplicación de las leyes sigue siendo insuficiente »} y que \esp{« faltan recursos en los juzgados »}. \\[4pt]
\textbf{Afirmación :} « La educación es la clave para prevenir la violencia. » \\
\textbf{Respuesta :} \textbf{Verdadero.} El texto concluye que \esp{« solo una educación en igualdad desde la infancia puede romper el ciclo de la violencia »}.
\end{exemplebox}

\begin{propriete}[Recherche d'équivalents et de contraires dans le texte]
\begin{itemize}
    \item \textbf{Équivalent (sinónimo) :} Trouvez dans le texte le mot/l'expression qui a le \textbf{même sens} que le mot donné, dans le même contexte.
    \item \textbf{Contraire (antónimo) :} Trouvez dans le texte le mot/l'expression de \textbf{sens opposé}.
    \item \textbf{Règle :} Recopiez \textbf{exactement la forme du texte} (masculin/féminin, singulier/pluriel, temps verbal). Citez entre guillemets.
\end{itemize}
\end{propriete}

\begin{center}
\begin{tabularx}{\linewidth}{|X|X|X|}
\hline
\rowcolor{popblueL}
\textbf{Mot donné (consigne)} & \textbf{Type} & \textbf{Réponse attendue (citation texte)} \\
\hline
\esp{igualdad} & Contraire & \esp{« discriminación »} / \esp{« desigualdad »} \\
\hline
\esp{maltrato} & Équivalent & \esp{« violencia »} / \esp{« abuso »} \\
\hline
\esp{denunciar} & Équivalent & \esp{« reportar »} / \esp{« acusar »} \\
\hline
\esp{silencio} & Contraire & \esp{« denuncia »} / \esp{« voz »} / \esp{« grito »} \\
\hline
\esp{víctima} & Équivalent & \esp{« mujer maltratada »} / \esp{« afectada »} \\
\hline
\esp{sumiso / sumisa} & Contraire & \esp{« rebelde »} / \esp{« empoderada »} (si dans texte) \\
\hline
\end{tabularx}
\end{center}

\courssub{Exercice résolu complet : simulation d'épreuve}

\begin{exoresolu}[Épreuve de compréhension type Bac (sur le texte ci-dessus)]
\textbf{Consignes :} Répondez en espagnol, phrases complètes, justifiez par une citation courte (\esp{« ... »}).

\textbf{1. Compréhension globale} \\
\textbf{a)} ¿Cuál es el tema central del artículo? \\
\textbf{b)} ¿Cuál es la tesis principal que defiende el autor? \\
\textbf{c)} ¿Con qué intención ha escrito el autor este texto?

\textbf{2. Compréhension détaillée} \\
\textbf{a)} ¿Qué cifra alarmante da la macroencuesta de 2023 sobre la violencia contra las mujeres? \\
\textbf{b)} Según el texto, ¿por qué el 70\% de las víctimas no denuncia? (Citez au moins trois raisons). \\
\textbf{c)} ¿Qué papel juega el acoso psicológico en el ciclo de la violencia? \\
\textbf{d)} ¿Qué tres medidas concretas exige el autor para en finir con la violence?

\textbf{3. Vrai / Faux justifié} \\
\textbf{a)} El machismo se define en el texto como un acto individual y esporádico. \\
\textbf{b)} La ley actual (Ley Orgánica 1/2004) se aplica con total eficacia según el autor. \\
\textbf{c)} La educación coeducativa se presenta como una medida preventiva fundamental.

\textbf{4. Lexique dans le texte} \\
Trouvez dans le texte : \\
\textbf{a)} Un synonyme de \esp{flaqueza / debilidad} (au sens de « plaie sociale », 1er paragraphe). \\
\textbf{b)} Un antonyme de \esp{superior} (1er paragraphe). \\
\textbf{c)} Un synonyme de \esp{denunciar} (dernier paragraphe). \\
\textbf{d)} L'antonyme de \esp{silencio} (dernier paragraphe, forme verbale).
\tcblower
\textbf{SOLUTION DÉTAILLÉE}

\textbf{1. Compréhension globale} \\
\textbf{a)} \textbf{El tema central del artículo es la violencia de género en el ámbito familiar y la persistencia del machismo como sistema de dominación.} \\
\textbf{b)} \textbf{La tesis principal es que el machismo estructural causa violencia contra las mujeres y que solo la educación en igualdad, la denuncia y una justicia eficaz pueden erradicarla.} \\
\textbf{c)} \textbf{La intención del autor es denunciar la impunidad y el silencio, concienciar sobre las causas estructurales y exigir medidas políticas y educativas urgentes.}

\textbf{2. Compréhension détaillée} \\
\textbf{a)} \textbf{La macroencuesta de 2023 revela que « una de cada tres mujeres mayores de 16 años ha sufrido violencia física o sexual ».} \\
\textbf{b)} \textbf{Las víctimas no denuncian « por miedo a las represalias », « por vergüenza », « por dependencia económica » y « porque no confían en una justicia que a menudo las revictimiza ».} \\
\textbf{c)} \textbf{El acoso psicológico « precede casi siempre a la agresión física » y incluye insultos, control del móvil y aislamiento familiar, aunque « es más difícil de probar ».} \\
\textbf{d)} \textbf{El autor exige: 1) « una educación coeducativa desde la escuela », 2) « juzgados especializados » con recursos, 3) « protección inmediata » y aplicación rigurosa de la Ley Orgánica 1/2004.}

\textbf{3. Vrai / Faux justifié} \\
\textbf{a)} \textbf{Falso.} El texto define el machismo como « un sistema de dominación », no como un acto aislado: \esp{« entendido no como un acto aislado sino como un sistema de dominación »}. \\
\textbf{b)} \textbf{Falso.} El autor critica su aplicación: \esp{« falta voluntad política para aplicar la Ley Orgánica 1/2004 con rigor »}. \\
\textbf{c)} \textbf{Verdadero.} El texto afirma que urge \esp{« una educación coeducativa desde la escuela que desmonte los estereotipos »}.

\textbf{4. Lexique dans le texte} \\
\textbf{a)} \esp{« lacra »} (synonyme de fléau/faiblesse structurelle). \\
\textbf{b)} \esp{« inferior »} (antonyme de superior : \esp{« considera a la mujer inferior al hombre »}). \\
\textbf{c)} \esp{« denunciar »} (répété au dernier paragraphe : \esp{« Romper el silencio es el primer paso: denunciar... »}). \\
\textbf{d)} \esp{« denunciar »} (verbe d'action opposée au silence : \esp{« Romper el silencio... denunciar »}). \textit{Note : « alzar la voz » serait aussi acceptable si présent, mais ici c'est le verbe « denunciar » qui structure l'opposition.}
\end{exoresolu}

\courssub{Tableau récapitulatif : Types de questions / Stratégies de réponse}

\begin{center}
\begin{tabularx}{\linewidth}{|>{\bfseries}l|X|X|}
\hline
\rowcolor{popblueL}
\textbf{Tipo de pregunta} & \textbf{Estrategia clave} & \textbf{Inicio de respuesta modelo} \\
\hline
Tema / Idea principal / Intención & Lectura global + síntesis & \esp{El tema es... / La idea principal es... / La intención es...} \\
\hline
Dato explícito (cifra, nombre, fecha) & Repérage visuel (scan) + citation exacte & \textbf{Según el texto,} \esp{la cifra es... / el dato es...} \\
\hline
Causa (\esp{¿Por qué?}) & Chercher connecteurs (\esp{por, porque, debido a, causa}) & \esp{La razón es que... / Se debe a que...} \\
\hline
Conséquence / Solution & Chercher \esp{por tanto, por eso, solución, propone, urge} & \esp{Como consecuencia... / La solución propuesta es...} \\
\hline
Verdadero / Falso & Décision + Citation obligatoire + Explication & \esp{Verdadero. El texto dice: « ... » / Falso. El texto afirma lo contrario: « ... »} \\
\hline
Sinónimo / Antónimo & Contexte immédiat + même catégorie grammaticale & \esp{El sinónimo de « X » es « Y » / El antónimo es « Z »} \\
\hline
Inférence (opinion auteur) & Indices : adjectifs, modalité (\esp{urge, falta, lamentablemente}) & \esp{El autor considera que... / Se deduce que...} \\
\hline
\end{tabularx}
\end{center}

\begin{experience}[Activité de classe : « Juego de roles : Entrevista a la experta »]
\begin{enumerate}
    \item Par binômes : l'un est \textbf{la periodista}, l'autre \textbf{la experta en igualdad} (psychologue, avocate, présidente d'association).
    \item La journaliste pose 5 questions basées sur le texte (globales + détaillées + opinion).
    \item L'experta répond \textbf{uniquement en s'appuyant sur le texte} (citations + reformulation) + une ouverture personnelle brève.
    \item Enregistrement audio (1 min 30) pour auto-évaluation : \emph{Ai-je cité ? Ai-je reformulé ? Mon espagnol est-il correct (ser/estar, subjonctif si opinion) ?}
    \item Échange de rôles.
\end{enumerate}
\emph{Objectif :} Oraliser la compréhension, s'entraîner à la citation orale (\esp{« Como dice el texto... »}), travailler la fluidité.
\end{experience}

\begin{savaistu}[Le contexte camerounais et la Guinée équatoriale]
Au Cameroun, la \textbf{Loi n° 2016/007 du 12 juillet 2016} portant Code pénal punit les violences conjugales (art. 356-357) et le harcèlement sexuel (art. 299-1). Cependant, le \emph{silence} reste fréquent par peur du stigmate, dépendance économique ou poids des traditions (dot, polygamie). En \textbf{Guinée équatoriale} (seul pays hispanophone d'Afrique), un cadre juridique inspiré de la \textbf{Ley Orgánica 1/2004} espagnole existe mais son application bute sur les mêmes obstacles : manque de juges spécialisés, pesanteur clanique, impunité des puissants. Le texte étudié résonne donc doublement pour vous : vocabulaire juridique commun (\esp{denuncia, juzgado, ley orgánica, protección}) et réalité sociologique partagée. \textbf{Le mot \cle{machismo} n'est pas un mot « espagnol d'Espagne » : il décrit une réalité vécue à Yaoundé, Douala, Malabo ou Bata.}
\end{savaistu}

\coursec{Lexique thématique : la familia y la condición de la mujer}

\courssub{Le noyau familial : vocabulaire de base et évolutions}

\begin{definition}[La familia : typologie et membres]
\begin{center}
\begin{tabularx}{\linewidth}{|X|X|X|}
\hline
\rowcolor{popblueL}
\textbf{Français} & \textbf{Español} & \textbf{Remarques / Exemples} \\
\hline
Famille nucléaire & \esp{familia nuclear} & Padre, madre, hijos. Modèle « traditionnel ». \\
\hline
Famille monoparentale & \esp{familia monoparental} & Un seul parent (\esp{madre soltera / padre soltero}). En hausse. \\
\hline
Famille recomposée & \esp{familia reconstituida} & \esp{Padrastro, madrastra, hermanastros}. \\
\hline
Famille élargie & \esp{familia extensa} & Abuelos, tíos, primos cohabitant. Fréquente en Amérique latine / Afrique. \\
\hline
Couple non marié & \esp{pareja de hecho} & Union libre, reconnaissance légale variable selon pays. \\
\hline
Mariage / Divorce & \esp{matrimonio / divorcio} & \esp{Casarse, separarse, divorciarse}. \\
\hline
Violences conjugales & \esp{violencia doméstica / de género} & \esp{Maltrato} (terme large), \esp{maltrato habitual} (délit spécifique Code pénal ES). \\
\hline
\end{tabularx}
\end{center}
\end{definition}

\courssub{Champ lexical de la violence de genre et de la protection}

\begin{center}
\begin{tabularx}{\linewidth}{|X|X|X|}
\hline
\rowcolor{poporangeL}
\textbf{Concepto clave} & \textbf{Definición / Contexto} & \textbf{Colocaciones frecuentes} \\
\hline
\esp{Violencia de género} & Violence structurelle hommes$\rightarrow$femmes (loi ES 2004). & \esp{ejercer / sufrir / denunciar / erradicar la violencia de género} \\
\hline
\esp{Machismo} & Idéologie de supériorité masculine / domination patriarcale. & \esp{machismo estructural, cultural, violento; combatir el machismo} \\
\hline
\esp{Maltrato} & Mauvais traitements (physiques, psychiques, économiques). & \esp{maltrato físico / psicológico / económico; víctima de maltrato} \\
\hline
\esp{Acoso} & Harcèlement (insultes, contrôle, surveillance, stalking). & \esp{acoso psicológico / sexual / callejero / ciberacoso} \\
\hline
\esp{Feminicidio / Femicidio} & Meurtre de femme \emph{parce que} femme (motif genre). & \esp{alarmante aumento de feminicidios; ley contra el feminicidio} \\
\hline
\esp{Denuncia} & Signalement officiel (police, juge). Acte juridique. & \esp{presentar / interponer / retirar una denuncia; denuncia falsa} \\
\hline
\esp{Orden de protección} & Mesure judiciaire d'éloignement / interdiction approche. & \esp{solicitar / conceder / violar una orden de protección} \\
\hline
\esp{Juzgado de Violencia sobre la Mujer} & Tribunal spécialisé (Espagne). Modèle exporté. & \esp{juzgado especializado, juez de guardia, fiscalía} \\
\hline
\esp{Revictimización} & Second traumatisme (procès, médias, entourage, institutions). & \esp{evitar la revictimización; trato digno a la víctima} \\
\hline
\esp{Coeducación} & Éducation non sexiste, égalitaire, mixte. & \esp{educación coeducativa, escuela coeducativa, perspectiva de género} \\
\hline
\esp{Empoderamiento} & Processus d'autonomisation, prise de pouvoir/parole. & \esp{empoderamiento de las mujeres, empoderarse, mujer empoderada} \\
\hline
\end{tabularx}
\end{center}

\courssub{Familles de mots et derivation (morphologie lexicale)}

\begin{propriete}[Dérivation nominale / verbale / adjectivale — À maîtriser pour l'expression]
\begin{center}
\begin{tabular}{|l|l|l|l|}
\hline
\rowcolor{popgreenL}
\textbf{Verbe} & \textbf{Nom (action/état)} & \textbf{Adjectif / Participe} & \textbf{Personne / Agent} \\
\hline
\esp{maltratar} & \esp{maltrato} & \esp{maltratado / maltratador} & \esp{maltratador / víctima} \\
\hline
\esp{acosar} & \esp{acoso} & \esp{acosado / acosador} & \esp{acosador / acosada} \\
\hline
\esp{denunciar} & \esp{denuncia} & \esp{denunciado / denunciable} & \esp{denunciante} \\
\hline
\esp{proteger} & \esp{protección} & \esp{protegido / protector} & \esp{protector / protegida} \\
\hline
\esp{discriminar} & \esp{discriminación} & \esp{discriminado / discriminatorio} & \esp{discriminador} \\
\hline
\esp{igualar} & \esp{igualdad} & \esp{igual / igualitario} & \esp{—} \\
\hline
\esp{dominar} & \esp{dominación} & \esp{dominado / dominante} & \esp{dominador} \\
\hline
\esp{silenciar} & \esp{silencio} & \esp{silenciado / silenciador} & \esp{—} \\
\hline
\esp{revictimizar} & \esp{revictimización} & \esp{revictimizado} & \esp{—} \\
\hline
\esp{empoderar} & \esp{empoderamiento} & \esp{empoderado} & \esp{mujer empoderada} \\
\hline
\end{tabular}
\end{center}
\emph{Astuce Bac :} Transformez une phrase nominale en verbale et inversement pour varier la syntaxe (ex. \esp{La denuncia es necesaria} $\leftrightarrow$ \esp{Es necesario denunciar}).
\end{propriete}

\coursec{Méthode du commentaire de texte}

\courssub{Structure type du commentaire composé (Épreuve écrite Bac)}

\begin{methode}[Les 3 étapes obligatoires du commentaire]
\textbf{Durée conseillée : 1h30 -- 2h. Longueur : 250--350 mots.}

\textbf{ÉTAPE 1 : L'INTRODUCTION (environ 50--70 mots)}
\begin{enumerate}
    \item \textbf{Accroche (Contexte) :} Une phrase générale sur le thème (violence, machisme, actualité). \esp{En la sociedad actual, la violencia de género sigue siendo un problema estructural...}
    \item \textbf{Présentation du texte :} Auteur (si connu), titre, source, date, nature. \esp{El artículo « Romper el silencio », publicado en [source/fecha], aborda...}
    \item \textbf{Problématique :} La question à laquelle le texte (et votre commentaire) répond. \esp{¿Hasta qué punto la educación y la denuncia pueden romper el ciclo del machismo?}
    \item \textbf{Annonce du plan :} \esp{En primer lugar analizaremos...; en segundo lugar...; finalmente...}
\end{enumerate}

\textbf{ÉTAPE 2 : LE DÉVELOPPEMENT (2 ou 3 parties, ~200 mots)}
\begin{itemize}
    \item \textbf{Partie A : Le constat alarmant (Paragraphes 1-2 du texte).} Idée directrice : \esp{La persistencia de la violencia y el silencio de las víctimas.}
    \begin{itemize}
        \item Argument + Citation textuelle (\esp{« ... »}) + Explication/Analyse + Exemple personnel / Contexte Cameroun/Guinée.
        \item Connecteurs : \esp{En primer lugar, además, como muestra el texto, por ejemplo, cabe destacar.}
    \end{itemize}
    \item \textbf{Partie B : Les causes structurelles (Paragraphes 1-2).} Idée directrice : \esp{El machismo como sistema de dominación y la impunidad.}
    \begin{itemize}
        \item Argument + Citation + Analyse (éducation patriarcale, stéréotypes, dépendance économique, justice défaillante).
        \item Connecteurs : \esp{En segundo lugar, la raíz del problema radica en..., debido a que, por tanto.}
    \end{itemize}
    \item \textbf{Partie C : Les pistes de solution / Appel à l'action (Paragraphe 3).} Idée directrice : \esp{La urgencia de una respuesta integral: educación, justicia, voluntad política.}
    \begin{itemize}
        \item Argument + Citation (\esp{urge, falta, hay que}) + Évaluation critique (faisabilité, limites).
        \item Connecteurs : \esp{Por último, en cuanto a las soluciones, es fundamental que..., sin embargo, no basta con...}
    \end{itemize}
\end{itemize}

\textbf{ÉTAPE 3 : LA CONCLUSION (environ 40--50 mots)}
\begin{enumerate}
    \item \textbf{Bilan synthétique :} Réponse directe à la problématique. \esp{En definitiva, el texto demuestra que...}
    \item \textbf{Ouverture (Commentaire personnel) :} Élargissement (autre pays, autre texte, actualité récente, rôle des hommes, éducation précoce, rôle des réseaux sociaux). \esp{Esta lucha nos concierne a todos: en Camerún, como en España, la educación es la llave...}
\end{enumerate}
\end{methode}

\courssub{Boîte à outils : Connecteurs logiques pour commenter}

\begin{center}
\begin{tabularx}{\linewidth}{|X|X|}
\hline
\rowcolor{poppurpleL}
\textbf{Fonction} & \textbf{Conectores en español} \\
\hline
Ajouter / Enchaîner & \esp{Además, asimismo, también, de igual modo, por otra parte, en primer lugar / segundo lugar / lugar.} \\
\hline
Opposer / Nuancer & \esp{Sin embargo, no obstante, por el contrario, en cambio, aunque, a pesar de (que + subj.).} \\
\hline
Illustrer / Citer & \esp{Por ejemplo, como ilustra el texto, « ... », véase el caso de..., cabe citar.} \\
\hline
Expliquer / Causalité & \esp{Porque, ya que, puesto que, debido a que, la razón es que, se debe a que.} \\
\hline
Conséquence & \esp{Por tanto, por consiguiente, así que, de modo que, en consecuencia.} \\
\hline
Conclure / Résumer & \esp{En conclusión, en definitiva, en resumen, para terminar, finalmente.} \\
\hline
Exprimer l'opinion / Évaluer & \esp{En mi opinión, me parece que, es evidente que, resulta fundamental que (+ subj.), es urgente que (+ subj.).} \\
\hline
\end{tabularx}
\end{center}

\begin{attention}[Pièges du commentaire : Ce que l'examinateur sanctionne]
\begin{itemize}
    \item \textbf{Le résumé (paraphrase) :} Raconter le texte paragraphe par paragraphe sans analyse. \emph{Interdit.}
    \item \textbf{L'absence de citations :} Un commentaire sans \esp{« ... »} est une dissertation hors sujet. \emph{Zéro point en compréhension.}
    \item \textbf{L'absence de « commentaire personnel » :} La conclusion ne doit pas être une simple répétition. Il faut \textbf{prendre position}, \textbf{confronter} le texte à votre réalité (Cameroun, Guinée, actualité), \textbf{critiquer} ou \textbf{nuancer}.
    \item \textbf{La langue française :} Tout le commentaire s'écrit \textbf{en espagnol}. Mots de liaison, verbes, accords en espagnol.
    \item \textbf{Le hors-sujet :} Parler de la condition de la femme en général sans s'ancrer dans \emph{ce} texte précis.
\end{itemize}
\end{attention}

\coursec{Commentaire modèle et production guidée}

\courssub{Commentaire modèle (sujet : « Romper el silencio »)}

\begin{textebox}[Commentaire type Bac — Niveau Excellence (environ 280 mots)]
\textbf{Introducción} \\
En la sociedad hispánica contemporánea, la violencia de género constituye una lacra persistente que trasciende fronteras y clases sociales. El artículo « Romper el silencio : la lucha contra el machismo », de factura periodística y tono denunciatorio, expone la dramática realidad de las víctimas y exige una respuesta estructural. La pregunta que guía nuestro análisis es: \emph{¿Qué mecanismos permiten sostener el machismo y cómo puede la sociedad desmantelarlos?} Analizaremos primero la pervivencia de la violencia y el silencio cómplice; luego, las causas patriarcales e institucionales; finalmente, las soluciones integrales que propone el autor.

\textbf{Desarrollo} \\
\textbf{1. La violencia invisible y el muro del silencio.} El texto abre con un dato escalofriante: « una de cada tres mujeres » ha sufrido violencia. Esta cifra, avalada por la macroencuesta de 2023, revela la magnitud del fenómeno. Sin embargo, el autor subraya que « el silencio sigue siendo el mayor cómplice ». El 70\% de las víctimas no denuncia por « miedo », « vergüenza » y « dependencia económica ». Este silencio no es pasivo; es el resultado de un terror racional y de una justicia que « a menudo las revictimiza ». En Camerún, como en España, el peso de la tradición y la falta de refugios seguros refuerzan esta omertà.

\textbf{2. El machismo: un sistema, no un acto.} El autor define el machismo como « un sistema de dominación » que considera a la mujer « inferior ». No es un hecho aislado, sino una estructura mental transmitida por « una educación patriarcal » y « estereotipos de género ». El acoso psicológico —control del móvil, aislamiento— « precede casi siempre a la agresión física ». Esta escalada recuerda el « ciclo de la violencia » teorizado por Lenore Walker: tensión, explosión, luna de miel.

\textbf{3. Urgencia de una respuesta integral: educar, proteger, juzgar.} La solución « no es solo penal ». El texto exige tres pilares: « una educación coeducativa desde la escuela » para « desmontar los estereotipos »; « juzgados especializados » con recursos; y « voluntad política » para aplicar la Ley Orgánica 1/2004. La coeducación es la vacuna preventiva; la justicia especializada, el remedio curativo.

\textbf{Conclusión} \\
En definitiva, el texto demuestra que romper el silencio exige un pacto social: denunciar, acompañar, educar. La igualdad « no es un favor, es un derecho humano innegociable ». En Yaoundé como en Madrid, la lucha contra el machismo es la gran tarea democrática de nuestro tiempo: requiere leyes, jueces, pero sobre todo una escuela que enseñe a los niños que no son « fuertes » y a las niñas que no son « cuidadoras », sino seres humanos libres.
\legende{Modèle de commentaire structuré : Introduction (accroche, présentation, problématique, plan), Développement en 3 parties argumentées (texte + connaissance), Conclusion (bilan + ouverture personnelle contextualisée).}
\end{textebox}

\courssub{Production guidée : À vous de jouer !}

\begin{exercice}[Sujet de commentaire d'entraînement (Devoir maison / Bac blanc)]
\textbf{Consigne :} Rédigez un commentaire composé du texte « Romper el silencio » en suivant \textbf{strictement la méthode en 3 étapes} (Introduction, Développement 2-3 parties, Conclusion avec ouverture).
\begin{itemize}
    \item \textbf{Longueur :} 250--300 mots.
    \item \textbf{Langue :} Espagnol exclusivement.
    \item \textbf{Contraintes :} Minimum 4 citations textuelles (\esp{« ... »}). Utiliser au moins 6 connecteurs logiques du tableau ci-dessus. Intégrer une référence au contexte camerounais ou guinéen dans l'ouverture.
    \item \textbf{Problématique suggérée :} \esp{¿En qué medida la educación y la justicia son las claves para erradicar la violencia machista?}
\end{itemize}
\end{exercice}

\begin{methode}[Auto-évaluation avant remise (Grille Bac)]
Cochez chaque case avant de rendre votre copie :
\begin{itemize}
    \item [ ] \textbf{Introduction :} Accroche + Présentation texte + Problématique + Plan annoncé.
    \item [ ] \textbf{Développement :} 2 ou 3 parties distinctes (idées directrices claires). Chaque partie = Argument + \textbf{Citation «...»} + Analyse + Exemple/Contexte.
    \item [ ] \textbf{Connecteurs :} Au moins 6 utilisés correctement (orthographe, accent, place).
    \item [ ] \textbf{Conclusion :} Bilan (réponse à la problématique) + Ouverture personnelle/contextualisée (Cameroun/Guinée/Autre).
    \item [ ] \textbf{Langue :} 0 français. Accents (\esp{á, é, í, ó, ú, ñ, ü}). Ponctuation (\esp{¿ ? ¡ !}). \esp{Ser/Estar} correct. \esp{Por/Para} correct. Subjonctif après \esp{urge que, es necesario que, quiero que, es fundamental que}.
    \item [ ] \textbf{Propreté :} Paragraphes sautés, lisible, pas de ratures excessives.
\end{itemize}
\end{methode}



\begin{corrige}[Exercice n° Compréhension (rappel du bloc précédent)]
\begin{enumerate}
    \item \textbf{El texto define el machismo como « un sistema de dominación » que « considera a la mujer inferior al hombre » y justifica el maltrato como herramienta de control.} (Ce n'est pas un acte isolé, c'est un système).
    \item \textbf{El 70\% de las víctimas no denuncia por: 1) « miedo a las represalias », 2) « vergüenza », 3) « dependencia económica », 4) « no confían en una justicia que a menudo las revictimiza ».}
    \item \textbf{Hace falta « una educación coeducativa desde la escuela » para « desmontar los estereotipos » (niños fuertes / niñas cuidadoras) y prevenir la violencia desde la raíz.}
    \item \textbf{Falso.} El texto dice lo contrario: \esp{« El acoso psicológico... precede casi siempre a la agresión física »}.
    \item \textbf{Synonyme de \esp{imprescindible} :} \esp{« innegociable »} (dernier mot). \textbf{Antonyme de \esp{negociable} :} \esp{« innegociable »} (même mot, préfixe \esp{in-}).
    \item \textbf{Propuesta :} \esp{Romper el silencio es difícil porque las víctimas sienten miedo y vergüenza, tienen dependencia económica y no confían en la justicia que las revictimiza.}
\end{enumerate}
\end{corrige}

\begin{exercice}[Entraînement personnel : Questions sur le texte « Romper el silencio »]
Répondez en espagnol, phrase complète, avec citation (\esp{« ... »}).
\begin{enumerate}
    \item ¿Cómo define el texto el machismo? (Cita y explica con tus palabras).
    \item ¿Qué porcentaje de víctimas no denuncia y cuáles son las cuatro razones textuales?
    \item Según el autor, ¿qué tipo de educación hace falta y para qué?
    \item \textbf{V/F :} « El acoso físico es anterior al acoso psicológico. » (Justifiez).
    \item Trouvez dans le dernier paragraphe : un synonyme de \esp{imprescindible} et un antonyme de \esp{negociable}.
    \item \textbf{Expression écrite courte (30 mots) :} \esp{En tu opinión, ¿por qué es tan difícil « romper el silencio »? Usa: miedo, vergüenza, dependencia, justicia.}
\end{enumerate}
\end{exercice}

\coursec{Lexique thématique : la famille et la condition de la femme}

Après avoir travaillé la compréhension globale et détaillée du texte support, nous allons maintenant fixer et enrichir le vocabulaire indispensable pour parler de la famille, des violences de genre et de la lutte pour l'égalité. Ce lexique vous servira non seulement pour le commentaire de texte, mais aussi pour l'expression écrite et orale tout au long du module.

\courssub{Champ lexical 1 : La famille hispanique (\emph{la familia hispánica})}

Observons d'abord comment le texte support et la vie quotidienne nomment les membres et la structure de la famille. En espagnol, le masculin pluriel englobe souvent les deux sexes (\emph{los padres} = les parents, père et mère), mais l'usage tend aujourd'hui à expliciter le féminin (\emph{las madres y los padres}).

\begin{definition}[Vocabulaire de base : la familia]
\textbf{La familia} : la famille (au sens large). \\
\textbf{El hogar} : le foyer, le domicile (chargé d'affectivité). \\
\textbf{El esposo / La esposa} : l'époux / l'épouse (syn. \emph{el marido} / \emph{la mujer}, mais plus formel). \\
\textbf{Los padres} : les parents (père + mère). \emph{El padre} (père), \emph{la madre} (mère). \\
\textbf{Los hijos} : les enfants (fils + filles). \emph{El hijo} (fils), \emph{la hija} (fille). \\
\textbf{El matrimonio} : le mariage (institution ou cérémonie). \\
\textbf{La pareja} : le couple (mariés, concubins, partenaires). \\
\textbf{La educación} : l'éducation (au sens large : instruction + élevage). \\
\textbf{Los abuelos} : les grands-parents. \emph{El abuelo}, \emph{la abuela}. \\
\textbf{Los hermanos} : les frères et sœurs. \emph{El hermano}, \emph{la hermana}. \\
\textbf{El tío / La tía} : l'oncle / la tante. \\
\textbf{El primo / La prima} : le cousin / la cousine.
\end{definition}

\begin{exemplebox}[Exemples en contexte]
\esp{En mi hogar, mis padres nos enseñan el respeto mutuo.} « Dans mon foyer, nos parents nous enseignent le respect mutuel. » \\
\esp{La pareja decidió casarse tras diez años de convivencia.} « Le couple a décidé de se marier après dix ans de vie commune. » \\
\esp{La educación de los hijos es responsabilidad de ambos progenitores.} « L'éducation des enfants est la responsabilité des deux parents. »
\end{exemplebox}

\begin{attention}[Faux amis et pièges fréquents]
\begin{itemize}
\item \textbf{Parientes} $\neq$ « parents » (père/mère). \emph{Parientes} = « parents » au sens de \textbf{proches, membres de la famille élargie} (oncles, cousins…). Pour « père et mère », on dit \emph{padres} ou \emph{progenitores}.
\item \textbf{Consorte} existe en espagnol (\emph{el consorte}) mais est très formel (souvent royal) ; préférez \emph{esposo/a} ou \emph{pareja}.
\item \textbf{Embarazada} = « enceinte » (et non « embarrassée» !). « Embarrassée » = \emph{avergonzada} / \emph{incómoda}.
\end{itemize}
\end{attention}

\courssub{Champ lexical 2 : Les violences et discriminations (\emph{las violencias y discriminaciones})}

Ce vocabulaire est au cœur du texte étudié. Il faut le maîtriser pour comprendre les articles de presse, les rapports officiels et les témoignages.

\begin{definition}[Vocabulaire des violences de genre]
\textbf{La violencia de género} : violence exercée contre une personne en raison de son sexe (généralement femmes). \\
\textbf{El machismo} : attitude de supériorité masculine, sexisme structurel. \\
\textbf{El maltrato} : maltraitance, mauvais traitements. \emph{Maltrato físico} / \emph{maltrato psicológico}. \\
\textbf{El acoso} : harcèlement (sexuel, moral, scolaire : \emph{acoso escolar}, cyberharcèlement : \emph{ciberacoso}). \\
\textbf{La discriminación} : discrimination. \emph{Discriminación por razón de sexo / género}. \\
\textbf{La víctima} : la victime (masc. \emph{el víctima} rare, on dit \emph{la víctima} pour les deux sexes). \\
\textbf{El agresor / La agresora} : l'agresseur / l'agresseuse. \\
\textbf{El abuso} : abus (souvent sexuel : \emph{abuso sexual}). \\
\textbf{El feminicidio / El femicidio} : meurtre d'une femme en raison de son sexe (terme juridique dans plusieurs pays). \\
\textbf{La trata de personas} : traite des êtres humains (souvent à des fins d'exploitation sexuelle).
\end{definition}

\begin{propriete}[Formation des noms abstraits en \texorpdfstring{\textbf{-dad}}{-dad} (équivalent fr. \textbf{-té})]
Beaucoup de noms féminins abstraits se forment en ajoutant \textbf{-dad} à l'adjectif (souvent avec changement vocalique) :
\begin{center}
\begin{tabularx}{\linewidth}{|X|X|X|}
\hline
\rowcolor{popblueL}
\textbf{Adjectif} & \textbf{Nom en -dad} & \textbf{Français} \\
\hline
\esp{igual} & \esp{igualdad} & égalité \\
\esp{libre} & \esp{libertad} & liberté \\
\esp{real} & \esp{realidad} & réalité \\
\esp{verdadero} $\to$ \esp{verdad} & \esp{verdad} & vérité \\
\esp{responsable} & \esp{responsabilidad} & responsabilité \\
\esp{capaz} & \esp{capacidad} & capacité \\
\esp{igualitario} & \esp{igualitarismo} (isme) & égalitarisme \\
\hline
\end{tabularx}
\end{center}
\emph{Règle} : adjectif en \emph{-o/-a} $\to$ base + \emph{-dad} (\emph{igual} $\to$ \emph{igualdad}) ; adjectif en \emph{-e} ou consonne $\to$ base + \emph{-dad} (\emph{libre} $\to$ \emph{libertad}). Attention à \emph{verdad} (nom primitif, de \emph{verdadero}) et \emph{comunidad} (de \emph{común}).
\end{propriete}

\begin{exemplebox}[Violences : exemples traduits]
\esp{María sufrió maltrato psicológico durante años.} « Marie a subi de la maltraitance psychologique pendant des années. » \\
\esp{El acoso laboral afecta a muchas trabajadoras.} « Le harcèlement au travail affecte de nombreuses travailleuses. » \\
\esp{La discriminación salarial persiste en muchos sectores.} « La discrimination salariale persiste dans de nombreux secteurs. » \\
\esp{El agresor fue detenido gracias a la denuncia de la víctima.} « L'agresseur a été arrêté grâce à la plainte de la victime. »
\end{exemplebox}

\courssub{Champ lexical 3 : La lutte et les droits (\emph{la lucha y los derechos})}

Pour commenter un texte ou rédiger un essai, il faut savoir nommer les actions de résistance, les instruments juridiques et les valeurs défendues.

\begin{definition}[Vocabulaire de la lutte et des droits]
\textbf{La denuncia} : la plainte, la dénonciation (acte juridique). \\
\textbf{Denunciar} : porter plainte, dénoncer (devant la justice/police). \\
\textbf{La igualdad} : l'égalité. \emph{Igualdad de género}, \emph{igualdad de oportunidades}. \\
\textbf{Los derechos de la mujer} : les droits des femmes. \\
\textbf{La ley} : la loi. \emph{Ley orgánica}, \emph{ley contra la violencia de género}. \\
\textbf{Proteger} : protéger. \emph{La ley protege a las víctimas.} \\
\textbf{Luchar contra} : lutter contre. \emph{Luchamos contra el machismo.} \\
\textbf{Defender} : défendre. \emph{Defendemos los derechos humanos.} \\
\textbf{La libertad} : la liberté. \\
\textbf{El respeto} : le respect. \\
\textbf{La sensibilización} : la sensibilisation. \\
\textbf{El empoderamiento} : l'autonomisation (empowerment). \\
\textbf{La orden de protección} : l'ordonnance de protection. \\
\textbf{El refugio} : le refuge, la maison d'accueil.
\end{definition}

\begin{exemplebox}[Lutte et droits : exemples traduits]
\esp{La víctima presentó una denuncia en la comisaría.} « La victime a déposé une plainte au commissariat. » \\
\esp{Las leyes de igualdad protegen a las mujeres frente a la discriminación.} « Les lois sur l'égalité protègent les femmes contre la discrimination. » \\
\esp{Las organizaciones feministas luchan contra la impunidad de los agresores.} « Les organisations féministes luttent contre l'impunité des agresseurs. » \\
\esp{Es fundamental romper el silencio y atreverse a hablar.} « Il est fondamental de briser le silence et d'oser parler. »
\end{exemplebox}

\courssub{Expressions utiles et familles de mots (derivation)}

Maîtriser ces collocations (associations de mots figées) et ces familles vous permettra de varier le vocabulaire et d'éviter les calques du français.

\begin{aretenir}[Expressions clés à apprendre par cœur]
\begin{itemize}
\item \textbf{Sufrir maltrato} (et non *avoir du maltrato) : subir des maltraitances.
\item \textbf{Presentar una denuncia} : déposer une plainte / porter plainte.
\item \textbf{Tener miedo de} + infinitif : avoir peur de…
\item \textbf{Atreverse a} + infinitif : oser…
\item \textbf{Romper el silencio} : briser le silence.
\item \textbf{Ser víctima de} : être victime de…
\item \textbf{Concienciar a la sociedad} : sensibiliser la société.
\item \textbf{Erradicar la violencia} : éradiquer la violence.
\end{itemize}
\end{aretenir}

\begin{propriete}[Familles de mots (derivation nom/verbe/adj.)]
\begin{center}
\begin{tabularx}{\linewidth}{|X|X|X|X|}
\hline
\rowcolor{poporangeL}
\textbf{Verbe} & \textbf{Nom (action/résultat)} & \textbf{Agent / Adjectif} & \textbf{Français (famille)} \\
\hline
\esp{igualar} & \esp{igualdad} & \esp{igualitario/a} & égaliser / égalité / égalitaire \\
\esp{denunciar} & \esp{la denuncia} & \esp{el/la denunciante} & dénoncer / plainte / plaignant \\
\esp{maltratar} & \esp{el maltrato} & \esp{el/la maltratador/a} & maltraiter / maltraitance / maltraitant \\
\esp{acosar} & \esp{el acoso} & \esp{el/la acosador/a} & harceler / harcèlement / harceleur \\
\esp{discriminar} & \esp{la discriminación} & \esp{discriminatorio/a} & discriminer / discrimination / discriminatoire \\
\esp{proteger} & \esp{la protección} & \esp{protector/a} & protéger / protection / protecteur \\
\esp{respetar} & \esp{el respeto} & \esp{respetuoso/a} & respecter / respect / respectueux \\
\hline
\end{tabularx}
\end{center}
\emph{Observation} : le suffixe \emph{-dor/-dora} forme l'agent (celui qui fait l'action) ; \emph{-ción/-sión} forme le nom d'action ; \emph{-orio/-oria} forme l'adjectif relationnel.
\end{propriete}

\courssub{Tableau récapitulatif visuel (carte mentale lexicale)}

\begin{popfigure}
\begin{tikzpicture}[
  node distance=1.2cm and 2cm,
  every node/.style={font=\small},
  column/.style={matrix of nodes, nodes={anchor=west, text width=4.5cm, minimum height=0.7cm}, column sep=0.5cm, row sep=0.2cm},
  header/.style={fill=popblue, text=white, font=\bfseries\small, align=center, minimum height=0.8cm},
  cellfam/.style={fill=popblueL, draw=popline},
  cellvio/.style={fill=poporangeL, draw=popline},
  celllut/.style={fill=popgreenL, draw=popline},
]
% En-têtes
\node[header, minimum width=5cm] (h1) at (0,0) {FAMILIA};
\node[header, minimum width=5cm] (h2) at (5.5,0) {VIOLENCIAS};
\node[header, minimum width=5cm] (h3) at (11,0) {LUCHA Y DERECHOS};

% Colonnes
\matrix[column, anchor=north west] at (0,-0.8) (fam) {
|[cellfam]| \esp{la familia} -- la famille \\
|[cellfam]| \esp{el hogar} -- le foyer \\
|[cellfam]| \esp{los padres} -- les parents \\
|[cellfam]| \esp{la pareja} -- le couple \\
|[cellfam]| \esp{la educación} -- l'éducation \\
|[cellfam]| \esp{los hijos} -- les enfants \\
};
\matrix[column, anchor=north west] at (5.5,-0.8) (vio) {
|[cellvio]| \esp{la violencia de género} -- violence de genre \\
|[cellvio]| \esp{el machismo} -- machisme/sexisme \\
|[cellvio]| \esp{el maltrato} -- maltraitance \\
|[cellvio]| \esp{el acoso} -- harcèlement \\
|[cellvio]| \esp{la discriminación} -- discrimination \\
|[cellvio]| \esp{la víctima / el agresor} -- victime / agresseur \\
};
\matrix[column, anchor=north west] at (11,-0.8) (lut) {
|[celllut]| \esp{la denuncia} -- la plainte \\
|[celllut]| \esp{la igualdad} -- l'égalité \\
|[celllut]| \esp{los derechos de la mujer} -- droits des femmes \\
|[celllut]| \esp{la ley / proteger} -- la loi / protéger \\
|[celllut]| \esp{luchar contra / defender} -- lutter contre / défendre \\
|[celllut]| \esp{la libertad / el respeto} -- liberté / respect \\
};
% Flèches décoratives
\draw[popblue, thick, ->] (h1.south) -- ++(0,-0.2);
\draw[poporange, thick, ->] (h2.south) -- ++(0,-0.2);
\draw[popgreen, thick, ->] (h3.south) -- ++(0,-0.2);
\end{tikzpicture}
\legende{Carte mentale lexicale : trois piliers du thème (Familia / Violencias / Lucha y derechos).}
\end{popfigure}

\courssub{Texte original d'entraînement : « Romper el ciclo »}

Lisez le texte suivant (environ 250 mots), repérez le lexique étudié, puis répondez aux questions.

\begin{textebox}[Texte original — Niveau Terminale A]
En la comunidad de El Progreso, a las afueras de Tegucigalpa, Doña Rosa dirige un pequeño refugio para mujeres que huyen del maltrato. Cada noche, las puertas se abren para recibir a madres con hijos pequeños que no tienen adónde ir. «Muchas llegan con el miedo grabado en la piel», explica Rosa mientras prepara café. «El machismo no es solo golpes; es el control del dinero, la prohibición de trabajar, el aislamiento de la familia. Es una jaula invisible».

Gracias a la nueva ley integral contra la violencia de género, las víctimas pueden solicitar una orden de protección inmediata. Sin embargo, presentar una denuncia sigue siendo un acto de valentía enorme. «Las mujeres temen la represalia del agresor, la vergüenza social y la pérdida de la custodia de sus hijos», añade Rosa. La impunidad, advierte, alimenta el ciclo de la violencia.

En el refugio, talleres de empoderamiento económico y charlas sobre derechos reproductivos ayudan a las usuarias a reconstruir su autoestima. «Aquí aprendemos a decir \emph{basta}. Romper el silencio es el primer paso hacia la libertad», dice Ana, una joven de veintidós años que logró escapar de una relación tóxica tras años de acoso psicológico.

Las organizaciones feministas exigen más recursos para los juzgados especializados y una educación afectivo-sexual obligatoria en las escuelas. «Solo la prevención y la igualdad real podrán erradicar esta lacra», concluye Rosa, con la mirada fija en el horizonte.
\end{textebox}

\begin{definition}[Glossaire des mots difficiles du texte]
\textbf{El Progreso / Tegucigalpa} : ville / capitale du Honduras (contexte centro-américain). \\
\textbf{Doña} : titre de respect pour une femme âgée (équivalent « Madame » mais plus familier/respectueux). \\
\textbf{Grabado en la piel} : gravé dans la peau (marqué profondément). \\
\textbf{Una jaula invisible} : une cage invisible (métaphore du contrôle coercitif). \\
\textbf{Ley integral} : loi globale / complète (couvrant prévention, protection, poursuite). \\
\textbf{Orden de protección} : ordonnance de protection (mesure judiciaire d'éloignement). \\
\textbf{Represalia} : représaille. \\
\textbf{Vergüenza social} : honte sociale (stigmatisation). \\
\textbf{Custodia} : garde (juridique : garde d'enfants). \\
\textbf{Impunidad} : impunité. \\
\textbf{Alimenta el ciclo} : alimente le cycle. \\
\textbf{Talleres de empoderamiento} : ateliers d'autonomisation (empowerment). \\
\textbf{Autoestima} : estime de soi. \\
\textbf{Relación tóxica} : relation toxique. \\
\textbf{Juzgados especializados} : tribunaux spécialisés (violences de genre). \\
\textbf{Educación afectivo-sexual} : éducation affective et sexuelle. \\
\textbf{Erradicar esta lacra} : éradiquer ce fléau (\emph{lacra} = plaie, fléau social).
\end{definition}

\begin{exercice}[Questions de compréhension et réemploi lexical]
\begin{enumerate}
\item Relevez dans le texte \textbf{5 noms} du champ « violences » et \textbf{5 noms} du champ « lutte/droits ».
\item Trouvez l'équivalent espagnol de : « briser le silence », « déposer une plainte », « subir des maltraitances », « ordonnance de protection », « estime de soi ».
\item Expliquez en espagnol (2-3 lignes) la métaphore « \emph{una jaula invisible} ».
\item Conjuguez au présent de l'indicatif : \emph{denunciar} (yo), \emph{proteger} (nosotros), \emph{luchar} (ellos), \emph{sufrir} (ella).
\item Transformez les adjectifs en noms en \emph{-dad} : \emph{responsable, capaz, libre, igual, real}.
\end{enumerate}
\end{exercice}

\courssub{Réemploi immédiat : phrases à compléter et mini-traductions}

\begin{exoresolu}[Exercice résolu : Complétion et traduction]
\textbf{Énoncé :} Complétez les phrases avec le mot approprié (choisissez dans la liste) puis traduisez en français.
\emph{Liste : denuncia, maltrato, igualdad, acoso, víctima, agresor, ley, refugio, empoderamiento, silencio.}

1. La \_\_\_\_\_\_ presentó una \_\_\_\_\_\_ contra su \_\_\_\_\_\_.
2. El \_\_\_\_\_\_ psicológico deja huellas invisibles.
3. La \_\_\_\_\_\_ de género garantiza los mismos derechos para hombres y mujeres.
4. El \_\_\_\_\_\_ ofrece seguridad y apoyo a las mujeres maltratadas.
5. El \_\_\_\_\_\_ económico permite a las mujeres ser independientes.
6. Es hora de romper el \_\_\_\_\_\_ y hablar.

\tcblower
\textbf{Solution détaillée :}
1. \textbf{La \esp{víctima} presentó una \esp{denuncia} contra su \esp{agresor}.} \\
   « La victime a déposé une plainte contre son agresseur. » \\
   \emph{Note} : \emph{presentar una denuncia} = locution figée (déposer une plainte). \emph{Víctima} est invariant en genre (la víctima / el víctima rare).

2. \textbf{El \esp{maltrato} psicológico deja huellas invisibles.} \\
   « La maltraitance psychologique laisse des traces invisibles. » \\
   \emph{Maltrato} = nom masculin. \emph{Huella} = trace, empreinte.

3. \textbf{La \esp{igualdad} de género garantiza los mismos derechos para hombres y mujeres.} \\
   « L'égalité de genre garantit les mêmes droits pour les hommes et les femmes. » \\
   \emph{Garantizar} = garantir. \emph{Mismos} = mêmes (adjectif placé avant le nom pour insister sur l'identité).

4. \textbf{El \esp{refugio} ofrece seguridad y apoyo a las mujeres maltratadas.} \\
   « Le refuge offre sécurité et soutien aux femmes maltraitées. » \\
   \emph{Maltratadas} = participe passé adjectivé (accord avec \emph{mujeres}).

5. \textbf{El \esp{empoderamiento} económico permite a las mujeres ser independientes.} \\
   « L'autonomisation économique permet aux femmes d'être indépendantes. » \\
   \emph{Permitir a + infinitif} = permettre à [qqn] de [faire].

6. \textbf{Es hora de romper el \esp{silencio} y hablar.} \\
   « Il est temps de briser le silence et de parler. » \\
   \emph{Es hora de + infinitif} = il est temps de… \emph{Romper} = briser (physique ou abstrait).
\end{exoresolu}

\begin{exercice}[Mini-traductions : Français $\to$ Espagnol]
Traduisez les phrases suivantes en utilisant le vocabulaire de la leçon.
\begin{enumerate}
\item Ma sœur a subi du harcèlement au travail.
\item Nous devons lutter contre le machisme dans la société.
\item La loi protège les victimes de violence de genre.
\item Elles ont osé briser le silence et porter plainte.
\item L'égalité réelle est un droit fondamental.
\item Le refuge accueille des mères et leurs enfants.
\end{enumerate}
\end{exercice}

\begin{corrige}[Exercice Mini-traductions]
\begin{enumerate}
\item \esp{Mi hermana sufrió acoso laboral.} (\emph{Sufrir acoso} / \emph{acoso laboral}).
\item \esp{Debemos luchar contra el machismo en la sociedad.} (\emph{Debemos} = nous devons / \emph{luchar contra}).
\item \esp{La ley protege a las víctimas de violencia de género.} (\emph{Proteger a} + personne = \emph{a} personnel).
\item \esp{Se atrevieron a romper el silencio y presentar una denuncia.} (\emph{Atreverse a} = oser / \emph{presentar una denuncia} = porter plainte).
\item \esp{La igualdad real es un derecho fundamental.}
\item \esp{El refugio acoge a madres y a sus hijos.} (\emph{Acoger} = accueillir / \emph{a} personnel devant \emph{madres}).
\end{enumerate}
\end{corrige}

\courssub{Point culture : Le saviez-vous ?}

\begin{savaistu}[Machismo, feminicidio et contexte africain]
\begin{itemize}
\item \textbf{Machismo} : vient de \emph{macho} (mâle). Désigne l'idéologie de la supériorité masculine. Son opposé revendiqué est \textbf{feminismo} (féminisme).
\item \textbf{Feminicidio / Femicidio} : terme juridique adopté au Mexique (2007), au Guatemala, en Argentine, en Espagne… pour qualifier le meurtre d'une femme \emph{parce qu'elle est femme}. Au Cameroun, on parle aussi de \emph{féminicide} dans les rapports ONU / société civile.
\item \textbf{Guinée équatoriale} : seul pays africain hispanophone. La \emph{Ley de Igualdad de Género} (2011) existe mais son application reste perfectible ; les violences conjugales y sont encore souvent traitées comme affaires privées.
\item \textbf{Cameroun} : la loi n°2016/007 du 12 juillet 2016 portant code pénal réprime les violences conjugales (art. 356-357), mais le silence culturel et la dépendance économique freinent les dépôts de plainte (\emph{denuncias}). Des ONG comme \emph{ALVF} (Association de lutte contre les violences faites aux femmes) mènent un travail de \emph{sensibilización} comparable aux \emph{talleres} du texte.
\end{itemize}
\end{savaistu}

\begin{attention}[Synthèse des pièges à éviter pour le commentaire]
\begin{itemize}
\item Ne confondez pas \textbf{denunciar} (porter plainte, acte juridique) et \textbf{criticar} (critiquer). \emph{Denunciar} implique une autorité (police, juge).
\item \textbf{Violencia de género} $\neq$ \textbf{violencia doméstica} (violence domestique = dans le foyer, peut toucher hommes, enfants, personnes âgées). La violence de genre vise spécifiquement les femmes \emph{en tant que femmes}.
\item \textbf{Igualdad} (nom fém.) vs \textbf{igualitario/a} (adj.). Écrivez \emph{políticas igualitarias} (politiques égalitaires), pas *\emph{políticas igualdades}.
\item \textbf{Accents} : \emph{denuncia} (nom fém., llana terminée par voyelle $\to$ \textbf{pas d'accent écrit}) ; \emph{denuncie} (subj. présent, llana $\to$ pas d'accent) ; \emph{denuncié} (participe, oxytone $\to$ accent sur \emph{é}). \emph{Mamá, papá, sofá} portent l'accent sur la dernière syllabe (oxytons terminés par voyelle).
\item \textbf{Ser víctima de} + nom / \textbf{sufrir} + nom (maltrato, acoso). Pas *\emph{tener maltrato}.
\end{itemize}
\end{attention}

\courssub{Méthode : Intégrer le lexique dans un commentaire structuré}

\begin{methode}[Réemployer le vocabulaire thématique dans votre commentaire]
\begin{enumerate}
\item \textbf{Identifiez} les idées directrices du texte : ex. « La violence structurelle (machismo) », « L'arsenal juridique (ley, denuncia, orden de protección) », « L'accompagnement des victimes (refugio, empoderamiento) ».
\item \textbf{Sélectionnez} 8-10 mots-clés précis des trois champs lexicaux pour les citer entre guillemets dans votre développement.
\item \textbf{Articulez} vos arguments avec des connecteurs logiques : \emph{En primer lugar, además, sin embargo, por consiguiente, en conclusión}.
\item \textbf{Illustrez} chaque point par un exemple du texte (citation courte) ou par votre connaissance du contexte (Cameroun, Guinée équatoriale, Espagne, Amérique latine).
\item \textbf{Concluez} en ouvrant sur la nécessité de la \emph{educación afectivo-sexual} et de l'\emph{igualdad real} pour \emph{erradicar la lacra}.
\end{enumerate}
\end{methode}

\begin{exemplebox}[Commentaire modèle guidé (structure type)]
\textbf{Introduction} : Présentation du texte (source, thème, thèse : la violence de genre comme cycle structurel brisé par la loi et l'accompagnement). Annonce du plan en trois axes.

\textbf{Développement} :
\begin{itemize}
\item \textbf{Axe 1 : La violence structurelle (\emph{el machismo}).} Le texte montre que \emph{el machismo} n'est pas seulement physique (\emph{golpes}) mais psychologique et économique (\emph{control del dinero, aislamiento}). Citation : « \emph{Es una jaula invisible} ». Exemple contextuel : au Cameroun, la dépendance économique freine les \emph{denuncias}.
\item \textbf{Axe 2 : La réponse juridique (\emph{la ley, la denuncia, la orden de protección}).} La \emph{ley integral} offre des outils (\emph{orden de protección}), mais la \emph{impunidad} et la \emph{vergüenza social} entravent l'accès à la justice. Citation : « \emph{presentar una denuncia sigue siendo un acto de valentía enorme} ».
\item \textbf{Axe 3 : L'accompagnement et l'autonomisation (\emph{refugio, empoderamiento, autoestima}).} Le \emph{refugio} n'est pas qu'un abri : \emph{talleres de empoderamiento}, \emph{derechos reproductivos}. Citation : « \emph{Romper el silencio es el primer paso hacia la libertad} ».
\end{itemize}

\textbf{Conclusion} : Bilan : la loi est nécessaire mais insuffisante sans \emph{sensibilización} et \emph{educación afectivo-sexual}. Ouverture : l'\emph{igualdad real} conditionne l'\emph{erradicación de la lacra}.
\end{exemplebox}

\begin{experience}[Activité orale guidée : « Tribunal de la parole »]
En binômes : l'un joue le rôle d'une victime qui \emph{rompe el silencio} et \emph{presenta una denuncia} ; l'autre est l'agent de police / l'avocat / le juge qui note la \emph{denuncia}. Utilisez obligatoirement : \emph{sufrir maltrato, tener miedo de, atreverse a, orden de protección, agresor, víctima, ley}. Changez de rôle. Durée : 5 min par binôme. L'enseignant circule et note les emplois corrects du subjonctif après \emph{es necesario que}, \emph{quiero que}… (révision B2).
\end{experience}

\coursec{Méthode du commentaire de texte}

Dans la section précédente, nous avons constitué ensemble le lexique thématique de la famille et de la condition de la femme : \esp{violencia de género}, \esp{machismo}, \esp{maltrato}, \esp{denuncia}, \esp{igualdad}, \esp{acoso}. Ces mots ne sont pas une fin en soi : ce sont des \cle{outils}. Au baccalauréat, on ne vous demandera pas seulement de comprendre un texte, mais de le \emph{commenter}, c'est-à-dire de le présenter, d'en dégager les idées et de donner votre avis argumenté. Cette section vous donne la méthode complète, étape par étape, avec un texte d'entraînement, un début de commentaire modèle et les critères exactes de l'évaluation.

\courssub{Observer d'abord : un texte support et un extrait de commentaire}

Avant toute règle, lisons un texte original sur le thème de la leçon, puis observons comment un bon élève commence son commentaire.

\begin{textebox}[La lucha contra la violencia de género]
En muchos países del mundo, las mujeres siguen sufriendo discriminaciones y violencias dentro y fuera de la familia. La violencia de género no es un problema privado : es un problema de toda la sociedad. Según las asociaciones de defensa de los derechos de la mujer, muchas víctimas no denuncian los malos tratos por miedo, por vergüenza o por dependencia económica.

En primer lugar, el machismo sigue presente en la vida cotidiana : en algunas familias, las niñas dejan la escuela antes que los niños, y las mujeres asumen solas las tareas domésticas. Además, el acoso en el trabajo y en la calle sigue siendo una realidad para muchas jóvenes.

Sin embargo, hay señales de esperanza. Cada vez más mujeres estudian, trabajan y participan en la vida política. En España y en América Latina, las leyes contra la violencia machista son cada vez más estrictas, y las denuncias aumentan porque las víctimas ya no se sienten solas. La educación de los niños y de las niñas en la igualdad es, sin duda, la clave del futuro.

La pregunta esencial es la siguiente : ¿cómo transformar las mentalidades para que la igualdad entre hombres y mujeres sea una realidad y no solamente un discurso ?
\end{textebox}

\textbf{Glossaire.} \esp{los malos tratos} : les mauvais traitements ; \esp{por vergüenza} : par honte ; \esp{las tareas domésticas} : les tâches ménagères ; \esp{cada vez más} : de plus en plus ; \esp{estrictas} : strictes, sévères ; \esp{las mentalidades} : les mentalités ; \esp{un discurso} : un discours.

\textbf{Questions de compréhension.}
\begin{enumerate}
\item ¿Por qué muchas víctimas no denuncian los malos tratos ?
\item ¿Qué señales de esperanza menciona el texto ?
\item ¿Cuál es, según el autor, la clave del futuro ?
\end{enumerate}

\textbf{Réponses attendues.} 1. \esp{No denuncian por miedo, por vergüenza o por dependencia económica.} 2. \esp{Cada vez más mujeres estudian, trabajan y participan en la vida política ; las leyes son más estrictas y las denuncias aumentan.} 3. \esp{La clave del futuro es la educación de los niños y de las niñas en la igualdad.}

Observons maintenant le début d'un commentaire d'élève sur ce texte :

\begin{exemplebox}[Un début de commentaire réussi]
\esp{El texto es un artículo de opinión sobre la violencia de género. El autor denuncia el machismo y las discriminaciones que sufren las mujeres, pero también muestra señales de esperanza. La problemática es la siguiente : ¿cómo transformar las mentalidades para que la igualdad sea una realidad ?}

« Le texte est un article d'opinion sur la violence de genre. L'auteur dénonce le machisme et les discriminations que subissent les femmes, mais il montre aussi des signes d'espérance. La problématique est la suivante : comment transformer les mentalités pour que l'égalité devienne une réalité ? »
\end{exemplebox}

Que remarque-t-on ? En trois phrases, l'élève a donné la \cle{nature} du texte, son \cle{thème} et la \cle{problématique}. C'est exactement la première étape de la méthode.

\courssub{Les quatre étapes du commentaire de texte}

\begin{popfigure}
\begin{tikzpicture}[font=\small]
\node[draw=popblue, fill=popblueL, rounded corners, minimum width=3.4cm, minimum height=1cm, align=center] (e1) at (0,4.5) {\textbf{1. Presentación}\\ naturaleza, fuente, tema,\\ problemática};
\node[draw=poporange, fill=poporangeL, rounded corners, minimum width=3.4cm, minimum height=1cm, align=center] (e2) at (4.2,3) {\textbf{2. Ideas directrices}\\ 2 o 3 ideas, cada una\\ con citas del texto};
\node[draw=popgreen, fill=popgreenL, rounded corners, minimum width=3.4cm, minimum height=1cm, align=center] (e3) at (8.4,1.5) {\textbf{3. Comentario personal}\\ tu opinión argumentada,\\ comparación con Camerún};
\node[draw=poppurple, fill=poppurpleL, rounded corners, minimum width=3.4cm, minimum height=1cm, align=center] (e4) at (12.6,0) {\textbf{4. Conclusión}\\ balance y apertura};
\draw[-{Stealth}, very thick, popblue] (e1.south east) -- (e2.north west);
\draw[-{Stealth}, very thick, poporange] (e2.south east) -- (e3.north west);
\draw[-{Stealth}, very thick, popgreen] (e3.south east) -- (e4.north west);
\end{tikzpicture}
\legende{Les quatre étapes du commentaire de texte, en escalier : chaque étape descend vers la suivante.}
\end{popfigure}

\courssub{Étape 1 — La présentation (2 à 4 phrases)}

\begin{methode}[Rédiger la présentation]
\begin{enumerate}
\item \textbf{Nature du texte} : indique le genre du document. \esp{El texto es un artículo de prensa / un diálogo / un poema / un fragmento de novela.}
\item \textbf{Source} : si elle est connue, cite-la. \esp{El texto está extraído del periódico... / aparece en la revista...} Si elle est inconnue, écris simplement \esp{un texto periodístico}.
\item \textbf{Thème} : annonce le sujet en une phrase. \esp{El tema es la violencia de género dentro de la familia.}
\item \textbf{Problématique} : formule la question centrale du texte, sous forme interrogative. \esp{¿Cómo luchar contra las violencias que sufren las mujeres ?}
\end{enumerate}
\end{methode}

\begin{methode}[Formuler la problématique]
La problématique est toujours une \cle{question}. Pour la trouver, demande-toi : « À quelle question l'auteur essaie-t-il de répondre ? » Utilise les mots interrogatifs espagnols : \esp{¿Cómo...?} (comment), \esp{¿Por qué...?} (pourquoi), \esp{¿En qué medida...?} (dans quelle mesure), \esp{¿Qué soluciones...?} (quelles solutions). N'oublie jamais le point d'interrogation ouvrant \esp{¿} au début de la phrase ni l'accent sur les mots interrogatifs (\esp{cómo}, \esp{por qué}, \esp{qué}).
\end{methode}

\courssub{Étape 2 — Le dégagement des idées directrices}

\begin{methode}[Dégager les idées directrices]
\begin{enumerate}
\item Relis le texte et repère sa structure : en général, \cle{une idée = un paragraphe} ou un groupe de paragraphes.
\item Formule chaque idée en \cle{une seule phrase complète}, avec tes propres mots.
\item Justifie chaque idée par une \cle{citation} du texte, placée entre guillemets (« ... ») et introduite par un verbe comme \esp{el autor dice}, \esp{afirma}, \esp{denuncia}, \esp{subraya}.
\item Annonce tes idées avec les connecteurs : \esp{En primer lugar... En segundo lugar... Por último...}
\end{enumerate}
\end{methode}

\begin{exemplebox}[Deux idées directrices justifiées]
\esp{En primer lugar, el autor denuncia la gravedad del machismo en la sociedad : « el machismo sigue presente en la vida cotidiana ».}

« Premièrement, l'auteur dénonce la gravité du machisme dans la société : "le machisme reste présent dans la vie quotidienne". »

\esp{En segundo lugar, subraya que existen señales de esperanza : « cada vez más mujeres estudian, trabajan y participan en la vida política ».}

« Deuxièmement, il souligne qu'il existe des signes d'espérance : "de plus en plus de femmes étudient, travaillent et participent à la vie politique". »
\end{exemplebox}

\begin{attention}[Citer sans recopier tout le texte]
Une citation doit être \cle{courte} (une ligne maximum) et \cle{exacte} : on ne modifie jamais les mots de l'auteur. Erreur fréquente des francophones : oublier les guillemets ou transformer la citation en résumé. La citation sert de \emph{preuve} ; c'est ta phrase d'introduction qui exprime l'idée.
\end{attention}

\courssub{Étape 3 — Le commentaire personnel}

C'est l'étape la plus importante et la plus valorisée au baccalauréat. Tu dois donner \cle{ton avis}, l'argumenter et, si possible, comparer la situation décrite dans le texte avec celle du Cameroun, de l'Afrique ou d'un autre pays.

\begin{attention}[Ne pas confondre résumé et commentaire]
Le \emph{résumé} répète le texte ; le \emph{commentaire personnel} exige TON avis argumenté. Si ton correcteur peut retrouver toutes tes phrases dans le texte, tu as fait un résumé, pas un commentaire. Signe d'un vrai commentaire : la présence de \esp{en mi opinión}, \esp{pienso que}, \esp{me parece que}, \esp{estoy de acuerdo con}, suivis d'arguments personnels et d'exemples concrets.
\end{attention}

\begin{exemplebox}[Exemples traduits de phrases d'opinion]
\esp{En mi opinión, la violencia de género es un problema grave porque destruye a las familias.}

« À mon avis, la violence de genre est un problème grave parce qu'elle détruit les familles. »

\esp{Estoy de acuerdo con el autor : la educación en la igualdad es la clave del futuro.}

« Je suis d'accord avec l'auteur : l'éducation à l'égalité est la clé du futur. »

\esp{En Camerún, la situación es parecida : muchas mujeres no denuncian por miedo, pero las asociaciones y las redes sociales ayudan cada vez más a las víctimas.}

« Au Cameroun, la situation est semblable : beaucoup de femmes ne portent pas plainte par peur, mais les associations et les réseaux sociaux aident de plus en plus les victimes. »
\end{exemplebox}

\begin{aretenir}[Après les verbes d'opinion négatifs, le subjonctif]
Attention à la grammaire : après \esp{pienso que}, \esp{creo que}, \esp{me parece que} affirmatifs, on utilise l'\cle{indicatif} : \esp{Pienso que la educación es la solución.} « Je pense que l'éducation est la solution. » Mais après la forme négative, on passe au \cle{subjonctif} : \esp{No pienso que la situación sea fácil.} « Je ne pense pas que la situation soit facile. » Remarque la forme \esp{sea}, subjonctif présent du verbe irrégulier \esp{ser}.
\end{aretenir}

\courssub{Étape 4 — La conclusion}

La conclusion fait le \cle{bilan} en une ou deux phrases et s'achève par une \cle{ouverture} : une question plus large, un appel, un espoir.

\begin{exemplebox}[Une conclusion modèle]
\esp{Para concluir, este texto denuncia las violencias contra las mujeres pero abre una puerta a la esperanza : la educación y las leyes pueden cambiar las mentalidades. La igualdad no es un sueño : es un objetivo para toda la sociedad.}

« Pour conclure, ce texte dénonce les violences contre les femmes mais ouvre une porte à l'espérance : l'éducation et les lois peuvent changer les mentalités. L'égalité n'est pas un rêve : c'est un objectif pour toute la société. »
\end{exemplebox}

\courssub{Les connecteurs du commentaire}

Les connecteurs sont la charpente de ton texte. Sans eux, le commentaire ressemble à une liste de phrases sans lien.

\begin{center}
\begin{tabularx}{\linewidth}{|X|X|X|}
\hline
\rowcolor{popblueL} \textbf{Español} & \textbf{Français} & \textbf{Fonction} \\
\hline
\esp{en primer lugar / en segundo lugar} & premièrement / deuxièmement & annoncer les idées \\
\hline
\esp{por un lado / por otro lado} & d'un côté / de l'autre côté & opposer deux aspects \\
\hline
\esp{sin embargo} & cependant & introduire une opposition \\
\hline
\esp{además} & de plus, en outre & ajouter un argument \\
\hline
\esp{en mi opinión / pienso que} & à mon avis / je pense que & donner son avis \\
\hline
\esp{por ejemplo} & par exemple & illustrer un argument \\
\hline
\esp{para concluir / en conclusión} & pour conclure / en conclusion & ouvrir la conclusion \\
\hline
\end{tabularx}
\end{center}

\begin{attention}[Pièges des francophones]
\begin{itemize}
\item \esp{Sin embargo} seul signifie « cependant » ; ne traduis jamais « cependant » par \esp{sin pero} !
\item \esp{Además} prend un accent sur le \emph{a} final ; \esp{ademas} sans accent est une faute.
\item \esp{En mi opinión} : pas d'article, contrairement au français « à \emph{mon} avis ». On ne dit jamais \esp{en la mi opinión}.
\item \esp{Por un lado} et non \esp{por una lado} : \esp{lado} est masculin.
\item \esp{Denunciar} signifie « dénoncer » mais aussi « porter plainte » ; le nom \esp{la denuncia} signifie « la plainte » ou « la dénonciation ». Ne confonds pas avec \esp{anunciar} (annoncer).
\end{itemize}
\end{attention}

\courssub{Exercice résolu et critères d'évaluation au Bac}

\begin{exoresolu}[Rédiger la présentation et une idée directrice]
Énoncé. À partir du texte « La lucha contra la violencia de género » étudié plus haut, rédigez : a) la présentation complète (nature, thème, problématique) ; b) une idée directrice justifiée par une citation.
\tcblower
Solution détaillée.

a) Présentation : \esp{El texto es un artículo de opinión sobre las discriminaciones y las violencias contra las mujeres. El autor denuncia el machismo, pero también muestra los progresos de la igualdad. La problemática es : ¿cómo transformar las mentalidades para que la igualdad sea una realidad ?}

b) Idée directrice : \esp{En primer lugar, el autor explica que muchas víctimas callan : afirma que « muchas víctimas no denuncian los malos tratos por miedo, por vergüenza o por dependencia económica ».}

Remarques méthodologiques : la présentation tient en trois phrases (nature + thème, position de l'auteur, problématique interrogative avec \esp{¿...?}) ; l'idée directrice est formulée en une phrase avec les mots de l'élève (\esp{muchas víctimas callan}), puis prouvée par une citation courte et exacte entre guillemets, introduite par le verbe \esp{afirma}.
\end{exoresolu}

\begin{aretenir}[Critères d'évaluation au baccalauréat]
Le correcteur évalue ton commentaire selon quatre critères :
\begin{enumerate}
\item \textbf{Respect de la structure} : présentation, idées directrices, commentaire personnel, conclusion — les quatre étapes doivent être présentes et dans l'ordre.
\item \textbf{Exactitude de la langue} : conjugaison correcte, accords, accents, ponctuation espagnole (¿ ¡), connecteurs bien employés.
\item \textbf{Pertinence des arguments} : ton avis doit être clair, argumenté et en rapport avec le texte ; la comparaison avec le Cameroun est valorisée.
\item \textbf{Citations du texte} : chaque idée directrice doit s'appuyer sur une citation exacte entre guillemets.
\end{enumerate}
Un commentaire sans citation perd des points ; un commentaire sans avis personnel aussi. Les deux sont indispensables.
\end{aretenir}

\begin{exercice}[À ton tour]
Rédige un commentaire complet (15 à 20 lignes) du texte « La lucha contra la violencia de género ». Respecte les quatre étapes, utilise au moins cinq connecteurs du tableau, deux citations exactes, et compare la situation du texte avec celle du Cameroun dans ton commentaire personnel.
\end{exercice}

Tu disposes maintenant de toute la méthode : dans la section suivante, tu verras un commentaire modèle entièrement rédigé, puis tu produiras le tien de façon guidée.

\coursec{Commentaire modèle et production guidée}

Après avoir acquis la méthode du commentaire de texte (présentation, idées directrices, commentaire personnel, conclusion), nous passons à l'application concrète. Cette section vous guide pas à pas : vous observerez d'abord un commentaire modèle complet, vous en repérerez la structure et les outils linguistiques, puis vous produirez votre propre commentaire en trois temps, avant de l'améliorer par une relecture croisée. Un complément « Bac » vous montre comment enrichir votre avis personnel par une comparaison Cameroun / monde hispanique.

\courssub{Lecture commentée d'un commentaire modèle}

Voici un commentaire modèle rédigé sur le texte support étudié en section 1 (article de presse sur les violences de genre dans le monde hispanique). Observez sa structure en quatre parties, l'emploi des connecteurs logiques et l'intégration de citations courtes.

\begin{textebox}[Commentaire modèle — Extrait complet]
El texto es un artículo periodístico publicado en un diario digital español que trata de la violencia de género en el mundo hispano. El autor denuncia la persistencia del machismo y presenta datos sobre las denuncias y las víctimas mortales.

En primer lugar, la idea directriz principal es que el machismo es la causa estructural del maltrato. El texto afirma que « el machismo considera a la mujer inferior al hombre y justifica el control sobre su vida ». Esta cita muestra que la violencia no es un hecho aislado, sino el resultado de una cultura patriarcal arraigada. Además, el autor explica que el acoso psicológico suele preceder a la violencia física.

En segundo lugar, el texto destaca las respuestas institucionales y sociales : las leyes de igualdad, los juzgados especializados y las asociaciones feministas que ayudan a las víctimas a romper el silencio. La frase « la denuncia es el primer paso hacia la libertad » resume esta idea.

En mi opinión, este problema también existe en Camerún, aunque con matices culturales distintos. Muchas mujeres sufren en silencio por miedo al rechazo familiar o a la pérdida de la custodia de los hijos. Pienso que la educación, desde la escuela primaria, es la mejor arma para cambiar las mentalidades y prevenir el maltrato. Los hombres deben ser aliados en esta lucha.

Para concluir, la lucha por la igualdad es tarea de todos : Estados, familias, escuelas y medios de comunicación. Solo una transformación cultural profunda podrá erradicar la violencia de género.
\end{textebox}

\noindent\textbf{Analyse de la structure :}
\begin{itemize}
  \item \textbf{Présentation} (paragraphe 1) : nature du texte, source, thème, thèse de l'auteur.
  \item \textbf{Idée directrice 1} (paragraphe 2) : cause structurelle (\emph{machismo}) + citation intégrée + explication.
  \item \textbf{Idée directrice 2} (paragraphe 3) : réponses/solutions + citation + développement.
  \item \textbf{Opinion personnelle} (paragraphe 4) : avis argumenté, ouverture sur le Cameroun, proposition concrète.
  \item \textbf{Conclusion} (paragraphe 5) : synthèse + perspective large.
\end{itemize}

\courssub{Repérage des connecteurs et des citations}

Le commentaire modèle utilise des connecteurs pour organiser la pensée et guider le lecteur. Le tableau ci-dessous répertorie les principaux connecteurs employés, leur fonction et leur équivalent français.

\begin{center}
\begin{tabularx}{\linewidth}{|X|X|X|}
\hline
\rowcolor{popblueL}
\textbf{Connecteur (ES)} & \textbf{Fonction} & \textbf{Équivalent FR} \\
\hline
En primer lugar / En segundo lugar & Ordonner les idées & En premier lieu / En second lieu \\
\hline
Además & Ajouter un argument & De plus / En outre \\
\hline
Esta cita muestra que / La frase ... resume & Introduire l'analyse d'une citation & Cette citation montre que / La phrase ... résume \\
\hline
En mi opinión / Pienso que & Marquer l'opinion personnelle & À mon avis / Je pense que \\
\hline
Aunque & Concession (idée adverse) & Bien que / Alors que \\
\hline
Por miedo a / Para & Cause / But & Par peur de / Pour \\
\hline
Para concluir / En definitiva & Conclure & Pour conclure / En définitive \\
\hline
Solo ... podrá & Condition restrictive & Seule ... pourra \\
\hline
\end{tabularx}
\end{center}

\noindent\textbf{Règle d'intégration des citations :} Une citation courte (moins de 3 lignes) s'insère dans la phrase entre guillemets espagnols \esp{« ... »} ou \esp{" ... "}, précédée de deux-points ou intégrée syntaxiquement. Exemple : \esp{El autor afirma que « la denuncia es el primer paso hacia la libertad ».} « L'auteur affirme que "la dénonciation est le premier pas vers la liberté". »

\courssub{Production guidée en trois temps}

Vous allez maintenant rédiger votre propre commentaire sur le même texte support (ou sur un texte analogue fourni par l'enseignant). Suivez la méthode ci-dessous.

\begin{methode}[Production guidée du commentaire — 3 temps]
\textbf{Temps 1 — Au brouillon : construire le plan} (10 min)
\begin{enumerate}
  \item Relisez le texte et soulignez 2 idées forces + 1 citation par idée.
  \item Remplissez le « plan vierge » (figure ci-dessous) : présentation, idée 1 + citation, idée 2 + citation, opinion + conclusion.
  \item Sélectionnez au moins \textbf{6 mots du lexique étudié} (\cle{violencia de género}, \cle{machismo}, \cle{maltrato}, \cle{denuncia}, \cle{igualdad}, \cle{acoso}, \cle{patriarcado}, \cle{víctima}, \cle{ley}, \cle{asociación}) à réemployer.
\end{enumerate}

\textbf{Temps 2 — Rédaction guidée} (25 min)
\begin{enumerate}
  \item \textbf{Présentation} (3-4 lignes) : \esp{El texto es un artículo que trata de ... El autor denuncia ...}
  \item \textbf{Idée 1} (5-6 lignes) : \esp{La primera idea es que ... : « citation ». Esto significa que ... Además, ...}
  \item \textbf{Idée 2} (5-6 lignes) : \esp{La segunda idea destaca ... : « citation ». Esto demuestra que ...}
  \item \textbf{Opinion personnelle} (8-10 lignes) : \esp{En mi opinión, ... porque ... En Camerún, ... Pienso que la mejor solución es ...}
  \item \textbf{Conclusion} (2-3 lignes) : \esp{Para concluir, ... Es tarea de ...}
\end{enumerate}

\textbf{Temps 3 — Relecture et enrichissement} (10 min)
\begin{enumerate}
  \item Vérifiez la présence des 6 mots de lexique minimum.
  \item Remplacez les répétitions par des synonymes (\esp{mujer} → \esp{la mujer}, \esp{la víctima}, \esp{ella} ; \esp{violencia} → \esp{maltrato}, \esp{abuso}, \cle{agresión}).
  \item Contrôlez les accords (genre/nombre), le subjonctif après \esp{es importante que}, \esp{es necesario que}, et les accents.
\end{enumerate}
\end{methode}

\noindent\textbf{Banque de phrases de départ (à adapter) :}
\begin{itemize}
  \item \esp{El texto trata de la violencia de género en ...}
  \item \esp{El autor denuncia el machismo como causa principal del maltrato.}
  \item \esp{La primera idea directriz es que ... : « ... »}
  \item \esp{La segunda idea se refiere a las soluciones : las leyes de igualdad y las asociaciones.}
  \item \esp{En mi opinión, este problema también existe en Camerún porque ...}
  \item \esp{Estoy de acuerdo con el autor cuando dice que ... porque ...}
  \item \esp{No estoy de acuerdo con la idea de que ... ya que ...}
  \item \esp{Pienso que la educación es fundamental para prevenir la violencia.}
  \item \esp{Para concluir, la lucha por la igualdad requiere la participación de todos.}
\end{itemize}

\courssub{Relecture et correction par les pairs — Grille simple}

Échangez votre production avec un camarade. Utilisez cette grille pour corriger mutuellement. Cochez \textbf{Oui} ou \textbf{Non} et notez une observation si nécessaire.

\begin{center}
\begin{tabularx}{\linewidth}{|X|c|c|X|}
\hline
\rowcolor{popblueL}
\textbf{Critère} & \textbf{Oui} & \textbf{Non} & \textbf{Observation / Correction} \\
\hline
Structure : 4 parties visibles (présentation, 2 idées, opinion + conclusion) & & & \\
\hline
Au moins 6 mots du lexique thématique réemployés & & & \\
\hline
2 citations courtes intégrées correctement (guillemets « ») & & & \\
\hline
Connecteurs logiques variés (al menos 5) & & & \\
\hline
Opinion personnelle argumentée (8-10 lignes, pas un résumé) & & & \\
\hline
Grammaire : accords, subjonctif après expression de nécessité/opinion, accents & & & \\
\hline
Orthographe espagnole (ñ, á, é, í, ó, ú, ü, ¿, ¡) & & & \\
\hline
Pas de mots français déguisés (violencia ≠ violence, denuncia ≠ dénonciation) & & & \\
\hline
\end{tabularx}
\end{center}

\noindent Après correction, réécrivez votre version finale sur votre cahier.

\courssub{Complément Bac : comparaison Cameroun / monde hispanique}

\begin{savaistu}[Le commentaire personnel au Baccalauréat A]
Au Bac, le commentaire de texte est noté sur trois critères : \textbf{compréhension} (idées justes, citations pertinentes), \textbf{langue} (correction, richesse lexicale, structures variées) et \textbf{pertinence du point de vue personnel}. Un avis clair, argumenté et ancré dans une réalité connue (votre pays, votre quartier, votre lycée) vaut mieux qu'un long résumé. Comparer la situation du Cameroun (lois, coutumes, associations, éducation) avec celle de l'Espagne ou de l'Amérique latine (lois intégrales, juzgados de violencia sobre la mujer, movimientos como \esp{Ni una menos}) montre votre capacité à \textbf{contextualiser} et à \textbf{réfléchir en citoyen}. Exemple d'ouverture : \esp{En Camerún, la ley castiga la violencia doméstica, pero muchas mujeres no denuncian por presión social, como en algunos países hispanos donde el « qué dirán » frena la denuncia.}
\end{savaistu}

\courssub{Plan vierge à remplir — Figure récapitulative}

Complétez ce schéma au brouillon avant de rédiger. Chaque case correspond à un paragraphe de votre commentaire.

\begin{popfigure}
\begin{tikzpicture}[
  box/.style={draw=popblue, thick, fill=popblueL, rounded corners=6pt, minimum height=3.2cm, minimum width=7.2cm, align=center, font=\small},
  title/.style={font=\bfseries\small, text=popdark, anchor=south, yshift=2pt},
  arrow/.style={-Stealth, thick, poporange}
]
\node[box, align=center] (p1) at (0,0){
  \textbf{Presentación}\\[2pt]
  Naturaleza del texto:\\
  Fuente / Autor:\\
  Tema central:\\
  Tesis del autor:
};
\node[box, align=center] (p2) at (8.5,0){
  \textbf{Idea 1 + Cita}\\[2pt]
  Idea directriz:\\
  Cita textual: « ... »\\
  Análisis / Explicación:\\
  Conector: \emph{Además} ...
};
\node[box, align=center] (p3) at (0,-4.2){
  \textbf{Idea 2 + Cita}\\[2pt]
  Idea directriz:\\
  Cita textual: « ... »\\
  Análisis / Explicación:\\
  Conector: \emph{Por otra parte} ...
};
\node[box, align=center] (p4) at (8.5,-4.2){
  \textbf{Opinión personal + Conclusión}\\[2pt]
  \emph{En mi opinión} ... (Camerún / mundo hispano)\\
  Argumento 1:\\
  Argumento 2 / Propuesta:\\
  \emph{Para concluir} ... \emph{Es tarea de todos}
};
\draw[arrow] (p1.east) -- (p2.west);
\draw[arrow] (p1.south) -- (p3.north);
\draw[arrow] (p2.south) -- (p4.north);
\draw[arrow] (p3.east) -- (p4.west);
\end{tikzpicture}
\legende{Plan vierge du commentaire de texte — À remplir au brouillon avant la rédaction.}
\end{popfigure}

\courssub{Exercice résolu : production guidée pas à pas}

\begin{exoresolu}[Rédiger la présentation et la première idée directrice]
\textbf{Consigne :} À partir de l'extrait de texte suivant, rédigez (1) la présentation du texte en 3-4 lignes, (2) la première idée directrice avec une citation intégrée et son analyse (5-6 lignes). Utilisez au moins 3 mots du lexique : \cle{violencia de género}, \cle{machismo}, \cle{denuncia}.

\begin{textebox}[Extrait de texte support (article fictif)]
En España, la violencia de género sigue siendo una lacra social. En 2023, más de 30.000 mujeres presentaron denuncia por maltrato. Los expertos coinciden: el machismo es la raíz del problema. « El machismo enseña al hombre a dominar y a la mujer a obedecer », afirma la socióloga María López. Las leyes de igualdad han avanzado, pero la cultura patriarcal persiste en los hogares, en el trabajo y en la escuela.
\end{textebox}

\tcblower
\textbf{Solution détaillée :}

\textbf{(1) Présentation :}
\esp{El texto es un artículo periodístico que trata de la violencia de género en España. El autor denuncia la persistencia del machismo como causa principal del maltrato y cita datos oficiales sobre las denuncias presentadas en 2023.}

\textbf{(2) Première idée directrice :}
\esp{La primera idea directriz es que el machismo es la raíz cultural de la violencia de género. El texto incluye la declaración de la socióloga María López: « El machismo enseña al hombre a dominar y a la mujer a obedecer ». Esta cita muestra que la violencia no es un acto individual aislado, sino el resultado de una socialización diferenciada que legitima el control del hombre sobre la mujer. Además, el autor precisa que, aunque las leyes de igualdad han avanzado, la cultura patriarcal persiste en los hogares, en el trabajo y en la escuela.}

\textbf{Analyse de la solution :}
\begin{itemize}
  \item Présentation : nature (artículo), thème (violencia de género), thèse (machismo causa principal), donnée chiffrée (30.000 denuncias).
  \item Idée 1 : phrase thèse + citation intégrée avec guillemets « » + analyse (socialisation, contrôle) + lien avec la persistance malgré les lois.
  \item Lexique réemployé : \cle{violencia de género}, \cle{machismo}, \cle{denuncia}, \cle{maltrato}, \cle{leyes de igualdad}, \cle{cultura patriarcal} (6 mots minimum).
  \item Connecteurs : \esp{Además} (ajout), \esp{pero} (concession), \esp{que} (subordonnée).
\end{itemize}
\end{exoresolu}

\courssub{Texte support original pour entraînement — « Romper el silencio »}

Voici un texte original (environ 250 mots) sur le même thème, rédigé pour vous permettre de vous entraîner au commentaire complet. Il intègre le lexique étudié et des références au contexte camerounais et hispanique.

\begin{textebox}[Romper el silencio — Artículo de opinión original]
En Yaoundé, como en Madrid o en Buenos Aires, el silencio es el mejor aliado del maltrato. Cada día, miles de mujeres sufren violencia de género en la intimidad de sus hogares, lejos de las cámaras y de las estadísticas oficiales. El machismo, esa ideología que considera a la mujer inferior al hombre, no conoce fronteras: se expresa en los insultos, en el control del dinero, en la prohibición de trabajar o de estudiar, y en los golpes que dejan marcas invisibles en el alma.

En el mundo hispano, el movimiento \emph{Ni una menos} ha logrado visibilizar el feminicidio y presionar a los gobiernos para que aprueben leyes integrales. España cuenta con juzgados especializados y la ley orgánica 1/2004; Argentina ha creado el Ministerio de las Mujeres, Géneros y Diversidad. Sin embargo, la impunidad sigue siendo alta: muchas víctimas retiran la denuncia por miedo, por dependencia económica o por presión familiar.

En Camerún, la ley castiga la violencia doméstica desde 2016, pero las costumbres pesan más que los códigos. En algunos pueblos, la mujer que denuncia a su esposo es rechazada por su propia familia; el « qué dirán » comunitario actúa como una cárcel sin barrotes. Las asociaciones como \emph{ALVF} (Association de Lutte contre les Violences faites aux Femmes) o \emph{RENATA} (Réseau National des Associations de Tantines) acompañan a las víctimas, pero les faltan recursos.

La educación es la única vía duradera. Enseñar la igualdad desde la escuela primaria, formar a la policía y a los jueces, y involucrar a los hombres como aliados son pasos indispensables. Romper el silencio no es solo un acto de valentía individual: es una exigencia colectiva de justicia.
\end{textebox}

\noindent\textbf{Glossaire des mots difficiles :}
\begin{itemize}
  \item \esp{lacra} : fléau, plaie sociale
  \item \esp{visibilizar} : rendre visible, médiatiser
  \item \esp{feminicidio} : féminicide (meurtre d'une femme parce qu'elle est femme)
  \item \esp{ley orgánica} : loi organique (loi fondamentale en Espagne)
  \item \esp{impunidad} : impunité
  \item \esp{retirar la denuncia} : retirer sa plainte
  \item \esp{dependencia económica} : dépendance économique
  \item \esp{qué dirán} : « ce qu'ils diront » (pression de l'opinion publique)
  \item \esp{aliados} : alliés
  \item \esp{exigencia colectiva} : exigence collective
\end{itemize}

\noindent\textbf{Questions de compréhension (pour préparer le commentaire) :}
\begin{enumerate}
  \item Quel est le thème central du texte ? Citez la phrase qui l'exprime le mieux.
  \item Relevez deux manifestations du \cle{machismo} mentionnées dans le premier paragraphe.
  \item Quelles réponses institutionnelles sont citées pour l'Espagne et l'Argentine ?
  \item Pourquoi, selon le texte, la loi camerounaise de 2016 reste-t-elle peu efficace ?
  \item Quel rôle jouent les associations \emph{ALVF} et \emph{RENATA} ?
  \item Quelle solution durable l'auteur propose-t-il ? Justifiez par une citation.
\end{enumerate}

\courssub{Attention : pièges fréquents des francophones au Bac}

\begin{attention}[Erreurs à éviter absolument dans le commentaire]
\begin{itemize}
  \item \textbf{Mots français déguisés :} N'écrivez jamais \esp{violencia} (correct) → \esp{violence} (incorrect) ; \esp{denuncia} → \esp{dénonciation} ; \esp{igualdad} → \esp{égalité} ; \esp{maltrato} → \esp{maltraitance} ; \esp{acoso} → \esp{harcèlement} ; \esp{víctima} → \esp{victime} (sauf si le mot est identique : \esp{víctima} s'écrit avec accent). \textbf{Au Bac, tout mot français non adapté fait perdre des points de langue.}
  \item \textbf{Guillemets :} Utilisez \esp{« ... »} (Alt+0171 / Alt+0187) ou \esp{" ... "}. Jamais les guillemets français \og ... \fg{} dans un texte espagnol.
  \item \textbf{Subjonctif après opinion/nécessité :} \esp{Es importante que la mujer \textbf{denuncie}} (subj. présent), \esp{Es necesario que los hombres \textbf{cambien}} (subj. présent). Indicatif \esp{creo que} / \esp{pienso que} + indicatif ; \esp{no creo que} / \esp{dudo que} + subjonctif.
  \item \textbf{Accents sur les monosyllabes :} \esp{sí} (oui) ≠ \esp{si} (si) ; \esp{él} (lui) ≠ \esp{el} (le) ; \esp{tú} (toi) ≠ \esp{tu} (ton) ; \esp{mí} (moi) ≠ \esp{mi} (mon) ; \esp{dé} (donne - subj.) ≠ \esp{de} (de).
  \item \textbf{Faux ami \emph{realizar} :} \esp{realizar} = accomplir, réaliser (un projet), \textbf{pas} « réaliser » au sens de « comprendre » (→ \esp{darse cuenta de}).
  \item \textbf{Ser / Estar :} \esp{La violencia \textbf{es} un problema estructural} (identité) ; \esp{La mujer \textbf{está} sufriendo} (action en cours) ; \esp{Está \textbf{maltratada}} (état résultant).
\end{itemize}
\end{attention}

\courssub{Exercices d'entraînement progressif}

\begin{exercice}[Entraînement 1 — Compléter la présentation]
Complétez ce paragraphe de présentation avec les mots manquants (nature, thème, thèse, chiffre) :
\esp{El texto es un \_\_\_\_\_\_\_\_ que trata de la \_\_\_\_\_\_\_\_ en el mundo hispano. El autor denuncia que el \_\_\_\_\_\_\_\_ es la causa principal. Cita el dato de \_\_\_\_\_\_\_\_ denuncias en 2023.}
\begin{itemize}
  \item Mots à placer : \esp{artículo periodístico}, \esp{violencia de género}, \esp{machismo}, \esp{más de 30.000}
\end{itemize}

\end{exercice}

\begin{exercice}[Entraînement 2 — Intégrer une citation]
Transformez ces deux phrases en une seule phrase avec citation intégrée :
\esp{El autor afirma una idea importante. « La denuncia es el primer paso hacia la libertad ».}
\begin{itemize}
  \item Réponse attendue : \esp{El autor afirma que « la denuncia es el primer paso hacia la libertad ».}
\end{itemize}

\end{exercice}

\begin{exercice}[Entraînement 3 — Rédiger l'opinion personnelle (8-10 lignes)]
Rédigez votre opinion personnelle sur le texte \emph{Romper el silencio} en respectant la méthode (avis + argumentation + ouverture Cameroun + proposition). Utilisez au moins 5 mots du lexique : \cle{violencia de género}, \cle{machismo}, \cle{denuncia}, \cle{igualdad}, \cle{educación}, \cle{asociación}, \cle{impunidad}, \cle{ley}. Commencez par : \esp{En mi opinión, ...}
\end{exercice}

\begin{corrige}[Exercice 1]
\esp{El texto es un \textbf{artículo periodístico} que trata de la \textbf{violencia de género} en el mundo hispano. El autor denuncia que el \textbf{machismo} es la causa principal. Cita el dato de \textbf{más de 30.000} denuncias en 2023.}
\end{corrige}

\begin{corrige}[Exercice 2]
\esp{El autor afirma que « la denuncia es el primer paso hacia la libertad ».}
\end{corrige}

\begin{corrige}[Exercice 3 — Proposition de correction]
\esp{En mi opinión, el texto refleja una realidad que también se vive en Camerún. La violencia de género no es solo un problema español o argentino: en nuestros barrios, muchas mujeres sufren el machismo en silencio. La impunidad persiste porque las víctimas temen retirar la denuncia por presión familiar o dependencia económica. Estoy de acuerdo con el autor: la educación es la clave. Si desde la escuela se enseña la igualdad y el respeto, los niños de hoy no serán los maltratadores de mañana. Las asociaciones como ALVF o RENATA hacen un trabajo esencial, pero necesitan más recursos y el apoyo real de la ley. Los hombres deben ser aliados, no enemigos. Para concluir, romper el silencio es un deber colectivo: solo juntos podremos construir una sociedad sin violencia.}
\end{corrige}

\newpage

\coursec{Activité d'intégration}

\coursec{Lectura y comentario : discriminations et violences contre les femmes (la familia hispánica)}

\courssub{Objectifs de l'activité}

À l'issue de cette activité d'intégration, l'élève sera capable de :
\begin{itemize}
\item \cle{Réemployer} dans un contexte authentique le lexique de la condition féminine étudié dans la leçon : \esp{violencia de género}, \esp{machismo}, \esp{maltrato}, \esp{denuncia}, \esp{igualdad}, \esp{acoso} ;
\item \cle{Produire} un texte court et engagé (slogan, phrases de sensibilisation, paragraphe d'opinion) en espagnol correct ;
\item \cle{Mobiliser} la méthode du commentaire personnel (prise de position argumentée, comparaison) à l'échelle d'un court paragraphe ;
\item \cle{S'exprimer à l'oral} pendant deux minutes pour présenter et défendre une production collective ;
\item \cle{Travailler en équipe} et évaluer la production des autres groupes selon des critères explicites.
\end{itemize}

\courssub{Situation}

Dans le cadre de la semaine du \esp{8 de marzo}, \esp{Día Internacional de la Mujer}, le club d'espagnol de ton lycée de Yaoundé a été sollicité par la direction. Le proviseur souhaite décorer le hall de l'établissement avec des affiches de sensibilisation rédigées en espagnol, en lien avec la campagne internationale \esp{« Rompe el silencio »} menée dans de nombreux pays hispanophones. Les meilleures affiches seront également envoyées à l'ambassade d'Espagne et publiées sur la page du club. Voici le message que la présidente du club a reçu et qu'elle vous lit en classe :

\begin{textebox}[Message de la coordinadora de la campaña]
Estimados alumnos y alumnas del club de español:

Con motivo del 8 de marzo, Día Internacional de la Mujer, lanzamos la campaña « Rompe el silencio ». Os invitamos a crear carteles de sensibilización contra la violencia de género. Cada cartel debe incluir un eslogan original, cinco frases con el vocabulario del tema (violencia de género, machismo, denuncia, igualdad, acoso) y un breve comentario personal que compare la situación de la mujer en Camerún y en el mundo hispánico.

Recordad: el silencio protege al agresor, nunca a la víctima. Denunciar es un derecho y un deber. ¡Vuestra voz puede cambiar las cosas!

Un saludo cordial,\\
La coordinadora
\end{textebox}

\begin{definition}[Glossaire du message]
\begin{itemize}
\item \esp{con motivo de} : à l'occasion de ;
\item \esp{cartel de sensibilización} : affiche de sensibilisation ;
\item \esp{eslogan} : slogan ;
\item \esp{el agresor / la víctima} : l'agresseur / la victime ;
\item \esp{un derecho y un deber} : un droit et un devoir.
\end{itemize}
\end{definition}

\courssub{Consignes}

Travail en groupes de 4 élèves. Durée totale : 45 minutes de préparation, puis 2 minutes de présentation orale par groupe.

\begin{enumerate}
\item \textbf{Relecture et repérage (5 min).} Relisez le message de la coordinadora et le texte support de la leçon. Soulignez dans les deux documents les mots du lexique de la condition féminine et notez deux idées fortes que vous pourriez reprendre.

\item \textbf{Création du slogan (5 min).} Inventez un slogan court, percutant et original en espagnol. Il doit contenir un impératif (par exemple : \esp{Rompe}, \esp{Denuncia}, \esp{Alza la voz}) ou une formule frappante. Interdiction de recopier tel quel le slogan de la campagne.

\item \textbf{Rédaction des cinq phrases (10 min).} Rédigez cinq phrases complètes et correctes utilisant chacune un des cinq mots-clés : \esp{violencia de género}, \esp{machismo}, \esp{denuncia}, \esp{igualdad}, \esp{acoso}. Chaque phrase doit avoir un sens complet (pas une simple définition) et viser à sensibiliser le lecteur.

\item \textbf{Commentaire personnel (15 min).} Rédigez ensemble un paragraphe de 5 à 6 lignes comparant la situation de la femme au Cameroun et dans le monde hispanique (Espagne, Amérique latine, Guinée équatoriale). Structure attendue : une phrase d'introduction (constat), deux arguments ou exemples, une phrase de conclusion avec votre opinion (\esp{en mi opinión}, \esp{creo que}, \esp{es necesario que} + subjonctif).

\item \textbf{Mise en page de l'affiche (5 min).} Disposez sur la feuille : le slogan en grand, les cinq phrases, le paragraphe, et un dessin ou symbole simple (poing levé, ruban, silhouette).

\item \textbf{Présentation orale (2 min par groupe).} Chaque membre du groupe prend la parole : un élève présente le slogan, deux élèves lisent et expliquent les phrases, un élève lit le commentaire. Parlez fort, regardez l'auditoire.

\item \textbf{Vote de la classe.} Chaque élève vote (hors de son groupe) selon trois critères : correction de la langue, force du message, qualité de la présentation orale.
\end{enumerate}

\begin{attention}[Pièges à éviter dans vos phrases]
\begin{itemize}
\item Ne traduis pas « les violences faites aux femmes » mot à mot : dis \esp{la violencia contra las mujeres} ou \esp{la violencia de género}.
\item \esp{Denunciar} se construit sans préposition devant la chose : \esp{denunciar el maltrato} (et non \esp{denunciar de el maltrato}).
\item Attention au faux ami : \esp{una denuncia} = une dénonciation / une plainte (juridique), pas « une dénonciation » au sens de calomnie.
\item Après \esp{es necesario que} + sujet différent, emploie le subjonctif : \esp{Es necesario que la sociedad reaccione.}
\item N'oublie pas l'accord : \esp{la igualdad entre hombres y mujeres}, \esp{las víctimas tienen derechos}.
\item Ponctuation espagnole : pas d'espace avant les deux-points (\esp{denuncia:}) ni avant le point-virgule (\esp{;}).
\end{itemize}
\end{attention}

\courssub{Ressources à mobiliser}

\begin{aretenir}[Boîte à outils de la campagne]
\textbf{Lexique :} \esp{la violencia de género}, \esp{el machismo}, \esp{el maltrato}, \esp{maltratar}, \esp{la denuncia / denunciar}, \esp{la igualdad}, \esp{el acoso (escolar, sexual, laboral)}, \esp{la víctima}, \esp{el agresor}, \esp{los derechos de la mujer}, \esp{la discriminación}, \esp{luchar contra}, \esp{proteger a las víctimas}, \esp{alzar la voz}, \esp{romper el silencio}.

\textbf{Structures d'opinion :} \esp{en mi opinión}, \esp{creo / pienso que + indicatif}, \esp{me parece que}, \esp{es necesario / importante / urgente que + subjonctif}, \esp{hay que + infinitif}.

\textbf{Structures de comparaison :} \esp{en Camerún... mientras que en España...}, \esp{al contrario que}, \esp{tanto en... como en...}, \esp{sin embargo}, \esp{no obstante}.

\textbf{Impératifs utiles pour les slogans (tú) :} \esp{Rompe el silencio}, \esp{Denuncia}, \esp{Alza la voz}, \esp{No te calles}, \esp{Actúa}, \esp{Di no a la violencia}.

\textbf{Méthode du commentaire personnel :} constat $\rightarrow$ arguments ou exemples $\rightarrow$ prise de position $\rightarrow$ conclusion.
\end{aretenir}

\courssub{Proposition de production et corrigé commenté}

\begin{exoresolu}[Affiche modèle du groupe « Las Voces Libres »]
Énoncé : Rédiger l'affiche complète (slogan, cinq phrases, commentaire personnel) attendue pour la campagne \esp{« Rompe el silencio »}.
\tcblower
\textbf{Eslogan :}\\
\esp{« Ni silencio ni miedo: denuncia y vive libre. »}

\textbf{Las cinco frases :}
\begin{enumerate}
\item \esp{La violencia de género es un problema de toda la sociedad, no solo de las mujeres.} « La violence de genre est un problème de toute la société, pas seulement des femmes. »
\item \esp{El machismo se combate desde la escuela, educando a los niños y a las niñas en el respeto.} « Le machisme se combat dès l'école, en éduquant les garçons et les filles au respect. »
\item \esp{Si sufres maltrato, presenta una denuncia: callar protege al agresor.} « Si tu subis de la maltraitance, dépose une plainte: se taire protège l'agresseur. »
\item \esp{La igualdad entre hombres y mujeres es un derecho, no un regalo.} « L'égalité entre hommes et femmes est un droit, pas un cadeau. »
\item \esp{El acoso escolar deja heridas invisibles: hablemos de ello.} « Le harcèlement scolaire laisse des blessures invisibles: parlons-en. »
\end{enumerate}

\textbf{Comentario personal :}\\
\esp{En mi opinión, la situación de la mujer ha mejorado tanto en Camerún como en el mundo hispánico, pero queda mucho por hacer. En nuestro país, muchas mujeres trabajan en los mercados y sostienen a la familia; sin embargo, algunas sufren maltrato y no se atreven a denunciar. En España y en América Latina existen leyes contra la violencia de género, pero el machismo sigue presente en la vida cotidiana. Creo que es necesario que los jóvenes, aquí y allá, alcemos la voz. Romper el silencio es el primer paso hacia la igualdad.}

\textit{Traduction : « À mon avis, la situation de la femme s'est améliorée tant au Cameroun que dans le monde hispanique, mais il reste beaucoup à faire. Dans notre pays, beaucoup de femmes travaillent dans les marchés et soutiennent la famille ; cependant, certaines subissent de la maltraitance et n'osent pas porter plainte. En Espagne et en Amérique latine, il existe des lois contre la violence de genre, mais le machisme reste présent dans la vie quotidienne. Je crois qu'il est nécessaire que les jeunes, ici et là-bas, élevions la voix. Briser le silence est le premier pas vers l'égalité. »}

\textbf{Commentaire des choix de langue (méthode) :}
\begin{itemize}
\item Le slogan utilise deux noms (\esp{silencio}, \esp{miedo}) et deux impératifs (\esp{denuncia}, \esp{vive}) : rythme binaire, facile à mémoriser, sans recopier le slogan officiel. Respect de la ponctuation espagnole (pas d'espace avant les deux-points).
\item Chaque phrase contient bien le mot-clé exigé, employé avec sa construction correcte : \esp{presentar una denuncia}, \esp{la igualdad entre... y...}, \esp{el acoso escolar}.
\item Le commentaire suit la structure du commentaire personnel : constat (\esp{ha mejorado... pero queda mucho por hacer}), comparaison avec \esp{tanto... como} et \esp{sin embargo} (précédé d'un point-virgule), opinion avec \esp{en mi opinión} et \esp{creo que}, et subjonctif après \esp{es necesario que} : \esp{alcemos} (1re personne du pluriel, \emph{nosotros}).
\item Note culturelle : le passé composé \esp{ha mejorado} (pretérito perfecto compuesto) et la tournure \esp{queda mucho por hacer} sont typiques de l'espagnol de presse ; \esp{no se atreven a denunciar} illustre la construction \esp{atreverse a + infinitif}.
\end{itemize}
\end{exoresolu}

\courssub{Bilan de l'activité}

Cette activité a permis de vérifier l'acquisition des trois savoir-faire de la leçon : la \cle{compréhension} (relecture ciblée des documents), le \cle{réemploi du lexique} (les cinq mots-clés dans des phrases de sensibilisation) et la \cle{production d'un commentaire personnel} structuré et comparatif. L'oral a en outre entraîné la prise de parole en public, compétence essentielle à l'examen.

Pour aller plus loin, la classe peut :
\begin{itemize}
\item afficher les productions dans le hall et organiser une permanence d'information le 8 mars ;
\item transformer le meilleur commentaire en article pour le journal du lycée, en appliquant la méthode complète du commentaire de texte (présentation, idées directrices, commentaire personnel) ;
\item préparer un débat : \esp{« ¿Son suficientes las leyes para acabar con el machismo ? »}, qui servira d'entraînement à l'expression orale en interaction.
\end{itemize}

\begin{savaistu}[Le 8 mars dans le monde hispanophone]
Depuis quelques années, les manifestations du \esp{8 de marzo} rassemblent des millions de personnes en Espagne, au Mexique, en Argentine ou au Chili, avec des mots d'ordre comme \esp{« Ni una menos »} (pas une de moins), né en Argentine contre les féminicides. En Guinée équatoriale, pays hispanophone voisin du Cameroun, la journée de la femme est aussi l'occasion de campagnes officielles pour l'éducation des filles. En participant à cette affiche, tu rejoins donc un véritable mouvement de la francophonie... hispanique !
\end{savaistu}

\newpage

\coursec{Méthodes et exercices résolus}

\courssub{Méthode : Lecture active et compréhension globale d'un texte d'actualité}

\begin{methode}[Lire, annoter, comprendre un texte journalistique sur la violence de genre]
\textbf{Objectif :} Extraire l'information essentielle, identifier la thèse de l'auteur et la structure argumentative.
\begin{enumerate}
    \item \textbf{Paratexte :} Identifier la source (journal, blog institutionnel, ONG), la date, le titre, les intertitres, les images. \cle{Question clé :} « Quel est le public visé ? Quel est le but : informer, dénoncer, sensibiliser ? »
    \item \textbf{Première lecture (globale) :} Lire sans s'arrêter au vocabulaire inconnu. Repérer le \cle{sujet} (violencia de género), la \cle{thèse} (ex: « La loi ne suffit pas, il faut éduquer ») et le \cle{ton} (indigné, objectif, espoir).
    \item \textbf{Deuxième lecture (détaillée, crayon en main) :}
        \begin{itemize}
            \item Souligner les \cle{mots-clés} récurrents (mujer, denuncia, impunidad, patriarcado, ley).
            \item Entourer les \cle{connecteurs logiques} (por tanto, sin embargo, además, en consecuencia) pour voir l'enchaînement des idées.
            \item Noter en marge l'idée-force de chaque paragraphe (une phrase en français ou en espagnol).
        \end{itemize}
    \item \textbf{Troisième lecture (vérification) :} Reconstituer le plan du texte (Introduction / Arguments / Conclusion). Vérifier la compréhension des données chiffrées (dates, pourcentages) et des noms propres (lois, pays).
    \item \textbf{Glossaire personnel :} Lister 10--15 mots nouveaux, les traduire, les classer par famille (verbes d'action, noms abstraits, adjectifs qualificatifs).
\end{enumerate}
\textbf{Reflexe Bac :} Ne jamais traduire mot à mot. Chercher le sens par le contexte (champ lexical, cognats, structure de la phrase).
\end{methode}

\begin{exoresolu}[Compréhension détaillée : « Ni una menos : la lucha continúa »]
\textbf{Énoncé :} Lisez le texte ci-dessous et répondez aux questions en espagnol (phrases complètes).

\begin{textebox}[Texte original pour l'exercice -- Simulacro de prueba de prensa]
La violencia machista no es un problema privado, es una pandemia social que atraviesa todas las clases, edades y fronteras. En España, desde 2003, más de 1.200 mujeres han sido asesinadas por sus parejas o exparejas. En América Latina, la cifra es escalofriante: catorce de los veinticinco países con las tasas más altas de feminicidio del mundo se encuentran en esta región. 

Ante esta realidad, los movimientos feministas han logrado avances legislativos históricos. La \emph{Ley Orgánica 1/2004 de Medidas de Protección Integral contra la Violencia de Género} en España o la tipificación del \emph{feminicidio} en los códigos penales de México, Argentina o Chile son hitos fundamentales. Sin embargo, las leyes no bastan si la justicia es lenta, si los recursos para las víctimas son insuficientes y si la educación sigue reproduciendo estereotipos de dominación. 

El \emph{machismo} mata, pero también el silencio de las instituciones y la indiferencia de la sociedad. La denuncia es el primer paso, pero la protección real exige juzgados especializados, casas de acogida financiadas y una educación afectivo-sexual obligatoria desde la infancia. \textbf{Ni una menos, vivas nos queremos.}
\end{textebox}

\textbf{Questions :}
\begin{enumerate}
    \item ¿Cómo define el autor la violencia machista en el primer párrafo? Cite dos características.
    \item ¿Qué dato numérico se da sobre España y qué comparación se hace con América Latina?
    \item ¿Cuáles son los dos limites de la legislación actual mencionados en el segundo párrafo?
    \item Según el tercer párrafo, ¿qué tres medidas concretas propone el autor para una protección real?
    \item Explique el sentido del eslogan final: « \emph{Ni una menos, vivas nos queremos} ».
\end{enumerate}

\tcblower
\textbf{Solution détaillée :}

\textbf{Raisonnement général :} Les réponses doivent être des phrases complètes en espagnol, reprenant les mots du texte ou les reformulant avec le vocabulaire de la leçon. On évite le « Sí / No ».

\textbf{1.} \emph{El autor define la violencia machista como una « pandemia social » (no un problema privado) que « atraviesa todas las clases, edades y fronteras ».}
\begin{itemize}
    \item \textit{Justification :} Reprise des termes « pandemia social » et de l'énumération « clases, edades y fronteras ». Attention à l'accord de l'adjectif « privada » (féminin) si on cite « problema privado ».
\end{itemize}

\textbf{2.} \emph{En España, desde 2003, más de 1.200 mujeres han sido asesinadas por sus parejas o exparejas. En América Latina, catorce de los veinticinco países con las tasas más altas de feminicidio del mundo se encuentran en esta región.}
\begin{itemize}
    \item \textit{Grammaire :} Passif « han sido asesinadas » (accord du participe avec « mujeres »). « Tasas más altas » (comparatif de supériorité).
\end{itemize}

\textbf{3.} \emph{Las leyes no bastan porque la justicia es lenta, los recursos para las víctimas son insuficientes y la educación sigue reproduciendo estereotipos de dominación.}
\begin{itemize}
    \item \textit{Synthèse :} Trois causes : justice lente, ressources insuffisantes, éducation stéréotypée. Utilisation de « porque » pour lier.
\end{itemize}

\textbf{4.} \emph{Propone juzgados especializados, casas de acogida financiadas y una educación afectivo-sexual obligatoria desde la infancia.}
\begin{itemize}
    \item \textit{Vocabulaire :} « Juzgados especializados », « casas de acogida » (centres d'hébergement), « educación afectivo-sexual ».
\end{itemize}

\textbf{5.} \emph{El eslogan « Ni una menos, vivas nos queremos » significa que no se acepta ni una sola muerte más de mujer por violencia de género (\"Ni una menos\") y que la reivindicación central es el derecho a la vida y a la libertad de las mujeres (\"Vivas nos queremos\"), frente a la cultura de la muerte que impone el patriarcado.}
\begin{itemize}
    \item \textit{Analyse :} « Ni una menos » = « Pas une de moins » (référence au mouvement argentin). « Vivas nos queremos » = « Nous nous voulons vivantes » (indicatif présent \emph{queremos} : affirmation d'une volonté collective, structure \emph{querer + infinitif} avec pronom réfléchi « nos » et adjectif prédicatif « vivas »).
\end{itemize}

\textbf{Conclusion :} Le texte articule constat chiffré, analyse juridique et revendication sociale. Le ton est engagé (lexique : « escalofriante », « silencio », « indiferencia »).
\end{exoresolu}

\courssub{Méthode : Repérage et appropriation du lexique thématique (famille, genre, violences)}

\begin{methode}[Construire son répertoire lexical actif : classer, dériver, contextualiser]
\textbf{Objectif :} Passer de la compréhension passive à l'expression écrite/orale précise.
\begin{enumerate}
    \item \textbf{Repérage dans le texte :} Surbrillance de trois catégories :
        \begin{itemize}
            \item \cle{Noms d'acteurs / rôles :} \esp{la víctima, el agresor, el maltratador, la denunciante, el testigo, los hijos/as expuestos}.
            \item \cle{Noms d'actes / concepts juridiques :} \esp{la denuncia, la orden de protección, el feminicidio, la violencia vicaria, la custodia, la patria potestad}.
            \item \cle{Adjectifs / États :} \esp{vulnerable, impune, revictimizada, machista, coeducativo, igualitario}.
        \end{itemize}
    \item \textbf{Travail morphologique (dérivation) :} Transformer pour enrichir.
        \begin{center}
        \begin{tabularx}{\linewidth}{|X|X|X|}\hline
        \rowcolor{popblueL} \textbf{Verbe} & \textbf{Nom (action/état)} & \textbf{Adjectif / Participe} \\ \hline
        \esp{maltratar} & \esp{el maltrato} & \esp{maltratado / maltratador} \\ \hline
        \esp{denunciar} & \esp{la denuncia} & \esp{denunciado / denunciante} \\ \hline
        \esp{acosar} & \esp{el acoso} & \esp{acosado / acosador} \\ \hline
        \esp{discriminar} & \esp{la discriminación} & \esp{discriminado / discriminatorio} \\ \hline
        \esp{igualar} & \esp{la igualdad} & \esp{igual / igualitario} \\ \hline
        \esp{proteger} & \esp{la protección} & \esp{protegido / protector} \\ \hline
        \end{tabularx}
        \end{center}
    \item \textbf{Collocations (associations figées) à mémoriser :}
        \begin{itemize}
            \item \esp{presentar / interponer una denuncia} (pas *faire une dénonciation).
            \item \esp{dictar una orden de protección / de alejamiento}.
            \item \esp{romper el silencio / el ciclo de la violencia}.
            \item \esp{erradicar la violencia de género / el machismo}.
            \item \esp{garantizar los derechos de las víctimas}.
        \end{itemize}
    \item \textbf{Faux amis \& Pièges (Francophone) :}
        \begin{itemize}
            \item \cle{Violación} = viol (sexuel) ; \emph{ne pas confondre avec} \esp{violencia} (violence). \esp{Violación de derechos} = violation de droits.
            \item \cle{Maltrato} = maltraitance (physique/psychologique) ; \esp{maltratar} = maltraiter.
            \item \cle{Acoger} = accueillir / héberger ; \esp{casa de acogida} = centre d'hébergement / foyer. \emph{Pas *acogida* comme nom principal pour le lieu}.
            \item \cle{Género} = genre (social) ; \emph{sexe biologique} = \esp{sexo}.
        \end{itemize}
    \item \textbf{Entraînement actif :} Écrire 5 phrases personnelles utilisant 5 collocations différentes. Les faire corriger.
\end{enumerate}
\end{methode}

\begin{exoresolu}[Réemploi lexical : classification, dérivation, phrase contextuelle]
\textbf{Énoncé :}
\begin{enumerate}
    \item \textbf{Classification :} Rangez les mots suivants dans le tableau ci-dessous selon leur catégorie grammaticale et sémantique : \esp{feminicidio, denunciar, vulnerables, protección, acosador, igualdad, maltrato, coeducativo, impunidad, víctima, agresor, legislación, romper, ciclo, estereotipo}.
    \item \textbf{Dérivation :} Complétez le tableau pour les verbes suivants (Nom d'action / Adjectif qualificatif / Agent).
    \item \textbf{Mise en contexte :} Rédigez en espagnol un paragraphe de 5 lignes minimum sur « La importancia de la denuncia » en utilisant \textbf{au moins 8 mots} de la liste ci-dessus (toutes catégories confondues). Soulignez les mots utilisés.
\end{enumerate}

\tcblower
\textbf{Solution détaillée :}

\textbf{1. Classification :}
\begin{center}
\begin{tabularx}{\linewidth}{|X|X|X|X|}\hline
\rowcolor{popblueL} \textbf{Noms (Concepts/Actes)} & \textbf{Noms (Acteurs/Rôles)} & \textbf{Verbes (Action)} & \textbf{Adjectifs / Qualificatifs} \\ \hline
\esp{feminicidio, protección, igualdad, maltrato, impunidad, legislación, ciclo, estereotipo} & \esp{víctima, agresor, acosador, denunciante (implícito)} & \esp{denunciar, romper, acosar (implícito), maltratar (implícito), legislar (implícito)} & \esp{vulnerables, coeducativo, igualitario (implícito), machista (implícito)} \\ \hline
\end{tabularx}
\end{center}
\textit{Note :} \esp{acosador} est un nom d'agent (qui peut fonctionner comme adjectif). \esp{vulnerables} est adjectif (invariant en genre, marque pluriel).

\textbf{2. Dérivation :}
\begin{center}
\begin{tabular}{|l|l|l|l|}\hline
\rowcolor{popblueL} \textbf{Verbe} & \textbf{Nom (Action/État)} & \textbf{Adjectif / Participe} & \textbf{Agent / Personne} \\ \hline
\esp{denunciar} & \esp{la denuncia} & \esp{denunciado} & \esp{el/la denunciante} \\ \hline
\esp{acosar} & \esp{el acoso} & \esp{acosado} & \esp{el/la acosador/a} \\ \hline
\esp{maltratar} & \esp{el maltrato} & \esp{maltratado} & \esp{el/la maltratador/a} \\ \hline
\esp{proteger} & \esp{la protección} & \esp{protegido} & \esp{el/la protector/a} \\ \hline
\esp{legislar} & \esp{la legislación} & \esp{legislado} & \esp{el/la legislador/a} \\ \hline
\esp{estereotipar} & \esp{el estereotipo} & \esp{estereotipado} & --- \\ \hline
\end{tabular}
\end{center}

\textbf{3. Production écrite (Modèle de réponse) :}
\begin{quote}
\emph{La \textbf{denuncia} es el primer paso indispensable para \textbf{romper} el \textbf{ciclo} de la \textbf{violencia} y la \textbf{impunidad}. Muchas \textbf{víctimas} no se atreven a \textbf{denunciar} a su \textbf{agresor} por miedo a la \textbf{revictimización} o a perder la \textbf{custodia} de sus hijos. Sin embargo, la \textbf{legislación} actual ofrece \textbf{protección} mediante las \textbf{órdenes de alejamiento} y los juzgados especializados. Es fundamental una \textbf{coeducación} afectivo-sexual desde la infancia para prevenir el \textbf{maltrato} y construir una sociedad \textbf{igualitaria}. \textbf{Ni una menos}.}
\end{quote}
\textbf{Vérification :} 12 mots utilisés (dépasse le minimum de 8). Accords respectés (\emph{víctimas, agresor, custodia, legislación, protección, coeducación, maltrato, igualitaria}). Subjonctif non nécessaire ici (indicatif pour faits réels). Structure \emph{Es fundamental + infinitif}.
\end{exoresolu}

\courssub{Méthode : Structure et rédaction du commentaire de texte (Épreuve du Baccalauréat)}

\begin{methode}[Le commentaire composé en 4 temps : Présenter, Analyser, Interpréter, Conclure/Ouvrir]
\textbf{Objectif :} Produire un devoir structuré (250--300 mots env.), en espagnol, noté sur 20 pts (Compréhension 8, Analyse 7, Expression 5).
\begin{enumerate}
    \item \textbf{INTRODUCTION (env. 40--50 mots) -- \cle{La "carte d'identité" + Problématique + Plan annoncé}}
        \begin{itemize}
            \item \textbf{Accroche :} Contexte général (ex: \emph{La violencia de género es una lacra global...}).
            \item \textbf{Présentation du texte :} Auteur (si connu) / Source / Date / Titre / Nature (article de presse, discours, rapport ONG).
            \item \textbf{Sujet :} De quoi parle le texte ? (Un nom + complément : \emph{la lucha contra el feminicidio en América Latina}).
            \item \textbf{Thèse :} Quelle est la position de l'auteur ? (Une phrase : \emph{El autor denuncia la impunidad y exige una educación transformadora.}).
            \item \textbf{Problématique :} Question ouverte à laquelle le texte répond (ex: \emph{¿Por qué las leyes actuales son insuficientes para erradicar la violencia machista?}).
            \item \textbf{Annonce du plan :} \emph{En primer lugar analizaremos...; en segundo lugar...; finalmente...}
        \end{itemize}
    \item \textbf{DÉVELOPPEMENT (2 ou 3 parties, env. 180--200 mots) -- \cle{Idée directrice + Citation + Analyse linguistique/thématique}}
        \begin{itemize}
            \item Une idée directrice par paragraphe (titre mental).
            \item \textbf{Citation courte} (entre guillemets espagnols \esp{« ... »}) intégrée à la phrase.
            \item \textbf{Analyse :} Expliquer \emph{comment} l'auteur dit les choses (champ lexical, figures de style, temps verbaux, connecteurs, ton). Pas de paraphrase !
            \item \textbf{Lien :} \esp{Por tanto / En consecuencia / Además / Sin embargo}.
        \end{itemize}
    \item \textbf{CONCLUSION (env. 40 mots) -- \cle{Bilan + Réponse à la problématique + Ouverture}}
        \begin{itemize}
            \item Résumer les 2--3 grands axes.
            \item Répondre explicitement à la problématique.
            \item Ouverture : lien avec l'actualité camerounaise (\esp{Ley de protección en Camerún}), autre œuvre (\esp{La casa de Bernarda Alba}), ou réflexion personnelle.
        \end{itemize}
    \item \textbf{RELECTURE (5 min) :} Accords (genre/nombre), accents (\esp{análisis, política, educación}), subjonctifs après \emph{es necesario que, quiero que, denuncio que}, connecteurs, ponctuation \esp{¿ ? ¡ !}.
\end{enumerate}
\textbf{Structure type de paragraphe d'analyse (méthode "C.I.A.") :}
\cle{C}itation -- \cle{I}nterprétation (sens) -- \cle{A}nalyse (forme/langue).
\end{methode}

\begin{exoresolu}[Commentaire modèle : « Ni una menos : la lucha continúa »]
\textbf{Énoncé :} Rédigez un commentaire structuré du texte étudié previously (environ 250 mots). Respectez la méthode en 4 temps.

\tcblower
\textbf{Solution détaillée :}

\textbf{INTRODUCTION}
\emph{La violencia de género constituye una de las violaciones más flagrantes de los derechos humanos en el mundo hispánico y más allá. El texto anónimo que comentamos, titulado « Ni una menos : la lucha continúa », parece extraído de un artículo de opinión o un manifiesto feminista reciente. Su tema central es la persistencia del feminicidio a pesar de los avances legislativos. La tesis defendida es clara: las leyes son necesarias pero insuficientes sin una transformación educativa y judicial profunda. Nos preguntaremos: \textbf{¿Por qué la respuesta legal sigue siendo ineficaz frente a la pandemia machista?} Para responder, analizaremos primero el diagnóstico alarmante que hace el autor, luego los límites de la vía jurídica y finalmente las alternativas propuestas.}

\textbf{DÉVELOPPEMENT}

\textbf{1. Un diagnóstico alarmante : la violencia como pandemia estructural.}
\emph{En primer lugar, el autor define la violencia machista no como un hecho aislado (« problema privado »), sino como una « pandemia social » que « atraviesa todas las clases, edades y fronteras ».} Le choix du mot \emph{pandemia} (champ lexical médical) souligne l'ampleur et la contagiosité du phénomène. Les chiffres (\emph{más de 1.200 mujeres... catorce de los veinticinco países}) ancrent le propos dans une réalité statistique incontestable, donnant une dimension objective à l'indignation. L'énumération \emph{clases, edades, fronteras} (polysyndeton) insiste sur l'universalité du fléau.

\textbf{2. Les limites de la voie législative : l'impunité structurelle.}
\emph{En segundo lugar, el texto reconoce « avances legislativos históricos » (Ley Orgánica 1/2004, tipificación del feminicidio), pero introduce un giro argumentativo decisivo con el conector adversativo « Sin embargo ».} La structure \emph{Si la justicia es lenta... si los recursos... son insuficientes... y si la educación sigue reproduciendo...} (trois propositions conditionnelles \emph{si} + indicatif présent pour hypothèse réelle) démontre la causalité multiple de l'échec. Le champ lexical de l'insuffisance (\emph{no bastan, lenta, insuficientes, estereotipos}) montre que la loi reste lettre morte sans moyens.

\textbf{3. Vers une protection intégrale : l'urgence éducative et institutionnelle.}
\emph{Finalmente, el autor pasa de la denuncia a la propuesta concreta.} La phrase nominale finale \emph{Juzgados especializados, casas de acogida financiadas y una educación afectivo-sexual obligatoria} dresse un programme d'action en trois piliers : justice, hébergement, prévention. L'usage de l'adjectif \emph{obligatoria} et du participe \emph{financiadas} marque la volonté politique requise. Le slogan final \emph{« Ni una menos, vivas nos queremos »} (indicatif \emph{queremos} : affirmation collective forte, « nous nous voulons vivantes ») résume l'éthique du mouvement : la vie comme droit non négociable.

\textbf{CONCLUSION}
\emph{En definitiva, el texto demuestra que la erradicación de la violencia de género exige un pacto de Estado que vaya más allá del papel legislativo. La respuesta a nuestra problématica es que la ley falla cuando la justicia, los recursos y la educación no la acompañan. Como apertura, esta reflexión resuena con fuerza en Camerún, donde la lucha contra les violences basées sur le genre (VBG) et pour l'application effective de la loi sur la protection de la femme reste un défi majeur pour la société civile et les pouvoirs publics.}

\textbf{Analyse de la correction :} Respect du plan annoncé. Citations intégrées. Vocabulaire précis (\emph{lacra, flagrante, tipificación, giro argumentativo, polysyndeton, causalité, éthique}). Connecteurs variés. Subjonctif correct (\emph{vaya} après \emph{exige que}). Ouverture contextualisée (Cameroun).
\end{exoresolu}

\courssub{Méthode : Traduction Thème / Version : pièges grammaticaux ciblés}

\begin{methode}[Stratégies pour le Thème (Français $\rightarrow$ Espagnol) et la Version (Espagnol $\rightarrow$ Français)]
\textbf{Objectif :} Éviter les calques syntaxiques et lexicaux, maîtriser les points grammaticaux "pièges" du programme Terminale (Subjonctif, Ser/Estar, Por/Para, Passifs, Discours indirect).
\begin{enumerate}
    \item \textbf{Analyse de la phrase source :} Identifier le \cle{sujet}, le \cle{verbe noyau}, les \cle{compléments}, le \cle{temps/aspect/modalité}. Détecter les \cle{mots-outils} (prépositions, conjonctions, pronoms relatifs).
    \item \textbf{Transposition grammaticale (Thème) :}
        \begin{itemize}
            \item \cle{Ser vs Estar :} \esp{Ser} = identité, définition, origine, possession, heure, matière, nationalité (\emph{Es médica, Es de Madrid, Son las dos}). \esp{Estar} = localisation, état resultant, état temporaire, progressif (\emph{Está en casa, Está cansada, Está lloviendo}). \textbf{Piège :} \esp{La mujer está maltratada} (état resultant) vs \emph{La mujer es maltratada} (passif actionnel = par quelqu'un).
            \item \cle{Subjonctif :} Déclencheur = \emph{Volonté / Émotion / Doute / Négation / Impersonnel (Es necesario que...) / Conjonctions (para que, sin que, antes de que...)}. Ex: \emph{Es urgente que el gobierno \textbf{financie} (subj.) las casas de acogida.}
            \item \cle{Por vs Para :} \esp{Por} = cause, moyen, durée, échange, agent passif, mouvement vague. \esp{Para} = but, destination, destinataire, échéance, opinion (\emph{para mí}). Ex: \emph{Luchan \textbf{por} la igualdad} (cause) / \emph{Esta ley es \textbf{para} proteger a las víctimas} (but).
            \item \cle{Passif réfléchi (Se + V 3e pers.) vs Passif sérén (Ser + Participe) :} \emph{Se denuncian muchos casos} (impersonnel/général) vs \emph{Muchos casos son denunciados por las víctimas} (agent explicite).
            \item \cle{Discours indirect :} Changement de personnes, temps (si verbe introducteur au passé), indicateurs spatio-temporels. \emph{Dice: « Denuncio »} $\rightarrow$ \emph{Dice que denuncia / denunció}.
        \end{itemize}
    \item \textbf{Recherche lexicale précise :} Éviter les faux amis (\esp{realizar} $\neq$ réaliser/réussir mais \emph{effectuer}; \esp{asistir} $\neq$ assister mais \emph{aller à}; \esp{pretender} $\neq$ prétendre mais \emph{viser/avoir l'intention}). Utiliser le dictionnaire bilingue \textbf{inversé} (Français $\rightarrow$ Espagnol) pour le thème.
    \item \textbf{Rédaction de la cible :} Écrire une phrase espagnole naturelle, pas une phrase française en mots espagnols.
    \item \textbf{Vérification (Check-list) :} Accords (GN), Accents (tildes), Ponctuation (\esp{¿ ?}), Subjonctifs, Ser/Estar, Por/Para, Prépositions verbales (\esp{tratar de, conseguir, lograr, dejar de}).
\end{enumerate}
\end{methode}

\begin{exoresolu}[Thème et Version ciblés : Violences, Loi, Éducation]
\textbf{Énoncé :}
\begin{enumerate}
    \item \textbf{Thème :} Traduisez en espagnol.
    \begin{enumerate}
        \item « Il est indispensable que le gouvernement finance davantage de centres d'hébergement pour les femmes victimes de violences. »
        \item « Les féministes luttent depuis des décennies pour que la société change ses mentalités machistes. »
        \item « Cette loi a été votée pour protéger les victimes, mais elle n'est pas appliquée correctement par les juges. »
        \item « Après avoir dénoncé son agresseur, elle a dû quitter son domicile par peur des représailles. »
    \end{enumerate}
    \item \textbf{Version :} Traduisez en français.
    \emph{« A pesar de las campañas de concienciación, el número de denuncias por acoso sexual en el entorno laboral no deja de crecer. Muchas mujeres callan porque temen perder su empleo o ser estigmatizadas por sus compañeros. Es fundamental que las empresas implementen protocolos de actuación eficaces y que la inspección de trabajo sancione severamente a los infractores. Solo así se romperá el círculo de la impunidad. »}
\end{enumerate}

\tcblower
\textbf{Solution détaillée :}

\textbf{1. Thème}

\textbf{a)} \emph{« Il est indispensable que le gouvernement finance davantage de centres d'hébergement pour les femmes victimes de violences. »}
\begin{itemize}
    \item \textbf{Analyse :} \emph{Il est indispensable que} $\rightarrow$ \esp{Es indispensable / Es necesario / Es fundamental que} + \textbf{Subjonctif présent} (volonté/nécessité). \emph{Gouvernement} = \esp{el gobierno}. \emph{Finance} = \esp{financie} (subj. \emph{financiar}). \emph{Davantage de} = \esp{más}. \emph{Centres d'hébergement} = \esp{centros de acogida / casas de acogida}. \emph{Femmes victimes de violences} = \esp{mujeres víctimas de violencia / maltrato}.
    \item \textbf{Proposition :} \esp{Es indispensable que el gobierno \textbf{financie} más \textbf{casas de acogida} para las mujeres \textbf{víctimas de violencia}.}
    \item \textit{Attention :} \emph{Víctimas de violencia} (singulier générique) ou \emph{violencias}. \emph{Casas de acogida} (terme institutionnel standard). Subjonctif \textbf{obligatoire} après \emph{Es indispensable que}.
\end{itemize}

\textbf{b)} \emph{« Les féministes luttent depuis des décennies pour que la société change ses mentalités machistes. »}
\begin{itemize}
    \item \textbf{Analyse :} \emph{Luttent depuis des décennies} = \emph{Llevan décadas luchando} (perifrase durative \emph{Llevar + gérondif}) OU \emph{Luchan desde hace décadas}. \emph{Pour que} = \esp{para que} + \textbf{Subjonctif}. \emph{Change} = \esp{cambie} (subj. \emph{cambiar}). \emph{Mentalités machistes} = \esp{mentalidades machistas}.
    \item \textbf{Proposition :} \esp{Las feministas \textbf{llevan décadas luchando} \textbf{para que} la sociedad \textbf{cambie} sus mentalidades machistas.}
    \item \textit{Alternative :} \esp{Luchan desde hace décadas para que...} (Moins idiomatique pour "depuis + durée" avec action continue). \emph{Para que} + subjonctif \textbf{obligatoire}.
\end{itemize}

\textbf{c)} \emph{« Cette loi a été votée pour protéger les victimes, mais elle n'est pas appliquée correctement par les juges. »}
\begin{itemize}
    \item \textbf{Analyse :} Passif composé \emph{a été votée} $\rightarrow$ \esp{ha sido votada / fue votada} (Passif sérén \emph{Ser + Participe} car agent implicite ou focus sur résultat). \emph{Pour protéger} = \esp{para proteger} (But = \emph{Para} + Infinitif). \emph{Mais} = \esp{pero / sin embargo}. \emph{N'est pas appliquée} = \esp{no es aplicada} (Passif sérén \emph{Ser + Participe} car agent explicite « par les juges »). \emph{Par les juges} = \esp{por los jueces} (Agent passif = \emph{Por}).
    \item \textbf{Proposition :} \esp{Esta ley \textbf{fue votada} \textbf{para proteger} a las víctimas, \textbf{pero no es aplicada} correctamente \textbf{por los jueces}.}
    \item \textit{Note :} \emph{Fue votada} (Prétérit indéfini : action passée achevée). \emph{No es aplicada} (Présent : vérité générale actuelle, passif sérén avec agent). \emph{Por los jueces} (Agent).
\end{itemize}

\textbf{d)} \emph{« Après avoir dénoncé son agresseur, elle a dû quitter son domicile par peur des représailles. »}
\begin{itemize}
    \item \textbf{Analyse :} \emph{Après avoir dénoncé} = \esp{Después de haber denunciado} (Gérondif composé / Infinitif composé prépositionnel). \emph{Son agresseur} = \esp{a su agresor / a su maltratador}. \emph{A dû} = \esp{tuvo que} (Obligation passée ponctuelle = Prétérit \emph{Tener que}) OU \emph{debió} (Devoir moral). \emph{Quitter} = \esp{dejar / abandonar}. \emph{Son domicile} = \esp{su domicilio / su casa}. \emph{Par peur de} = \esp{por miedo a} (Cause = \emph{Por}). \emph{Représailles} = \esp{represalias}.
    \item \textbf{Proposition :} \esp{\textbf{Después de haber denunciado} a su agresor, \textbf{tuvo que abandonar} su domicilio \textbf{por miedo a} represalias.}
    \item \textit{Grammaire :} \emph{Después de + Infinitif composé} (antériorité). \emph{Tuvo que} (Prétérit indéfini). \emph{Por miedo a} (Cause).
\end{itemize}

\textbf{2. Version}

\emph{« Malgré les campagnes de sensibilisation, le nombre de plaintes pour harcèlement sexuel dans l'environnement professionnel ne cesse de croître. Beaucoup de femmes se taisent car elles craignent de perdre leur emploi ou d'être stigmatisées par leurs collègues. Il est fondamental que les entreprises mettent en place des protocoles d'action efficaces et que l'inspection du travail sanctionne sévèrement les contrevenants. Ce n'est qu'ainsi que se brisera le cercle de l'impunité. »}

\textbf{Justifications grammaticales / lexicales :}
\begin{itemize}
    \item \esp{A pesar de} + Nom (concession) = « Malgré ».
    \item \esp{No deja de crecer} (Périphrase aspectuelle : action continue/ininterrompue) = « Ne cesse de croître ».
    \item \esp{Callan} (Verbe \emph{callar} = se taire, ne pas parler) = « Se taisent ». \textit{Variante courante :} \emph{se callan}.
    \item \esp{Temen perder / ser estigmatizadas} (Verbe \emph{temer} + Infinitif ; Passif réfléchi \emph{ser estigmatizadas} après préposition) = « Craignent de perdre / d'être stigmatisées ».
    \item \esp{Es fundamental que... implementen... sancione} (\emph{Es fundamental que} + \textbf{Subjonctif présent} pour les deux verbes : \emph{implementen} de \emph{implementar}, \emph{sancione} de \emph{sancionar}).
    \item \esp{Solo así se romperá} (Inversion sujet-verbe après \emph{Solo así} ; Futur simple \emph{romperá} = prédiction/certitude) = « Ce n'est qu'ainsi que se brisera ».
\end{itemize}
\end{exoresolu}

\courssub{Méthode : Expression orale / Débat : Argumenter sur l'égalité et la violence de genre}

\begin{methode}[Réussir l'épreuve orale (Exposé + Interaction) : Structure, Langage, Attitude]
\textbf{Objectif :} Présenter un point de vue argumenté (2--3 min) puis interagir avec l'examinateur (3--4 min) sur le thème de la famille / condition féminine.
\begin{enumerate}
    \item \textbf{Préparation (10 min) :} Notez des \cle{mots-clés} (pas de phrases entières). Structure : \emph{Introducción (Tesis) -- 2 Arguments + Exemples -- Conclusión + Apertura}.
    \item \textbf{Exposé (Monologue suivi) :}
        \begin{itemize}
            \item \textbf{Accroche :} \esp{Buenos días. Voy a hablar sobre... / El tema que voy a tratar es...}
            \item \textbf{Thèse claire :} \esp{Mi tesis es que la educación es la única herramienta para erradicar el machismo.}
            \item \textbf{Arguments structurés :} \esp{En primer lugar...; Además...; Por otro lado...; Un ejemplo claro es...; Los datos demuestran que...}
            \item \textbf{Conclusion + Ouverture :} \esp{En conclusión...; Para terminar, quisiera añadir que en Camerún...}
        \end{itemize}
    \item \textbf{Interaction (Dialogue) :} L'examinateur relance, contredit, demande des précisions.
        \begin{itemize}
            \item \textbf{Écoute active :} Ne pas couper. Reformuler si besoin : \esp{¿Quiere decir que...? / Si entiendo bien...}
            \item \textbf{Arguments de secours :} Préparer 2--3 idées "de réserve" (ex: rôle des hommes, éducation précoce, lois camerounaises, réseaux sociaux).
            \item \textbf{Langage de l'opinion / nuance :} \esp{En mi opinión...; Me parece que...; No estoy de acuerdo porque...; Es cierto, pero...; Hay que matizar...; Depende del contexto...}
        \end{itemize}
    \item \textbf{Prononciation / Prosodie :} Accent tonique respecté (\esp{mújer, violéncia, igualdád, machísmo}). Intonation montante (question), descendante (affirmation). Liaisons naturelles.
    \item \textbf{Critères d'évaluation (Grille Bac) :} Cohérence / Pertinence (5 pts) -- Correction grammaticale (5 pts) -- Richesse lexicale (5 pts) -- Prononciation / Fluidité (5 pts).
\end{enumerate}
\end{methode}

\begin{exoresolu}[Simulation d'épreuve orale : Exposé + Interaction]
\textbf{Énoncé :} \emph{Sujet : « La ley sola no basta para proteger a las mujeres : hace falta una revolución cultural ». Préparez un exposé de 2 minutes, puis simulez l'interaction avec l'examinateur (3 questions/réponses).}

\tcblower
\textbf{Solution détaillée :}

\textbf{PHASE 1 : EXPOSÉ (Modèle rédigé -- à apprendre par mots-clés)}

\emph{Buenos días. El tema que voy a tratar es la insuficiencia de la legislación frente a la violencia de género. Mi tesis es la siguiente: \textbf{la ley es el esqueleto, pero la educación y la cultura son la sangre que da vida a la protección real de las mujeres}.}

\emph{\textbf{En primer lugar}, las estadísticas son implacables. En España, desde la ley de 2004, los feminicidios no han desaparecido; en América Latina, las tasas son las más altas del mundo. \textbf{Por tanto}, el papel no detiene la mano del agresor si la impunidad judicial persiste y si los recursos (juzgados, casas de acogida) \textbf{faltan}.}

\emph{\textbf{En segundo lugar}, el machismo no es solo una conducta individual, es un sistema cultural. Se aprende en la familia, en la escuela, en los medios. \textbf{Por eso}, una \textbf{revolución cultural} implica una \textbf{coeducación obligatoria} desde la infancia que enseñe el consentimiento, la igualdad y la resolución pacífica de conflictos. \textbf{Un ejemplo}: en Suecia, la educación de género transversal ha reducido drásticamente la violencia de pareja.}

\emph{\textbf{En conclusión}, la ley fija el marco, pero la sociedad debe llenarlo de contenido. \textbf{Para terminar}, conecto esto con nuestra realidad en Camerún: tenemos leyes contra \textbf{la violencia de género}, pero el peso de las tradiciones patriarcales y el silencio social frenan su aplicación. \textbf{La lucha es global: ni una menos, vivas nos queremos. Gracias.}}

\textbf{PHASE 2 : INTERACTION (Simulation)}

\begin{center}
\begin{tabularx}{\linewidth}{|X|X|}\hline
\rowcolor{popblueL} \textbf{Examinateur (Question / Relance)} & \textbf{Candidat (Réponse argumentée)} \\ \hline
\textbf{Q1 :} \esp{« Usted dice que la educación es clave, pero ¿no cree que las penas más duras (cárcel perpetua, cadena perpetua) disuadirían más a los agresores? »} & \textbf{R1 :} \esp{Entiendo su punto, \textbf{pero matizo}. Las penas duras son necesarias para la justicia reparadora, pero no previenen el crimen \emph{in situ}. El agresor actúa en la intimidad, creyendo que no será pillado. La prevención real viene de cambiar la mentalidad \emph{antes} del acto. Además, en países con pena de muerte o cadena perpetua, la violencia no ha bajado. \textbf{La certeza del castigo (justicia rápida) disuade más que la dureza de la pena.}} \\ \hline
\textbf{Q2 :} \esp{« En el contexto camerunés, ¿cómo se puede conciliar la lucha contra la violencia de género con el respeto a las tradiciones y a la familia extensa? »} & \textbf{R2 :} \esp{Es un debate delicado. \textbf{No se trata de destruir la cultura, sino de eliminar las prácticas dañinas} (mutilación genital, matrimonio forzado, repudio). La familia extensa puede ser un \textbf{refugio} para la víctima si se sensibiliza a los líderes tradicionales (jefes, \emph{fons}, notables). En Camerún, la \emph{Ley de 2016 sobre la protección de la mujer} avanza, pero falta aplicación. \textbf{La clave: formar a los líderes comunitarios como aliados, no como enemigos.}} \\ \hline
\textbf{Q3 :} \esp{« ¿Qué papel deben jugar los hombres en esta « revolución cultural »? ¿No es un problema solo de mujeres? »} & \textbf{R3 :} \esp{\textbf{¡Error!} No es un problema de mujeres, es un problema de hombres que ejercen violencia. Los hombres deben ser \textbf{agentes de cambio}: denunciar a sus pares, educar a sus hijos en la igualdad, compartir las tareas domésticas y el cuidado (licencia de paternidad real). Movimientos como \emph{Hombres por la Igualdad} en España lo demuestran. \textbf{La masculinidad tóxica también encarcela al hombre.} La igualdad libera a ambos.} \\ \hline
\end{tabularx}
\end{center}

\textbf{Analyse de la performance :} Vocabulaire précis (\emph{insuficiencia, implacable, disuadir, matizo, prácticas dañinas, masculinidad tóxica}). Structures complexes (\emph{No se trata de... sino de...; La certeza del castigo...; Movimientos como...}). Interaction naturelle (reformulation, concession \emph{Entiendo su punto, pero...}, engagement personnel). Prononciation soignée des termes techniques.
\end{exoresolu}

\courssub{Méthode : Synthèse grammaticale ciblée : Subjonctif, Ser/Estar, Por/Para dans le contexte thématique}

\begin{methode}[Les 3 piliers grammaticaux du thème « Femmes, Famille, Violences » : Maîtriser les choix obligatoires]
\textbf{Objectif :} Automatiser les bons choix grammaticaux dans les productions écrites/orales (Thème, Commentaire, Expression).
\begin{enumerate}
    \item \textbf{LE SUBJONCTIF : Le mode de la subjectivité, de l'influence, de l'irréel.}
        \begin{itemize}
            \item \cle{Déclencheurs obligatoires (WEIRDO) :} \textbf{W}ish/Will (\emph{querer que, exigir que, pedir que, necesitar que, insistir en que}), \textbf{E}motion (\emph{sentir que, temer que, alegrarse de que, indignarse de que}), \textbf{I}mpersonal expressions (\emph{es necesario que, es urgente que, es fundamental que, basta que, hace falta que}), \textbf{R}ecommendation/Doubt (\emph{recomendar que, dudar que, no creer que}), \textbf{D}ependent clauses (\emph{para que, sin que, antes de que, a menos que, con tal de que}), \textbf{O}jalá.
            \item \textbf{Contexte leçon :} \emph{Exijo que el agresor \textbf{sea} juzgado. Temo que ella \textbf{no denuncie}. Es urgente que el Estado \textbf{financie} refugios. Luchamos para que las niñas \textbf{crezcan} libres.}
            \item \textbf{Attention :} \emph{Creer que / Pensar que / Estar seguro que} + \textbf{Indicatif} (certitude). \emph{No creer que / Dudar que} + \textbf{Subjonctif}.
        \end{itemize}
    \item \textbf{SER vs ESTAR : L'identité vs L'état/La localisation.}
        \begin{itemize}
            \item \cle{SER (Identité, Essence, Classification, Origine, Heure, Matériau, Possession, Nationalité, Prix) :} \esp{La violencia \textbf{es} un delito. Ella \textbf{es} víctima. El machismo \textbf{es} cultural. La casa \textbf{es} de la víctima. Son las tres.}
            \item \cle{ESTAR (Localisation, État resultant, État temporaire, Progressif, Actionnel passif "être en train de") :} \esp{La víctima \textbf{está} en un refugio. Ella \textbf{está} traumatizada (estado). La ley \textbf{está} siendo aplicada (pasiva de estado/proceso). \textbf{Está} prohibido maltratar.}
            \item \textbf{Piège majeur (Adjectifs changeant de sens) :}
                \begin{center}
                \begin{tabular}{|l|l|l|}\hline
                \rowcolor{popblueL} \textbf{Adjectif} & \textbf{SER} & \textbf{ESTAR} \\ \hline
                \esp{maltratado} & \emph{Es maltratado} (Passif actionnel : qqn le maltraite) & \emph{Está maltratado} (État : il porte des marques) \\ \hline
                \esp{protegido} & \emph{Es protegido} (Statut légal : personne protégée) & \emph{Está protegido} (Situation : en sécurité maintenant) \\ \hline
                \esp{violento} & \emph{Es violento} (Caractère, personnalité) & \emph{Está violento} (État momentané : en colère) \\ \hline
                \end{tabular}
                \end{center}
        \end{itemize}
    \item \textbf{POR vs PARA : Cause/Moyen vs But/Destination.}
        \begin{itemize}
            \item \cle{POR (Cause, Motif, Moyen, Durée, Échange, Agent passif, Mouvement vague, "Par erreur") :}
                \esp{Lucha \textbf{por} la igualdad (causa). Denunció \textbf{por} teléfono (medio). \textbf{Por} miedo (causa). Fue asesinada \textbf{por} su pareja (agente).}
            \item \cle{PARA (But, Finalité, Destinataire, Échéance, Opinion, Direction) :}
                \esp{La ley es \textbf{para} proteger (fin). Trabaja \textbf{para} una ONG (destinataire). Es \textbf{para} mañana (plazo). \textbf{Para} mí, es injusto (opinión). Salimos \textbf{para} Madrid (dirección).}
            \item \textbf{Astuce mnémotechnique :} \emph{POR} = \textbf{P}assé / \textbf{P}rovenance / \textbf{P}ar \textbf{P}... (Cause/Agent). \emph{PARA} = \textbf{P}rojet / \textbf{P}lus tard / \textbf{P}our \textbf{Q}ui (But/Destinataire).
        \end{itemize}
\end{enumerate}
\textbf{Entraînement :} Transformer 10 phrases français $\rightarrow$ espagnol en justifiant \emph{chaque} choix (Subj/Ind, Ser/Est, Por/Para).
\end{methode}

\begin{exoresolu}[Application grammaticale : Transformation et Justification]
\textbf{Énoncé :} Transformez les phrases suivantes selon les consignes. Justifiez chaque choix grammatical (Mode, Verbe, Préposition) en français.

\begin{enumerate}
    \item \textbf{Subjonctif / Indicatif :} Complétez avec le mode et le temps convenables (présent ou passé).
    \begin{enumerate}
        \item Es fundamental que el juez \_\_\_\_\_\_\_\_\_\_ (dictar) la orden de protección hoy mismo.
        \item La fiscal cree que la víctima \_\_\_\_\_\_\_\_\_\_ (decir) la verdad.
        \item No creo que el agresor \_\_\_\_\_\_\_\_\_\_ (cumplir) la condena íntegra.
        \item Dudo que haya \_\_\_\_\_\_\_\_\_\_ (haber) testigos presenciales.
    \end{enumerate}
    \item \textbf{Ser / Estar :} Choisissez et conjuguez au temps indiqué.
    \begin{enumerate}
        \item La violencia machista \_\_\_\_\_\_\_\_\_\_ (ser/estar - présent) una violación de derechos humanos.
        \item La mujer \_\_\_\_\_\_\_\_\_\_ (ser/estar - présent) en la casa de acogida desde ayer.
        \item El agresor \_\_\_\_\_\_\_\_\_\_ (ser/estar - présent) un hombre violento por naturaleza.
        \item El caso \_\_\_\_\_\_\_\_\_\_ (ser/estar - présent) siendo investigado por la policía.
    \end{enumerate}
    \item \textbf{Por / Para :} Complétez et justifiez.
    \begin{enumerate}
        \item Presentó la denuncia \_\_\_\_\_\_\_\_\_\_ proteger a sus hijos.
        \item Fue condenada \_\_\_\_\_\_\_\_\_\_ un tribunal especializado.
        \item Llevan años luchando \_\_\_\_\_\_\_\_\_\_ la igualdad.
        \item Esta campaña es \_\_\_\_\_\_\_\_\_\_ concienciar a la población.
    \end{enumerate}
\end{enumerate}

\tcblower
\textbf{Solution détaillée :}

\textbf{1. Subjonctif / Indicatif}
\begin{enumerate}
    \item \esp{dicté} (Subjonctif présent). \textbf{Justification :} \emph{Es fundamental que} (Expression impersonnelle de nécessité) $\rightarrow$ Subjonctif. Sujet différent (juez).
    \item \esp{dice} (Indicatif présent). \textbf{Justification :} \emph{Cree que} (Opinion/Certitude affirmative) $\rightarrow$ Indicatif.
    \item \esp{cumpla} (Subjonctif présent). \textbf{Justification :} \emph{No creo que} (Négation de l'opinion/Doute) $\rightarrow$ Subjonctif.
    \item \esp{haya} (Subjonctif présent). \textbf{Justification :} \emph{Dudo que} (Doute) $\rightarrow$ Subjonctif. \emph{Haber} (impersonnel) $\rightarrow$ \emph{haya}.
\end{enumerate}

\textbf{2. Ser / Estar}
\begin{enumerate}
    \item \esp{es} (Ser). \textbf{Justification :} Définition / Classification / Identité essentielle (\emph{una violación}). Caractère permanent, intrinsèque au concept.
    \item \esp{está} (Estar). \textbf{Justification :} Localisation physique (\emph{en la casa de acogida}). \emph{Desde ayer} marque la durée de l'état/localisation.
    \item \esp{es} (Ser). \textbf{Justification :} Trait de caractère / Personnalité (\emph{por naturaleza}). Identité définissante de l'individu.
    \item \esp{está} (Estar). \textbf{Justification :} Passif de procès / Action en cours (\emph{siendo investigado}). \emph{Estar + gérondif} = progressif ; \emph{Estar + participe} = passif d'état/processus (focus sur l'action en train de se faire). *Note : \emph{Es investigado} serait un passif d'état statique (statut), moins fréquent ici.*
\end{enumerate}

\textbf{3. Por / Para}
\begin{enumerate}
    \item \esp{para}. \textbf{Justification :} But / Finalité (\emph{proteger a sus hijos}). Répond à « ¿Para qué? ». Infinitif après préposition.
    \item \esp{por}. \textbf{Justification :} Agent du passif (\emph{Fue condenada por...}). Répond à « ¿Por quién? ».
    \item \esp{por}. \textbf{Justification :} Cause / Motif / Raison profonde (\emph{luchando por la igualdad}). Idée de « en faveur de », « à cause de », mouvement vague.
    \item \esp{para}. \textbf{Justification :} But / Destination de l'action (\emph{concienciar}). La campagne \emph{est pour} sensibiliser. Répond à « ¿Para qué? ».
\end{enumerate}

\textbf{Bilan :} La maîtrise de ces trois piliers (Subjonctif, Ser/Estar, Por/Para) représente 40\% des points de correction grammaticale au Bac. L'entraînement contextuel (thème violences/famille) ancre les automatismes.
\end{exoresolu}

\begin{attention}[Les 5 erreurs qui coûtent des points au Bac]
\begin{enumerate}
    \item \textbf{Accents oubliés ou mal placés (Tildes) :} \cle{¡Crucial !} \esp{violencia} (sans accent, paroxytonne en -a), \esp{denuncia} (nom, \textbf{sans accent}, paroxytonne en -a), \esp{dé} (verbe dar, pas *de), \esp{sé} (verbe saber, pas *se), \esp{tú} (pronom, pas *tu), \esp{él} (pronom, pas *el), \esp{más} (adv., pas *mas), \esp{está} (verbe, pas *esta), \esp{política} (accent sur i, esdrújula), \esp{análisis} (accent sur a, esdrújula). \textbf{Règle :} Mots oxytones (fin en voyelle, n, s) $\rightarrow$ accent sur dernière ; paroxytones (autres fins) $\rightarrow$ accent sur avant-dernière ; proparoxytones $\rightarrow$ toujours accent.
    \item \textbf{Faux amis « Transparents » :} \cle{Realizar} = Effectuer / Mener à bien (pas *Réaliser* au sens de comprendre). \cle{Asistir} = Assister à (aller à) / Aider (pas *Assister* au sens de regarder). \cle{Pretender} = Viser / Avoir l'intention / Revendiquer (pas *Prétendre* au sens d'affirmer faussement). \cle{Éxito} = Succès (pas *Exil*). \cle{Decepción} = Déception (pas *Déception* orthographe fr). \cle{Violación} = Viol (pas *Violation* sauf \emph{violación de derechos}). \cle{Maltrato} = Maltraitance (pas *Mauvais traitement* mot à mot).
    \item \textbf{Calques syntaxiques du français :}
        \begin{itemize}
            \item *\esp{La mujer es violada por su marido} $\rightarrow$ \textbf{Correct :} \esp{La mujer \textbf{fue} violada por su marido} (Passif sérén : action ponctuelle passée) OU \esp{La mujer \textbf{está siendo} violada} (en cours). \emph{Ser} pour action, \emph{Estar} pour état resultant.
            \item *\esp{Ella tiene 20 años} $\rightarrow$ \textbf{Correct.} Mais *\esp{Ella es 20 años} $\rightarrow$ \textbf{Erreur.}
            \item *\esp{Es importante para que...} $\rightarrow$ \textbf{Erreur.} \emph{Es importante QUE + Subj.} (Pas de \emph{para}).
            \item *\esp{Busco a alguien que me puede ayudar} $\rightarrow$ \textbf{Erreur.} Antécédent indéfini/négatif $\rightarrow$ \esp{Busco a alguien que me \textbf{pueda} ayudar} (Subjonctif).
        \end{itemize}
    \item \textbf{Accords oubliés (Genre / Nombre) :} \cle{La víctima} (F) $\rightarrow$ \emph{la víctima \textbf{herida}}, \emph{la víctima \textbf{denunciante}}. \cle{El agresor} (M) $\rightarrow$ \emph{el agresor \textbf{condenado}}. \cle{Las medidas} (F pl) $\rightarrow$ \emph{las medidas \textbf{adoptadas}}. \cle{El machismo} (M sg) $\rightarrow$ \emph{el machismo \textbf{estructural}}. \textbf{Piège :} \emph{La mano, el día, el mapa, el problema, el sistema, el tema, el idioma, el clima, el programa, el planeta} (Masculins en -\emph{a} / -\emph{ma} d'origine grecque).
    \item \textbf{Mauvais choix de temps / Aspect (Passé) :}
        \begin{itemize}
            \item \cle{Prétérit indéfini (Pretérito indefinido) :} Action passée, achevée, datée, ponctuelle. \emph{Ayer \textbf{denunció}.}
            \item \cle{Imparfait (Pretérito imperfecto) :} Description, habitude, action en cours, simultanéité. \emph{Mientras \textbf{esperaba}, \textbf{lloraba}.}
            \item \cle{Passé composé (Pretérito perfecto compuesto) :} Passé récent, lien avec le présent, « cette semaine/mois/année ». \emph{Este año \textbf{ha habido} muchos feminicidios.}
            \item \textbf{Erreur type :} *\emph{Ayer he ido} (Français *\emph{J'ai été hier}*). \textbf{Correct :} \emph{Ayer \textbf{fui}}.
        \end{itemize}
\end{enumerate}
\textbf{Conseil final :} Relisez \textbf{à l'envers} (dernière phrase $\rightarrow$ première) pour chasser les fautes d'inattention (accords, accents, \emph{ser/estar}, \emph{por/para}).
\end{attention}

\newpage

\coursec{Exercices}

\coursec{Lectura y comentario : Discriminaciones y violencias contra las mujeres}

\courssub{Texte support : lecture et découverte}

\begin{textebox}[La violencia de género: una lacra persistente]
La violencia de género no entiende de fronteras, ni de clases sociales, ni de nivel educativo. Afecta a mujeres de todas las edades y condiciones, aunque las más vulnerables —mujeres migrantes, con discapacidad, en situación de pobreza— sufren una doble discriminación. En Camerún, como en España o en Guinea Ecuatorial, el machismo sigue impregnando muchas relaciones familiares y laborales. Según datos del Observatorio contra la Violencia Doméstica y de Género, solo en España se registraron más de 30.000 denuncias por maltrato en 2023. Las víctimas sufren no solo agresiones físicas, sino también acoso psicológico, control económico y aislamiento social. La ley orgánica 1/2004 de medidas de protección integral fue un avance histórico, pero su aplicación sigue siendo desigual según las comunidades autónomas y, en muchos países, los recursos para la protección de las víctimas son insuficientes. \emph{Es urgente que la sociedad toda se movilice para romper el ciclo de la violencia.} La educación en igualdad desde la infancia es la mejor vacuna contra el machismo.
\end{textebox}

\begin{definition}[Glosario del texto]
\begin{itemize}
\item \textbf{La lacra} : le fléau, la plaie sociale.
\item \textbf{La doble discriminación} : la double discrimination (genre + origine/handicap/pauvreté).
\item \textbf{Impregnar} : imprégner, pénétrer profondément.
\item \textbf{El maltrato} : la maltraitance, les mauvais traitements.
\item \textbf{El aislamiento social} : l'isolement social.
\item \textbf{La protección integral} : la protection globale / intégrale (juridique, psychologique, sociale).
\item \textbf{Romper el ciclo} : briser le cycle.
\end{itemize}
\end{definition}

\courssub{Compréhension du texte}

\begin{exercice}[Questions de compréhension globale et détaillée]
Répondez en espagnol par des phrases complètes.
\begin{enumerate}
\item ¿Cuál es el tema principal del texto?
\item Según el texto, ¿qué grupos de mujeres son los más vulnerables?
\item ¿Qué cifra clave se da para España en 2023?
\item Enumere cuatro formas de violencia mencionadas en el texto además de la agresión física.
\item ¿Qué valoración se hace de la ley orgánica 1/2004?
\item ¿Cuál es la propuesta final del autor para erradicar el machismo?
\end{enumerate}
\end{exercice}

\courssub{Lexique thématique : La familia y la condición de la mujer}

\begin{definition}[Vocabulaire essentiel : La violencia de género y la familia]
\begin{center}
\begin{tabularx}{\linewidth}{|X|X|X|}
\hline
\rowcolor{popblueL} \textbf{Español} & \textbf{Français} & \textbf{Contexto / Ejemplo} \\
\hline
\emph{la violencia de género} & la violence de genre (basée sur le genre) & \esp{La violencia de género viola los derechos humanos.} \\
\hline
\emph{el machismo} & le machisme (idéologie de supériorité masculine) & \esp{El machismo justifica el control sobre la mujer.} \\
\hline
\emph{el maltrato} & la maltraitance, les violences & \esp{Denunció a su pareja por maltrato habitual.} \\
\hline
\emph{la denuncia} & la plainte, la dénonciation & \esp{Presentar una denuncia es el primer paso.} \\
\hline
\emph{el acoso (psicológico / sexual / laboral)} & le harcèlement (psychologique / sexuel / au travail) & \esp{El acoso psicológico deja secuelas invisibles.} \\
\hline
\emph{la víctima} & la victime & \esp{Las víctimas necesitan protección integral.} \\
\hline
\emph{la igualdad (real / efectiva)} & l'égalité (réelle / effective) & \esp{Luchar por la igualdad real es tarea de todos.} \\
\hline
\emph{la protección (integral / judicial / policial)} & la protection (globale / judiciaire / policière) & \esp{La orden de alejamiento es una medida de protección.} \\
\hline
\emph{la orden de alejamiento} & l'ordonnance d'éloignement & \esp{El juez dictó una orden de alejamiento de 500 metros.} \\
\hline
\emph{el juzgado de violencia sobre la mujer} & le tribunal des violences faites aux femmes & \esp{El juzgado especializado tramita la causa con rapidez.} \\
\hline
\emph{el techo de cristal} & le plafond de verre & \esp{Las mujeres chocan contra el techo de cristal en las empresas.} \\
\hline
\emph{la brecha salarial} & l'écart salarial & \esp{La brecha salarial en el sector tecnológico es del 18,5\%.} \\
\hline
\emph{la corresponsabilidad} & la coresponsabilité (partage des tâches) & \esp{La corresponsabilidad doméstica libera a la mujer.} \\
\hline
\emph{la doble discriminación} & la double discrimination & \esp{Sufre doble discriminación: por mujer y por migrante.} \\
\hline
\end{tabularx}
\end{center}
\end{definition}

\begin{attention}[Faux amis et pièges fréquents]
\begin{itemize}
\item \textbf{Violencia de género} $\neq$ \emph{violencia doméstica} (cette dernière inclut violences entre tous membres du foyer ; la première vise spécifiquement la femme \emph{por el hecho de ser mujer}).
\item \textbf{Machismo} : nom masculin (\emph{el machismo}), pas \emph{*la machisma}. Dérivé : \emph{macho} (adj./n.m.), \emph{machista} (adj./n. m/f).
\item \textbf{Denuncia} : nom féminin (\emph{la denuncia}). Verbe : \emph{denunciar}. Ne pas confondre avec \emph{denuncio} (1ère pers. sg. présent).
\item \textbf{Acoso} : nom masculin. \emph{Acoso escolar} (harcèlement scolaire), \emph{acoso laboral} (harcèlement moral au travail), \emph{acoso sexual}.
\item \textbf{Orden de alejamiento} : \emph{la orden} (fém.), mais \emph{el alejamiento} (masc.). Accord : \emph{órdenes de alejamiento}.
\item \textbf{Techo de cristal} : expression figée. Ne pas dire \emph{*techo de vidrio}.
\item \textbf{Brecha} : écart, fossé. \emph{Brecha digital, brecha de género, brecha salarial}.
\end{itemize}
\end{attention}

\courssub{Grammaire : Voix passive réfléchie et impersonnelle \emph{se}}

\begin{propriete}[Passif réfléchi (\emph{se} + verbe) vs Impersonnel \emph{se}]
\textbf{1. Passif réfléchi (Passiva refleja)} : Sujet patient + \emph{se} + verbe accordé au sujet. Équivaut au passif français (\emph{être} + participe). Utilisé quand l'agent n'est pas exprimé ou secondaire.
\begin{itemize}
\item \esp{Se venden casas.} (On vend des maisons / Des maisons sont vendues.)
\item \esp{Se detuvo al agresor.} (L'agresseur a été arrêté.)
\item \esp{Se dictaron órdenes de alejamiento.} (Des ordonnances ont été émises.)
\end{itemize}
\textbf{Règle d'accord} : Verbe au singulier si sujet singulier / infinitif ; pluriel si sujet pluriel.
\esp{Se necesita voluntarios.} (Sujet \emph{voluntarios} pluriel $\rightarrow$ verbe pluriel).

\textbf{2. Impersonnel \emph{se} (Impersonal refleja)} : Pas de sujet patient. \emph{Se} + verbe à la 3ème personne du singulier. Sujet indéterminé (\emph{on}, \emph{les gens}). Souvent avec verbes intransitifs ou transitifs sans objet patient explicite.
\begin{itemize}
\item \esp{Se vive bien aquí.} (On vit bien ici.)
\item \esp{En las reuniones se ignora a las mujeres.} (Dans les réunions, on ignore les femmes / les femmes sont ignorées.)
\item \esp{Se trabaja mucho en esta empresa.} (On travaille beaucoup dans cette entreprise.)
\end{itemize}

\textbf{Comment choisir ?} Transformez en passif classique (\emph{ser} + participe). Si possible $\rightarrow$ Passif réfléchi. Si impossible (verbe intransitif, ou COD personne introduit par \emph{a} qu'on ne veut pas mettre en sujet) $\rightarrow$ Impersonnel \emph{se}.
\end{propriete}

\begin{exemplebox}[Exemples contrastés]
\begin{itemize}
\item \esp{La policía detuvo al agresor.} $\rightarrow$ \esp{Se detuvo al agresor.} (Passif réfléchi : \emph{al agresor} reste COD, verbe au singulier car pas de sujet patient explicite en position sujet, ou accord avec \emph{agresor} si mis en sujet : \emph{El agresor fue detenido} / \emph{Se detuvo al agresor}).
\item \emph{Meilleur test} : \esp{Los jueces dictan órdenes.} $\rightarrow$ \esp{Se dictan órdenes.} (Sujet \emph{órdenes} pluriel $\rightarrow$ verbe pluriel = Passif réfléchi).
\item \esp{En la empresa se trabaja duro.} (Pas de sujet patient $\rightarrow$ Impersonnel).
\item \esp{A las víctimas se les atiende bien.} (COD personne \emph{a las víctimas} maintenu $\rightarrow$ Impersonnel \emph{se} + redondance \emph{les}).
\end{itemize}
\end{exemplebox}

\begin{exercice}[Transformation : Passif réfléchi ou Impersonnel \emph{se} ?]
Transformez les phrases actives selon la structure indiquée. Attention à l'accord du verbe.
\begin{enumerate}
\item La policía \emph{detuvo} al agresor anoche. $\rightarrow$ (Passif réfléchi)
\item Los jueces \emph{dictan} órdenes de alejamiento cada día. $\rightarrow$ (Passif réfléchi)
\item La sociedad \emph{condena} la violencia de género. $\rightarrow$ (Passif réfléchi)
\item Los medios de comunicación \emph{difunden} campañas de prevención. $\rightarrow$ (Passif réfléchi)
\item El gobierno \emph{aprobó} una nueva ley de protección. $\rightarrow$ (Passif réfléchi)
\item Los psicólogos \emph{atienden} a las víctimas en los centros. $\rightarrow$ (Impersonnel \emph{se} + \emph{les})
\item En la escuela \emph{educan} a los niños en igualdad. $\rightarrow$ (Impersonnel \emph{se})
\item Nadie \emph{olvida} el dolor de las víctimas. $\rightarrow$ (Impersonnel \emph{se})
\end{enumerate}
\end{exercice}

\courssub{Grammaire : Subjonctif présent : valeur de nécessité, doute, souhait}

\begin{propriete}[Emplois du subjonctif présent (révision ciblée)]
Le subjonctif s'emploie dans la subordonnée introduite par \emph{que} quand le principal exprime :
\begin{enumerate}
\item \textbf{Volonté, ordre, conseil, nécessité} : \esp{Es necesario que...}, \esp{Es urgente que...}, \esp{Queremos que...}, \esp{Exijo que...}, \esp{Recomiendo que...} $\rightarrow$ \textbf{Subjonctif}.
\item \textbf{Doute, négation de certitude} : \esp{No creo que...}, \esp{Dudo que...}, \esp{No es seguro que...} $\rightarrow$ \textbf{Subjonctif}. (Affirmation \emph{Creo que...} $\rightarrow$ Indicatif).
\item \textbf{Souhait, espoir} : \esp{Ojalá que...}, \esp{Espero que...}, \esp{Quiero que...} $\rightarrow$ \textbf{Subjonctif}.
\item \textbf{Émotion, valeur subjective} : \esp{Me duele que...}, \esp{Es injusto que...}, \esp{Me parece bien que...} $\rightarrow$ \textbf{Subjonctif}.
\item \textbf{Conjonctions de finalité, condition, temps (futur)} : \emph{para que, a fin de que, sin que, antes de que, cuando (futur), hasta que (futur)} $\rightarrow$ \textbf{Subjonctif}.
\end{enumerate}
\textbf{Formation} : Radical de la 1ère pers. sg. présent indicatif (\emph{yo}) + terminaisons opposées (\emph{-ar} $\rightarrow$ \emph{-e, -es, -e, -emos, -éis, -en} ; \emph{-er/-ir} $\rightarrow$ \emph{-a, -as, -a, -amos, -áis, -an}).
\end{propriete}

\begin{attention}[Irréguliers fréquents au subjonctif présent]
\begin{center}
\begin{tabular}{|l|l|l|l|}
\hline
\rowcolor{popblueL} \textbf{Infinitif} & \textbf{Yo (Ind. Prés.)} & \textbf{Radical Subj.} & \textbf{Formes (yo, tú, él, ns, vs, ell)} \\
\hline
\emph{Ser} & soy & se- & sea, seas, sea, seamos, seáis, sean \\
\hline
\emph{Estar} & estoy & est- & esté, estés, esté, estemos, estéis, estén \\
\hline
\emph{Tener} & tengo & teng- & tenga, tengas, tenga, tengamos, tengáis, tengan \\
\hline
\emph{Hacer} & hago & hag- & haga, hagas, haga, hagamos, hagáis, hagan \\
\hline
\emph{Poner} & pongo & pong- & ponga, pongas, ponga, pongamos, pongáis, pongan \\
\hline
\emph{Saber} & sé & sep- & sepa, sepas, sepa, sepamos, sepáis, sepan \\
\hline
\emph{Ir} & voy & vay- & vaya, vayas, vaya, vayamos, vayáis, vayan \\
\hline
\emph{Haber} & he & hay- & haya, hayas, haya, hayamos, hayáis, hayan \\
\hline
\end{tabular}
\end{center}
\end{attention}

\begin{exercice}[Conjugaison : Complétez avec le temps convenu]
Conjuguez les verbes entre parenthèses. Indiquez le mode et le temps utilisés.
\begin{enumerate}
\item Muchas mujeres \_\_\_\_\_\_\_\_ (sufrir) en silencio durante años antes de pedir ayuda. (Pluscuamperfecto de indicativo)
\item Es urgente que el gobierno \_\_\_\_\_\_\_\_ (tomar) medidas más eficaces contra el machismo. (Subjuntivo presente)
\item Si la sociedad \_\_\_\_\_\_\_\_ (educar) en igualdad, la violencia de género \_\_\_\_\_\_\_\_ (disminuir). (Imperfecto de subjuntivo / Condicional simple)
\item La ley \_\_\_\_\_\_\_\_ (entrar) en vigor en 2004 y desde entonces \_\_\_\_\_\_\_\_ (salvar) muchas vidas. (Pretérito indefinido / Pretérito perfecto compuesto)
\item No creo que el acoso psicológico \_\_\_\_\_\_\_\_ (ser) menos grave que el físico. (Subjuntivo presente)
\item Las víctimas \_\_\_\_\_\_\_\_ (necesitar) apoyo psicológico, jurídico y económico. (Presente de indicativo)
\item Cuando yo \_\_\_\_\_\_\_\_ (ser) joven, no \_\_\_\_\_\_\_\_ (haber) tanta conciencia sobre este problema. (Imperfecto indicativo / Imperfecto indicativo)
\item Ojalá que pronto \_\_\_\_\_\_\_\_ (lograr) una igualdad real entre hombres y mujeres. (Subjuntivo presente)
\end{enumerate}
\end{exercice}

\courssub{Connecteurs logiques : structurer l'argumentation}

\begin{definition}[Principaux connecteurs pour le commentaire]
\begin{center}
\begin{tabularx}{\linewidth}{|X|X|X|}
\hline
\rowcolor{popblueL} \textbf{Connecteur} & \textbf{Fonction} & \textbf{Exemple} \\
\hline
\emph{Por tanto, en consecuencia, por consiguiente} & Conséquence, conclusion & \esp{No hay recursos, \textbf{por tanto} las víctimas no denuncian.} \\
\hline
\emph{Sin embargo, no obstante, ahora bien} & Opposition, concession & \esp{La ley es buena, \textbf{sin embargo} su aplicación falla.} \\
\hline
\emph{Además, asimismo, también} & Addition, renchérissement & \esp{Hay acoso físico, \textbf{además} control económico.} \\
\hline
\emph{Por ejemplo, por poner un caso} & Illustration & \esp{Hay muchas leyes, \textbf{por ejemplo} la ley 1/2004.} \\
\hline
\emph{En definitiva, en resumen, en conclusión} & Synthèse, conclusion finale & \esp{\textbf{En definitiva}, la educación es la clave.} \\
\hline
\emph{En primer lugar, en segundo lugar, finalmente} & Organisation, énumération & \esp{\textbf{En primer lugar}, hay que legislar.} \\
\hline
\emph{Es decir, o sea, esto es} & Explication, reformulation & \esp{Es un delito, \textbf{es decir}, castigado por la ley.} \\
\hline
\end{tabularx}
\end{center}
\end{definition}

\begin{exercice}[Texte à trous : Conjugaison et connecteurs]
Complétez avec les verbes (temps convenu) et les connecteurs (\emph{por tanto, sin embargo, además, en consecuencia, por ejemplo, en definitiva}).
\begin{textebox}[La lucha contra la violencia de género]
La violencia de género \_\_\_\_\_\_\_\_ (ser) una violación de los derechos humanos que \_\_\_\_\_\_\_\_ (afectar) a millones de mujeres. \_\_\_\_\_\_\_\_, no solo \_\_\_\_\_\_\_\_ (manifestar) a través de golpes, \_\_\_\_\_\_\_\_ también mediante control económico y aislamiento. \_\_\_\_\_\_\_\_, muchas víctimas no \_\_\_\_\_\_\_\_ (denunciar) por miedo. \_\_\_\_\_\_\_\_, en España la ley \_\_\_\_\_\_\_\_ (garantizar) la protección, pero \_\_\_\_\_\_\_\_ (hacer) falta más recursos. \_\_\_\_\_\_\_\_, la educación en igualdad \_\_\_\_\_\_\_\_ (ser) la clave para erradicar el machismo.
\end{textebox}
\end{exercice}

\courssub{Expressions idiomatiques et locutions}

\begin{exercice}[Appariement : Expressions idiomatiques]
Reliez chaque expression (A) à son sens (B) et écrivez une phrase exemple en espagnol.
\begin{center}
\begin{tabularx}{\linewidth}{|X|X|}
\hline
\rowcolor{popblueL} \textbf{A. Expresión} & \textbf{B. Significado} \\
\hline
1. romper el silencio & a. Tomar medidas drásticas para resolver un problema de raíz. \\
\hline
2. poner el grito en el cielo & b. Denunciar públicamente una injusticia que se callaba. \\
\hline
3. pasar factura & c. Protestar enérgicamente ante algo que se considera injusto. \\
\hline
4. cortar de raíz & d. Tener consecuencias negativas a largo plazo. \\
\hline
5. alzar la voz & e. Expresar su opinión con valentía, no callar. \\
\hline
\end{tabularx}
\end{center}
\end{exercice}

\courssub{Méthode du commentaire de texte (Type Bac)}

\begin{methode}[Commentaire de texte : 3 étapes clés]
\textbf{Étape 1 : Présentation (Introducción)} \\
\emph{Contexte} : Auteur (si connu), source, date, type de document (article, témoignage, rapport). \\
\emph{Thème} : Sujet général (violencia de género, discriminación laboral, etc.). \\
\emph{Thèse / Problématique} : Idée directrice de l'auteur + votre question de départ. \\
\emph{Annonce du plan} : \esp{En primer lugar...; En segundo lugar...; Finalmente...}

\textbf{Étape 2 : Développement (Desarrollo / Análisis)} \\
2 à 3 paragraphes arguments. Chaque paragraphe = 1 idée force + citation courte + analyse (vocabulaire, figures de style, arguments de l'auteur) + lien avec connaissances personnelles (contexte Cameroun, Guinée équatoriale, actualité). \\
Utilisez connecteurs : \esp{Por un lado... por otro lado; Además; Sin embargo; En consecuencia.}

\textbf{Étape 3 : Conclusion (Conclusión)} \\
Synthèse des arguments. Réponse à la problématique. Ouverture (perspective, proposition d'action, citation). \\
\emph{Ne pas introduire de nouvelle idée.}
\end{methode}

\courssub{Je me prépare au Bac : Entraînement complet}

\begin{exercice}[Compréhension et analyse de texte inédit (Type Bac)]
\begin{textebox}[El machismo en el ámbito laboral: una barrera invisible]
María González, ingeniera informática de 34 años, trabaja en una empresa tecnológica de Madrid desde hace cinco años. A pesar de tener mejor expediente académico y más experiencia que varios de sus compañeros varones, lleva tres años estancada en el mismo puesto. «En las reuniones, mis propuestas se ignoran hasta que un hombre las repite. Mi jefe me pide que traiga café cuando vienen clientes, algo que nunca le pide a mis compañeros», denuncia. Su caso no es único: según el Instituto de la Mujer, el 67\% de las profesionales españolas afirma haber sufrido discriminación salarial o de ascenso por razón de género. La brecha salarial en el sector tecnológico alcanza el 18,5\%. «El techo de cristal no se ve, pero se siente», resume María. Las empresas que implementan planes de igualdad obligatorios desde 2021 notan mejoras, \textbf{pero es fundamental que la sociedad entera asuma la corresponsabilidad}. Los expertos insisten en que la corresponsabilidad y la educación no sexista son la única vía para lograr una igualdad real en el trabajo.
\end{textebox}

\textbf{Consignes :}
\begin{enumerate}
\item \textbf{Compréhension (réponses en espagnol, phrases complètes) :}
\begin{enumerate}
\item ¿Cuál es la profesión de María González y cuántos años lleva en la empresa?
\item ¿Qué trato desigual recibe María en las reuniones y ante los clientes?
\item ¿Qué porcentaje de profesionales españolas declara haber sufrido discriminación según el Instituto de la Mujer?
\item ¿A cuánto asciende la brecha salarial en el sector tecnológico?
\item ¿Qué medidas han mejorado la situación en algunas empresas desde 2021?
\end{enumerate}
\item \textbf{Vrai ou Faux ? Justifiez par une citation exacte du texte.}
\begin{enumerate}
\item María tiene menos experiencia que sus compañeros varones.
\item La brecha salarial en el sector tecnológico es inferior al 10\%.
\item La paternidad frena la carrera profesional de los hombres igual que la maternidad la de las mujeres.
\item Los planes de igualdad obligatorios han supuesto mejoras en las empresas que los aplican.
\end{enumerate}
\item \textbf{Analyse linguistique :}
\begin{enumerate}
\item Expliquez en français l'expression : «\emph{El techo de cristal no se ve, pero se siente}».
\item Expliquez : «\emph{La maternidad sigue siendo un freno para la carrera profesional}».
\end{enumerate}
\item \textbf{Relevé grammatical dans le texte :}
\begin{enumerate}
\item Relevez deux connecteurs logiques d'opposition.
\item Identifiez \textbf{un verbe au subjonctif présent}, citez la phrase complète et justifiez son emploi.
\item Relevez deux noms formés par dérivation préfixale (ex: \emph{des-igualdad}).
\end{enumerate}
\end{enumerate}
\end{exercice}

\begin{exercice}[Thème et Version (Type Bac)]

\textbf{A. Thème (Français $\rightarrow$ Espagnol)}
Traduisez le passage suivant en espagnol soigné.
« Au Cameroun comme en Espagne, la violence faite aux femmes reste un fléau silencieux. Trop de victimes n'osent pas porter plainte par peur des représailles ou par honte. Pourtant, la loi protège les femmes et prévoit des peines sévères pour les agresseurs. Il est urgent que la société toute entière — hommes et femmes — se mobilise pour briser le cycle de la violence. L'éducation à l'égalité dès l'enfance est le meilleur vaccin contre le machisme. »
\end{exercice}


\textbf{B. Version (Espagnol $\rightarrow$ Français)}
Traduisez en français le texte suivant.
\begin{textebox}
La violencia de género no entiende de fronteras, ni de clases sociales, ni de nivel educativo. Afecta a mujeres de todas las edades y condiciones, aunque las más vulnerables —mujeres migrantes, con discapacidad, en situación de pobreza— sufren una doble discriminación. Las instituciones deben garantizar no solo la persecución del delito, sino también la reparación integral de la víctima: asistencia jurídica gratuita, apoyo psicológico especializado, formación para el empleo y vivienda digna. Solo así se rompe el círculo de la violencia y se devuelve a la mujer su autonomía y su dignidad.
\end{textebox}

\begin{exercice}[Expression écrite : Commentaire personnel guidé (Type Bac)]
\textbf{Sujet :} \esp{¿Crees que la igualdad real entre hombres y mujeres es posible en tu país? Justifica tu respuesta con argumentos concretos y ejemplos.}

\textbf{Consignes de rédaction :}
\begin{itemize}
\item Longueur : 10 à 12 lignes (120-150 mots).
\item Structure : Introduction (accroche + thèse) / Développement (2 arguments + exemples) / Conclusion (ouverture).
\item Contraintes linguistiques obligatoires :
\begin{enumerate}
\item Minimum 5 mots du vocabulaire de la leçon (\emph{violencia de género, machismo, igualdad, denuncia, acoso, víctima, maltrato, protección, techo de cristal, brecha salarial, corresponsabilidad}).
\item Minimum 2 connecteurs logiques (\emph{por tanto, sin embargo, además, en consecuencia, en definitiva}).
\item Minimum 1 structure au subjonctif (\emph{es necesario que..., ojalá que..., no creo que..., es fundamental que...}).
\item Orthographe soignée (accents, \emph{ñ}, \emph{¿ ¡}).
\end{enumerate}
\end{itemize}

\end{exercice}

\begin{exemplebox}[Modèle de commentaire (Corrigé type)]
\esp{En mi opinión, la igualdad real es posible, pero exige un cambio profundo de mentalidades. En primer lugar, el machismo sigue siendo un obstáculo estructural: en Camerún, como en España, muchas mujeres sufren violencia de género y acoso laboral sin atreverse a presentar una denuncia por miedo o dependencia económica. Por tanto, es fundamental que el Estado refuerce la protección integral de las víctimas.}
\esp{Sin embargo, los avances legales existen: la ley contra la violencia de género y los planes de igualdad en las empresas empiezan a reducir la brecha salarial y a romper el techo de cristal. Además, la corresponsabilidad en el hogar permite a las mujeres desarrollarse profesionalmente.}
\esp{En definitiva, la educación no sexista desde la infancia es la clave. Ojalá que las nuevas generaciones construyan una sociedad donde ser mujer no sea un riesgo, sino un derecho pleno.}
\end{exemplebox}



\begin{aretenir}[Bilan de la leçon E01]
\begin{itemize}
\item \textbf{Lexique maîtrisé} : \emph{violencia de género, machismo, maltrato, denuncia, acoso, víctima, igualdad, protección, orden de alejamiento, techo de cristal, brecha salarial, corresponsabilidad, doble discriminación}.
\item \textbf{Grammaire} : Passif réfléchi (\emph{se} + verbe accordé) vs Impersonnel \emph{se} (verbe 3ème pers. sg.) ; Subjonctif présent après expressions de nécessité, doute, souhait, émotion.
\item \textbf{Méthode} : Commentaire structuré (Présentation - Analyse argumentée - Conclusion ouverte) + connecteurs logiques.
\item \textbf{Compétences Bac} : Compréhension fine (V/F + citations), Thème/Version (registre formel, vocabulaire précis), Expression écrite argumentée (120-150 mots, contraintes linguistiques).
\end{itemize}
\end{aretenir}

\begin{savaistu}[Cameroun et Guinée équatoriale : cadres juridiques]
\begin{itemize}
\item \textbf{Cameroun} : La Constitution de 1996 (préambule) garantit l'égalité. Loi n°2016/007 du 12 juil. 2016 (Code pénal) : criminalise les violences conjugales (art. 356-357), le harcèlement sexuel (art. 297-1). Existence de « \emph{Guichets genre} » dans certains tribunaux. Défis : application, coutumes discriminatoires (héritage, mariage précoce), impunité.
\item \textbf{Guinée équatoriale} : Seule nation hispanophone d'Afrique subsaharienne. Loi n°1/2011 sur l'égalité des chances et le traitement entre hommes et femmes. Ratification CEDAW. Problèmes persistants : violences domestiques sous-déclarées, faible représentation politique des femmes malgré quotas.
\item \textbf{Espagne} : Référence pionnière avec \emph{Ley Orgánica 1/2004} (protection intégrale), \emph{Ley Orgánica 3/2007} (égalité effective), \emph{Ley Orgánica 10/2022} (liberté sexuelle, « solo sí es sí »). \emph{Juzgados de Violencia sobre la Mujer} spécialisés. \emph{Observatorio contra la Violencia Doméstica y de Género} (statistiques officielles).
\end{itemize}
\end{savaistu}

\newpage

\coursec{Corrigés des exercices}

\begin{corrige}[Exercice 1 — Questions de compréhension globale et détaillée]
\textbf{Barème indicatif (6 pts)} : 1 point par réponse correcte, complète et rédigée en espagnol correct (0,5 pt si l'idée est juste mais la langue fautive).

\textbf{Réponses attendues :}
\begin{enumerate}
\item \esp{El tema principal es la violencia de género como una lacra persistente que afecta a mujeres de todas las condiciones sociales, y la necesidad de movilizar a toda la sociedad para erradicarla.} \\
\emph{Le sujet principal est la violence de genre comme fléau persistant qui touche les femmes de toutes conditions, et la nécessité de mobiliser toute la société pour l'éradiquer.}
\item \esp{Las mujeres más vulnerables son las migrantes, las mujeres con discapacidad y las que viven en situación de pobreza, porque sufren una doble discriminación.}
\item \esp{Según el texto, solo en España se registraron más de 30.000 denuncias por maltrato en 2023.}
\item \esp{Además de las agresiones físicas, el texto menciona el acoso psicológico, el control económico, el aislamiento social y la doble discriminación.}
\item \esp{La ley orgánica 1/2004 se considera un avance histórico, pero su aplicación sigue siendo desigual según las comunidades autónomas y los recursos son insuficientes.}
\item \esp{El autor propone la educación en igualdad desde la infancia como la mejor vacuna contra el machismo.}
\end{enumerate}

\textbf{Points de vigilance :}
\begin{itemize}
\item Rédiger des \textbf{phrases complètes} (sujet + verbe), jamais de réponse nominale : \esp{*Las migrantes.} est insuffisant.
\item Attention à l'accord : \esp{las mujeres más vulnerables} (fém. plur.), \esp{se registraron} (accord avec \esp{denuncias}).
\item Ne pas recopier le texte sans reformuler : le correcteur valorise la paraphrase.
\end{itemize}
\end{corrige}

\begin{corrige}[Exercice 2 — Transformation : passif réfléchi ou impersonnel \emph{se}]
\textbf{Barème indicatif (8 pts)} : 1 point par transformation correcte (0,5 pt si l'accord verbal est fautif).

\textbf{Transformations attendues :}
\begin{enumerate}
\item \esp{Se detuvo al agresor anoche.} \\
\emph{On a arrêté l'agresseur hier soir.} — COD de personne introduit par \esp{a} : le verbe reste au singulier (construction impersonnelle admise).
\item \esp{Se dictan órdenes de alejamiento cada día.} \\
\emph{Des ordonnances d'éloignement sont émises chaque jour.} — Sujet patient \esp{órdenes} pluriel $\rightarrow$ verbe au pluriel : passif réfléchi.
\item \esp{Se condena la violencia de género.} \\
\emph{La violence de genre est condamnée.} — Sujet patient singulier $\rightarrow$ verbe au singulier.
\item \esp{Se difunden campañas de prevención.} \\
\emph{Des campagnes de prévention sont diffusées.} — Sujet patient pluriel $\rightarrow$ \esp{difunden}.
\item \esp{Se aprobó una nueva ley de protección.} \\
\emph{Une nouvelle loi de protection a été adoptée.} — Sujet patient singulier $\rightarrow$ \esp{aprobó}.
\item \esp{Se atiende a las víctimas en los centros.} (ou mieux : \esp{A las víctimas se les atiende en los centros.}) \\
\emph{On s'occupe des victimes dans les centres.} — COD de personne maintenu avec \esp{a} : impersonnel, verbe au singulier, redondance pronominale \esp{les}.
\item \esp{En la escuela se educa a los niños en igualdad.} \\
\emph{À l'école, on éduque les enfants à l'égalité.}
\item \esp{No se olvida el dolor de las víctimas.} \\
\emph{On n'oublie pas la douleur des victimes.}
\end{enumerate}

\textbf{Points de vigilance :}
\begin{itemize}
\item \textbf{Test de l'accord} : remplacez mentalement par le passif classique (\esp{ser} + participe). Si c'est possible (\esp{Las órdenes son dictadas}), c'est un passif réfléchi $\rightarrow$ accord avec le sujet patient.
\item Erreur fréquente des francophones : \esp{*Se dictan las órdenes} avec l'article n'est pas fautif mais change le sens (défini vs indéfini) ; respectez l'article de la phrase active.
\item Ne jamais écrire \esp{*se detuvieron al agresor} : avec COD de personne introduit par \esp{a}, le verbe reste au singulier.
\end{itemize}
\end{corrige}

\begin{corrige}[Exercice 3 — Conjugaison : complétez avec le temps convenu]
\textbf{Barème indicatif (8 pts)} : 1 point par verbe correctement conjugué (0,5 pt si le mode est juste mais la forme fautive).

\textbf{Solutions :}
\begin{enumerate}
\item \esp{Muchas mujeres \textbf{habían sufrido} en silencio durante años antes de pedir ayuda.} \\
Pluscuamperfecto de indicativo : auxiliaire \esp{haber} à l'imparfait (\esp{habían}) + participe \esp{sufrido}. Antériorité par rapport à \esp{pedir ayuda}.
\item \esp{Es urgente que el gobierno \textbf{tome} medidas más eficaces contra el machismo.} \\
Subjuntivo presente de \esp{tomar} : expression impersonnelle de nécessité \esp{es urgente que} + subjonctif. Verbe en \esp{-ar} $\rightarrow$ terminaison \esp{-e}.
\item \esp{Si la sociedad \textbf{educara / educase} en igualdad, la violencia de género \textbf{disminuiría}.} \\
Phrase hypothétique d'irréel du présent : \esp{si} + imperfecto de subjuntivo $\rightarrow$ condicional simple. Jamais de conditionnel après \esp{si}.
\item \esp{La ley \textbf{entró} en vigor en 2004 y desde entonces \textbf{ha salvado} muchas vidas.} \\
\esp{Entró} : pretérito indefinido (action ponctuelle, datée, achevée). \esp{Ha salvado} : pretérito perfecto (bilan d'une période qui continue jusqu'à aujourd'hui, marqué par \esp{desde entonces}).
\item \esp{No creo que el acoso psicológico \textbf{sea} menos grave que el físico.} \\
Subjuntivo presente de \esp{ser} (irrégulier : \esp{sea}) après \esp{no creo que} (doute, négation de certitude).
\item \esp{Las víctimas \textbf{necesitan} apoyo psicológico, jurídico y económico.} \\
Presente de indicativo : vérité générale, fait réel, aucune marque de subjectivité.
\item \esp{Cuando yo \textbf{era} joven, no \textbf{había} tanta conciencia sobre este problema.} \\
Imperfecto de indicativo : description et situation habituelle dans le passé. \esp{Había} : forme impersonnelle de \esp{haber}, toujours au singulier (\esp{*habían conciencias} est fautif).
\item \esp{Ojalá que pronto \textbf{se logre} una igualdad real entre hombres y mujeres.} \\
Subjuntivo presente après \esp{ojalá (que)} (souhait). Forme pronominale \esp{lograrse} ou passif réfléchi : \esp{se logre}.
\end{enumerate}

\textbf{Points de vigilance :}
\begin{itemize}
\item \esp{Ojalá} est \textbf{toujours} suivi du subjonctif, sans exception : \esp{*ojalá que logra} est une faute éliminatoire au Bac.
\item Contrastez : \esp{Creo que es grave} (indicatif, affirmation) / \esp{No creo que sea grave} (subjonctif, doute).
\item \esp{Haber} impersonnel : singulier à tous les temps (\esp{hay, había, hubo, habrá, haya}).
\end{itemize}
\end{corrige}

\begin{corrige}[Exercice 4 — Texte à trous : conjugaison et connecteurs]
\textbf{Barème indicatif (6 pts)} : 0,5 point par case correctement remplie.

\textbf{Texte complété :}
\begin{textebox}[La lucha contra la violencia de género (versión corregida)]
La violencia de género \textbf{es} una violación de los derechos humanos que \textbf{afecta} a millones de mujeres. \textbf{Por ejemplo}, no solo \textbf{se manifiesta} a través de golpes, \textbf{sino} también mediante control económico y aislamiento. \textbf{Sin embargo}, muchas víctimas no \textbf{denuncian} por miedo. \textbf{Por tanto}, en España la ley \textbf{garantiza} la protección, pero \textbf{hace} falta más recursos. \textbf{En definitiva}, la educación en igualdad \textbf{es} la clave para erradicar el machismo.
\end{textebox}

\textbf{Justifications :}
\begin{itemize}
\item \esp{es / afecta} : présent de vérité générale ; \esp{afecta} s'accorde avec \esp{que} (= \esp{una violación}, singulier).
\item \esp{se manifiesta} : verbe pronominal \esp{manifestarse} ; tournure réfléchie équivalente à « elle se manifeste ».
\item \esp{no solo... sino también} : corrélation obligatoire ; ne pas écrire \esp{*no solo... pero también}.
\item \esp{denuncian} : présent d'habitude, accord avec \esp{víctimas}.
\item \esp{hace falta} : expression impersonnelle figée au singulier (\esp{*hacen falta recursos} est un hispanisme répandu mais à éviter en production écrite normée ; on préfère \esp{hacen falta} toléré ou \esp{se necesitan más recursos}).
\item Connecteurs : \esp{Por ejemplo} (illustration), \esp{Sin embargo} (opposition : la loi existe mais les victimes ont peur), \esp{Por tanto} (conséquence), \esp{En definitiva} (conclusion).
\end{itemize}
\end{corrige}

\begin{corrige}[Exercice 5 — Appariement : expressions idiomatiques]
\textbf{Barème indicatif (5 pts)} : 0,5 pt par appariement correct + 0,5 pt par phrase exemple pertinente et correcte.

\textbf{Appariements :} 1-b ; 2-c ; 3-d ; 4-a ; 5-e.

\textbf{Phrases exemples (modèles) :}
\begin{enumerate}
\item \esp{Después de años de silencio, la víctima decidió \textbf{romper el silencio} y denunciar a su agresor.} \\
\emph{Après des années de silence, la victime a décidé de briser le silence et de dénoncer son agresseur.}
\item \esp{Las asociaciones feministas \textbf{pusieron el grito en el cielo} cuando el agresor quedó en libertad.} \\
\emph{Les associations féministes ont poussé des cris d'indignation quand l'agresseur a été libéré.}
\item \esp{El maltrato psicológico \textbf{pasa factura}: muchas víctimas sufren depresión años después.} \\
\emph{La maltraitance psychologique laisse des traces : beaucoup de victimes souffrent de dépression des années après.}
\item \esp{Para \textbf{cortar de raíz} el machismo, hay que educar a los niños en la igualdad desde la infancia.} \\
\emph{Pour éradiquer le machisme à la racine, il faut éduquer les enfants à l'égalité dès l'enfance.}
\item \esp{Es hora de \textbf{alzar la voz} contra todas las formas de discriminación.} \\
\emph{Il est temps d'élever la voix contre toutes les formes de discrimination.}
\end{enumerate}

\textbf{Points de vigilance :}
\begin{itemize}
\item \esp{Poner el grito en el cielo} exprime une \textbf{indignation collective}, pas une plainte individuelle.
\item \esp{Pasar factura} : ne pas traduire littéralement « passer la facture » ; le sens est « avoir des conséquences néfastes à long terme ».
\item Ces expressions relèvent d'un registre soutenu/journalistique : parfaites pour enrichir un commentaire de texte.
\end{itemize}
\end{corrige}

\begin{corrige}[Exercice 6 — Bac blanc : compréhension et analyse de texte]
\textbf{Barème indicatif (20 pts)} : Compréhension 5 pts (1 pt/item) ; Vrai/Faux 4 pts (0,5 pt réponse + 0,5 pt citation) ; Analyse linguistique 4 pts (2 pts/item) ; Relevé grammatical 7 pts (2 pts connecteurs + 3 pts subjonctif + 2 pts dérivation).

\textbf{1. Comprensión}
\begin{enumerate}
\item \esp{María González es ingeniera informática y lleva cinco años trabajando en una empresa tecnológica de Madrid.}
\item \esp{En las reuniones, sus propuestas se ignoran hasta que un hombre las repite, y su jefe le pide que traiga café a los clientes, algo que nunca pide a sus compañeros varones.}
\item \esp{El 67\% de las profesionales españolas afirma haber sufrido discriminación salarial o de ascenso por razón de género.}
\item \esp{La brecha salarial en el sector tecnológico alcanza el 18,5\%.}
\item \esp{Los planes de igualdad obligatorios, implementados desde 2021, han mejorado la situación en las empresas que los aplican.}
\end{enumerate}

\textbf{2. Verdadero / Falso}
\begin{enumerate}
\item \textbf{Falso.} Citation : \esp{«A pesar de tener mejor expediente académico y más experiencia que varios de sus compañeros varones...»}
\item \textbf{Falso.} Citation : \esp{«La brecha salarial en el sector tecnológico alcanza el 18,5\%.»} (18,5 \% est supérieur à 10 \%).
\item \textbf{Falso.} Le texte ne dit rien de tel ; au contraire, la logique du texte (et la phrase de conclusion sur la corresponsabilité) montre que la charge familiale pèse sur les femmes. Toute affirmation non étayée par le texte est à juger fausse ou non vérifiable.
\item \textbf{Verdadero.} Citation : \esp{«Las empresas que implementan planes de igualdad obligatorios desde 2021 notan mejoras.»}
\end{enumerate}

\textbf{3. Analyse linguistique}
\begin{enumerate}
\item \esp{«El techo de cristal no se ve, pero se siente»} : métaphore qui désigne les barrières invisibles (préjugés, stéréotypes, réseaux informels masculins) qui empêchent les femmes d'accéder aux postes de responsabilité, sans qu'aucune règle explicite ne les en empêche. L'obstacle n'est pas visible (\esp{no se ve}) mais ses effets sont bien réels (\esp{se siente}) : stagnation de carrière, sous-représentation aux postes de direction.
\item \esp{«La maternidad sigue siendo un freno para la carrera profesional»} : la maternité est perçue par certains employeurs comme une indisponibilité ou un coût (congé maternité, absences), ce qui ralentit les promotions et les augmentations des femmes, alors que la paternité ne pénalise pas — voire valorise — la carrière des hommes.
\end{enumerate}

\textbf{4. Relevé grammatical}
\begin{enumerate}
\item Connecteurs d'opposition : \esp{pero} (\esp{«El techo de cristal no se ve, pero se siente»}) ; \esp{a pesar de} (\esp{«A pesar de tener mejor expediente...»}). On accepte aussi \esp{hasta que} non (limite temporelle) ; retenir \esp{pero} et \esp{a pesar de}.
\item Subjonctif présent : \esp{«...pero es fundamental que la sociedad entera \textbf{asuma} la corresponsabilidad.»} — \esp{Asuma} : subjonctif présent d'\esp{asumir}, justifié par la structure impersonnelle de nécessité \esp{es fundamental que} + subjonctif.
\item Dérivation préfixale : \esp{desigual} n'apparaît pas, mais on relève \esp{invisible} (préfixe \esp{in-} + \esp{visible}) et \esp{corresponsabilidad} (préfixe \esp{co-/cor-} + \esp{responsabilidad}). On accepte aussi \esp{discriminación} analysé avec le préfixe \esp{dis-}.
\end{enumerate}

\textbf{Points de vigilance :}
\begin{itemize}
\item Pour le Vrai/Faux, la \textbf{citation exacte} est obligatoire : une réponse sans justification vaut zéro au Bac.
\item Recopiez les citations entre guillemets, sans modifier l'orthographe du texte.
\item Pour le relevé du subjonctif, exigez la phrase complète + le nom du verbe à l'infinitif + la règle d'emploi.
\end{itemize}
\end{corrige}

\begin{corrige}[Exercice 7 — Thème et Version]
\textbf{Barème indicatif (10 pts par version)} : fidélité au sens 4 pts ; correction grammaticale (modes, temps, accords) 3 pts ; lexique précis 2 pts ; orthographe et ponctuation 1 pt.

\textbf{A. Thème — Traduction modèle :}
\begin{textebox}[Traducción propuesta]
En Camerún, como en España, la violencia contra las mujeres sigue siendo una lacra silenciosa. Demasiadas víctimas no se atreven a denunciar por miedo a las represalias o por vergüenza. Sin embargo, la ley protege a las mujeres y prevé penas severas para los agresores. Es urgente que la sociedad entera —hombres y mujeres— se movilice para romper el ciclo de la violencia. La educación en la igualdad desde la infancia es la mejor vacuna contra el machismo.
\end{textebox}

\textbf{Points de vigilance (thème) :}
\begin{itemize}
\item « Trop de victimes » : \esp{demasiadas víctimas} (accord féminin pluriel), pas \esp{*demasiados}.
\item « n'osent pas porter plainte » : \esp{no se atreven a denunciar} ou \esp{no se atreven a presentar una denuncia} ; verbe pronominal \esp{atreverse a} + infinitif.
\item « Il est urgent que... se mobilise » : \esp{es urgente que... se movilice} — \textbf{subjonctif obligatoire} ; \esp{*se moviliza} est une faute grave.
\item « dès l'enfance » : \esp{desde la infancia}, pas \esp{*desde la niñez} (possible mais moins idiomatique ici).
\item « le machisme » : \esp{el machismo} (masculin).
\end{itemize}

\textbf{B. Version — Traduction modèle :}

La violence de genre ne connaît ni frontières, ni classes sociales, ni niveau d'éducation. Elle touche des femmes de tous âges et de toutes conditions, bien que les plus vulnérables — les femmes migrantes, en situation de handicap, vivant dans la pauvreté — subissent une double discrimination. Les institutions doivent garantir non seulement la poursuite du délit, mais aussi la réparation intégrale de la victime : assistance juridique gratuite, soutien psychologique spécialisé, formation professionnelle et logement digne. C'est seulement ainsi que l'on brise le cercle de la violence et que l'on rend à la femme son autonomie et sa dignité.

\textbf{Points de vigilance (version) :}
\begin{itemize}
\item \esp{No entiende de fronteras} : ne pas traduire « ne comprend pas les frontières » ; sens : « ne connaît pas de frontières ».
\item \esp{Persecución del delito} : « la poursuite (judiciaire) du délit », pas « la persécution » (faux ami partiel).
\item \esp{Reparación integral} : « réparation intégrale/globale » (terme juridique).
\item \esp{Vivienda digna} : « logement décent/digne », pas « habitation digne ».
\end{itemize}
\end{corrige}

\begin{corrige}[Exercice 8 — Expression écrite : commentaire personnel guidé]
\textbf{Barème indicatif (20 pts)} : respect des contraintes (longueur, structure) 4 pts ; richesse du lexique thématique (5 mots minimum) 4 pts ; connecteurs (2 minimum) 2 pts ; subjonctif correctement employé 3 pts ; correction grammaticale et orthographe 4 pts ; pertinence des arguments 3 pts.

\textbf{Proposition de rédaction modèle :}
\begin{textebox}[Comentario modelo]
¿Es posible la igualdad real entre hombres y mujeres en mi país? En mi opinión, sí es posible, pero exige un cambio profundo de mentalidades. En primer lugar, el machismo sigue siendo un obstáculo estructural: en Camerún, como en España, muchas mujeres sufren violencia de género y acoso laboral sin atreverse a presentar una denuncia por miedo o por dependencia económica. Por tanto, es fundamental que el Estado refuerce la protección integral de las víctimas y que los agresores sean castigados. Sin embargo, los avances existen: cada vez más mujeres acceden a la educación superior y rompen el techo de cristal en la política y en las empresas. Además, la corresponsabilidad en el hogar empieza a liberar a las mujeres de las tareas domésticas. En definitiva, la educación en igualdad desde la infancia es la clave. Ojalá que las nuevas generaciones construyan una sociedad donde ser mujer no sea un riesgo, sino un derecho pleno.
\end{textebox}

\emph{Traduction :} L'égalité réelle entre hommes et femmes est-elle possible dans mon pays ? À mon avis, oui, mais elle exige un changement profond des mentalités. D'abord, le machisme reste un obstacle structurel : au Cameroun comme en Espagne, beaucoup de femmes subissent des violences de genre et du harcèlement au travail sans oser porter plainte, par peur ou par dépendance économique. Il est donc fondamental que l'État renforce la protection intégrale des victimes et que les agresseurs soient punis. Cependant, des progrès existent : de plus en plus de femmes accèdent à l'enseignement supérieur et brisent le plafond de verre en politique et dans les entreprises. De plus, la coresponsabilité au foyer commence à libérer les femmes des tâches domestiques. En définitive, l'éducation à l'égalité dès l'enfance est la clé. Puissent les nouvelles générations construire une société où être une femme ne soit pas un risque, mais un droit plein.

\textbf{Vérification des contraintes dans le modèle :}
\begin{itemize}
\item Lexique de la leçon (7 occurrences) : \esp{machismo, violencia de género, acoso, denuncia, víctimas, techo de cristal, corresponsabilidad}.
\item Connecteurs (4) : \esp{En primer lugar, Por tanto, Sin embargo, Además, En definitiva}.
\item Subjonctifs (2) : \esp{es fundamental que el Estado refuerce... sean castigados} ; \esp{Ojalá que... construyan}.
\item Longueur : environ 140 mots, conforme.
\end{itemize}

\textbf{Points de vigilance :}
\begin{itemize}
\item Sanctionner tout subjonctif manquant après \esp{es necesario que}, \esp{ojalá}, \esp{no creo que}.
\item Vérifier les accents : \esp{opinión, educación, política, económica} ; et les signes \esp{¿} \esp{¡} en ouverture de phrase.
\item Une opinion personnelle claire (\esp{en mi opinión, creo que}) est exigée : le commentaire ne doit pas être un simple résumé du texte support.
\end{itemize}
\end{corrige}

\newpage

\coursec{Fiche bilan}

\begin{popfigure}
\begin{center}
\begin{tikzpicture}[mindmap, grow cyclic, every node/.style={concept, align=center, text width=2.6cm, font=\small},
root concept/.append style={concept color=popblue, text=white, font=\bfseries\small, text width=3cm},
level 1 concept/.append style={level distance=4.2cm, sibling angle=60, font=\scriptsize, text width=2.4cm}]
\node[root concept]{Violencias y discriminaciones contra las mujeres}
child[concept color=poporange]{node[align=center]{Texto soporte\\ \esp{La otra cara del silencio}\\ lectura y descubrimiento}}
child[concept color=popgreen]{node[align=center]{Comprensión\\ preguntas, ideas directrices, vocabulario en contexto}}
child[concept color=poppurple]{node[align=center]{Léxico\\ \esp{violencia de género, machismo, maltrato, denuncia, igualdad, acoso}}}
child[concept color=poppink]{node[align=center]{Gramática\\ pasiva con \esp{se}, ser/estar, subjuntivo tras verbos de voluntad}}
child[concept color=popgold]{node[align=center]{Método del comentario\\ presentación, ideas, opinión personal}}
child[concept color=popblue!70]{node[align=center]{Producción\\ comentario guiado, debate, expresión escrita}};
\end{tikzpicture}
\end{center}
\legende{Carte mentale de la leçon : les six compétences travaillées autour du texte sur les violences faites aux femmes.}
\end{popfigure}

\begin{bilan}
\textbf{1. Lexique clé à connaître par cœur}
\begin{center}
\begin{tabularx}{\linewidth}{|X|X|}\hline
\rowcolor{popblueL} \textbf{Español} & \textbf{Français} \\ \hline
\esp{la violencia de género} & les violences faites aux femmes \\ \hline
\esp{el machismo} & le machisme \\ \hline
\esp{el maltrato (físico / psicológico)} & la maltraitance (physique / psychologique) \\ \hline
\esp{la denuncia / denunciar} & la dénonciation, la plainte / dénoncer \\ \hline
\esp{la igualdad (de derechos)} & l'égalité (des droits) \\ \hline
\esp{el acoso (laboral, escolar)} & le harcèlement (au travail, à l'école) \\ \hline
\esp{la víctima / el agresor} & la victime / l'agresseur \\ \hline
\esp{la brecha salarial} & l'écart de salaire \\ \hline
\esp{la concienciación} & la sensibilisation \\ \hline
\esp{el matrimonio forzado} & le mariage forcé \\ \hline
\esp{luchar por sus derechos} & lutter pour ses droits \\ \hline
\esp{romper el silencio} & briser le silence \\ \hline
\end{tabularx}
\end{center}

\textbf{2. Grammaire à revoir pour ce thème}
\begin{center}
\begin{tabularx}{\linewidth}{|X|X|}\hline
\rowcolor{popblueL} \textbf{Règle} & \textbf{Exemple} \\ \hline
Passif avec \esp{se} + verbe à la 3\ieme{} personne & \esp{Se denuncian muchos casos.} « On dénonce beaucoup de cas. » \\ \hline
\esp{ser} (identité, caractère) / \esp{estar} (état, situation) & \esp{La ley es justa.} / \esp{Ella está protegida.} \\ \hline
Subjonctif après les verbes de volonté et de sentiment & \esp{Espero que las leyes cambien.} \\ \hline
\esp{por} (cause, moyen) / \esp{para} (but, destination) & \esp{luchar por la igualdad} / \esp{leyes para proteger} \\ \hline
Futur et conditionnel pour projeter et nuancer & \esp{La situación mejorará.} / \esp{Habría que actuar.} \\ \hline
\end{tabularx}
\end{center}

\textbf{3. Méthode express du commentaire de texte}
\begin{enumerate}
\item \textbf{Présentation} : nature du texte, source, auteur, thème, problématique (\esp{El texto trata de...}).
\item \textbf{Idées directrices} : deux ou trois axes, chacun illustré par des citations courtes traduites.
\item \textbf{Commentaire personnel} : ton avis argumenté, comparaison avec le contexte camerounais ou africain, ouverture.
\item \textbf{Langue} : connecteurs (\esp{en primer lugar, sin embargo, por lo tanto, en conclusión}), reprise du lexique du thème.
\end{enumerate}

\textbf{4. Phrases utiles pour le commentaire}
\begin{itemize}
\item \esp{El autor denuncia que...} « L'auteur dénonce le fait que... »
\item \esp{Según el texto, las mujeres sufren...} « Selon le texte, les femmes subissent... »
\item \esp{En mi opinión, es fundamental educar a los niños en la igualdad.} « À mon avis, il est essentiel d'éduquer les enfants dans l'égalité. »
\item \esp{A mi juicio, todavía queda mucho por hacer.} « Selon moi, il reste encore beaucoup à faire. »
\end{itemize}
\end{bilan}

\begin{aretenir}[Checklist avant le Bac]
\begin{itemize}
\item[$\square$] Je sais définir et employer \esp{violencia de género}, \esp{machismo}, \esp{maltrato}, \esp{denuncia}, \esp{igualdad}, \esp{acoso}.
\item[$\square$] Je distingue sans erreur \esp{ser} et \esp{estar} dans mes phrases de commentaire.
\item[$\square$] Je sais former le passif avec \esp{se} : \esp{se denuncia / se denuncian}.
\item[$\square$] Je mets correctement le subjonctif après \esp{esperar que}, \esp{es importante que}, \esp{quiero que}.
\item[$\square$] Je choisis correctement entre \esp{por} (cause) et \esp{para} (but).
\item[$\square$] Je sais présenter un texte en trois phrases : nature, source, thème et problématique.
\item[$\square$] Je dégage deux ou trois idées directrices et je les justifie par des citations courtes.
\item[$\square$] Je sais exprimer mon opinion : \esp{en mi opinión}, \esp{a mi juicio}, \esp{me parece que}.
\item[$\square$] J'utilise des connecteurs logiques variés dans ma rédaction.
\item[$\square$] Je n'oublie ni les accents (\esp{género, víctima, igualdad}) ni les signes ¿ et ¡.
\item[$\square$] Je fais attention aux faux amis : \esp{denunciar} = dénoncer (pas « dénoncer une erreur » au sens de « signaler une faute ») ; \esp{la víctima} est toujours féminin.
\item[$\square$] Je peux comparer la situation en Espagne, en Amérique latine, en Guinée équatoriale et au Cameroun avec des exemples simples.
\end{itemize}
\end{aretenir}

\findecours

\end{document}
